% concatenated from 2026-03-25_-_ArXiV_version.tex, Packages.tex, Commands.tex
\documentclass[a4paper, 11pt, reqno]{amsart}

%Packages

\usepackage[utf8]{inputenc}
\usepackage[T1]{fontenc}
\DeclareUnicodeCharacter{03C8}{\ensuremath{\psi}}
\usepackage[english]{babel}
\usepackage{latexsym,amsmath,amssymb,amsfonts}
\usepackage{dsfont}
\usepackage{color}
\usepackage{graphicx}
\usepackage{float}
\usepackage{esint}
\usepackage{bm}
\usepackage{mathtools}
\usepackage[top=1in, bottom=1.25in, left=1in, right=1in]{geometry}
\usepackage{enumitem}
\usepackage{hyperref}
\usepackage{cleveref}
\usepackage{bbm}
\usepackage{tikz}                        
\usepackage{pgfplots}    
\usepackage{soul}         
%\usepackage{svg}
\usepackage{mathrsfs}
%\usepackage{adjustbox}
%\usepackage{transparent-nometadata}

\hypersetup{
  colorlinks = true,
  linkcolor = blue,
  citecolor = red
}

%%%%%%%%%%%%%%%%%%%%%%%%%%%%%%%%%%%%%%%%%%%%

\theoremstyle{plain}
\begingroup
\newtheorem{theorem}{Theorem}[section]
\newtheorem{lemma}[theorem]{Lemma}
\newtheorem{proposition}[theorem]{Proposition}
\newtheorem{corollary}[theorem]{Corollary}
\endgroup
\theoremstyle{definition}
\begingroup
\newtheorem{definition}[theorem]{Definition}
\newtheorem{remark}[theorem]{Remark}
\newtheorem{example}[theorem]{Example}
\newtheorem{exercise}[theorem]{Exercise}
\endgroup
\theoremstyle{remark}

%\setlength\parindent{0pt}
%Spaces

\newcommand{\N}{\mathbb N}
\newcommand{\Z}{\mathbb Z}
\newcommand{\Q}{\mathbb Q}
\newcommand{\R}{\mathbb R}
\newcommand{\C}{\mathbb C}
\newcommand{\M}{\mathbb M}
\newcommand{\T}{\mathcal T}
\renewcommand{\S}{\mathbb S}
\newcommand{\po}{\mathcal{P}(\Omega)}
\newcommand{\eps}{\varepsilon}
\DeclareMathOperator{\spt}{spt}

\newcommand{\abs}[1]{\left|#1\right|}
\newcommand{\norm}[1]{\left\|#1\right\|}
\newcommand{\jump}[1]{[\![#1]\!]}

%Math symbols

\newcommand{\tr}{{\rm trace\,}}
\renewcommand{\div}{{\rm div}}
\newcommand{\dist}{{\rm dist\,}}
\newcommand{\supp}{{\rm supp\,}}
\newcommand{\id}{\mathrm{Id}}
\newcommand{\diam}{{\rm diam\,}}
\newcommand{\ca}{\mathbb{1}}
\newcommand{\pj}{\mathcal{PJ}}



%Measures

\renewcommand{\l}{{\mathcal L}}
\renewcommand{\ln}{{\mathcal{L}^N}}
\newcommand{\hn}{{\mathcal{H}}^N}
\newcommand{\hno}{{\mathcal H}^{N-1}}
\newcommand{\hnt}{{\mathcal H}^{N-2}}
\newcommand{\ho}{{\mathcal H}^{1}}
\newcommand{\hz}{{\mathcal H}^{0}}
\renewcommand{\H}{\mathcal H}


%Integration

\newcommand{\dd}{\,\mathrm{d}}
\renewcommand{\l}{\,\dd{\mathcal L}}
\newcommand{\dhn}{\,\dd{\mathcal{H}}^N}
\newcommand{\dhno}{\,\dd{\mathcal H}^{N-1}}
\newcommand{\dhnt}{\,\dd{\mathcal H}^{N-2}}
\newcommand{\dho}{\,\dd{\mathcal H}^{1}}
\newcommand{\dhz}{\,\dd{\mathcal H}^{0}}




%Abbreviation

\newcommand{\mca}{\mathcal{A}}
\newcommand{\wto}{\rightharpoonup}
\newcommand{\wtom}{\stackrel{w*}{\rightharpoonup}}
\newcommand{\meno}{\!\setminus\!}
\newcommand{\e}{\varepsilon}
\newcommand{\p}{\varphi}
\renewcommand{\o}{\Omega}
\newcommand{\rga}{\rightarrow}
\renewcommand{\sp}{H^1(\o;\R^M)}
\newcommand{\dw}{\mathrm{d}_W}
\def\ca{\mathbbmss{1}}
\renewcommand{\arg}{\mathrm{arg}}


%Functionals

\newcommand{\f}{\mathcal{F}}
\newcommand{\g}{\mathcal{G}}
\newcommand{\per}{\mathcal{P}}
\newcommand{\nl}{\mathcal{NL}}
\newcommand{\ms}{\mathcal{MS}}


%Other stuff

\newcommand{\mloc}{m_{\rm loc}}
\newcommand{\mglob}{m_{\rm glob}}


\newcommand{\restr}{%
  \,\raisebox{-.127ex}{\reflectbox{\rotatebox[origin=br]{-90}{$\lnot$}}}\,%
}


%Per il comando di integrale mediato
\newcommand{\average}{{\mathchoice {\kern1ex\vcenter{\hrule height.4pt
width 6pt
depth0pt} \kern-9.7pt} {\kern1ex\vcenter{\hrule height.4pt width 4.3pt
depth0pt}
\kern-7pt} {} {} }}
\newcommand{\med}{\average\int}

\newcommand{\red}[1]{{\color{red} {#1}}}
\newcommand{\blue}[1]{{\color{blue} {#1}}}


\title[Low Regularity Potentials]{Low regularity potentials in heterogeneous Cahn--Hilliard functionals}

\author{Riccardo Cristoferi, Jakob Deutsch, Luca Pignatelli}

\DeclareGraphicsExtensions{.pdf}

\begin{document}


\begin{abstract}
In this paper, we study the prototypical model of liquid-liquid phase separation, the Cahn-Hilliard functional, in a highly irregular setting.
    Specifically, we analyze potentials with low regularity vanishing on space-dependent wells.
    Under remarkably weak hypotheses, we establish a robust compactness result.
    Strengthening the regularity of the wells and of the growth of the potential close to the wells only slightly, we completely characterize the asymptotic behavior of the associated family of functionals through a $\Gamma$-convergence analysis.
    As a notable technical result, we prove the existence of geodesics for a degenerate metric and establish a uniform bound on their Euclidean length.
\end{abstract}

\maketitle


\tableofcontents


\section{Introduction}

The Cahn--Hilliard functional is the prototypical model for liquid-liquid phase separation.
It was initially proposed by van der Waals in \cite{vanDW}, and later independently rediscovered by Cahn and Hilliard in \cite{CH}.

We start by describing the physical situation that we are dealing with.
We consider a physical system inside a container $\o \subset \R^N$.
We assume that the system is described by a phase (or order) parameter $u \colon \o \to \R^M$.
For instance, in the original derivation of the model, the system under investigation was a mixture of two fluids. In such a case, the phase parameter represents the density.
We are interested in describing stable equilibrium configurations of the system.
Using a variational point of view, these correspond to local minimizers of the Gibbs free energy.
Under \emph{homogeneous conditions}, and after a rescaling, this is given by a functional of the form
\[
F_\varepsilon(u) \coloneqq \int_\Omega \left[ \frac{1}{\varepsilon}W(u(x)) + \varepsilon|\nabla u(x)|^2 \right] \dd x.
\]
Here, $W\colon \R^M\to[0,\infty)$ is the free energy density.
Assuming for simplicity that the physical system has two stable phases, modeled by $a,b,\in\R^M$, we impose that the potential $W$ vanishes only at $a$ and $b$.
The second term of the functional penalizes sharp oscillations of the phase variable.
Finally, the parameter $\varepsilon>0$ is related to the scale of the transition region between bulk regions where the phase variable is close to the stable phases $a$ and $b$, and it is usually very small.

The mathematical study of the functional $F_\varepsilon$ was, and still is, a source of many deep and interesting problems.
In particular, here we focus on the study of its asymptotic behavior as $\varepsilon\to0$.
In the community of calculus of variations, such studies make use of $\Gamma$-convergence, the notion of variational convergence introduced in \cite{DGFra}, that allows to capture the asymptotic behavior of minimizers of the family of functionals $\{F_\varepsilon\}_{\varepsilon>0}$, as well as the asymptotic of the minimal energy.
A similar analysis can be performed also in the case where a \emph{mass constraint} is imposed. Namely, given $m\in\R^M$, we consider the problem
\begin{equation}\label{eq:min_pb_intro}
\min\left\{ F_\varepsilon(u) : u\in H^1(\Omega;\R^M), \, \int_\Omega u(x)\dd x = m  \right\}.
\end{equation}
Here, we recall some contributions, without claiming to be exhaustive.
We first focus on the case where the phase variable is scalar-valued, namely when $M=1$.
Gurtin (see \cite{Gur}) conjectured that in the limit as $\varepsilon\to0$, any family of minimizers of $\{F_\varepsilon\}_{\varepsilon>0}$ converge to a piecewise constant function that partitions $\Omega$ into two regions separated by an interface with minimal surface area.
This conjecture was proved rigorously 
by Carr, Gurtin, and Slemrod for $N = 1$ (see \cite{CarGurSle}), and by Modica (see \cite{modica87}) and by Sternberg (see \cite{Ste_Sing}) for $N \ge 2$ (see also \cite{ModMor_Limite, ModMor_Esempio}).

The case of vector-valued wells, namely $M\geq 1$, was considered in \cite{FonTar} by Fonseca and Tartar.
What the author proved is a compactness result with respect to the strong $L^1(\Omega;\R^M)$ topology for sequences of uniformly bounded energy, and that the $\Gamma$-limit of the sequence $\{F_\varepsilon\}_\varepsilon$ is given by
\[
F_0(u)\coloneqq \sigma\,\mathrm{Per}(\{u=a\};\Omega),
\]
for functions $u\in BV(\Omega;\R^M)$ taking values only in $\{a,b\}$, and $+\infty$ else in $L^1(\Omega;\R^M)$.
Here, $\mathrm{Per}(\{u=a\};\Omega)$ denotes the perimeter of the set $\{u=a\}$ in $\Omega$ (see Section \ref{sec:BV} for more details on sets of finite perimeter), and the constant $\sigma>0$ is given by
\begin{equation}\label{eq:sigma_intro}
\sigma\coloneqq \inf\left\{\int_{-1}^1 2\sqrt{W(\gamma(t))}|\gamma'(t)|\dd t \,:\, \gamma\in W^{1,1}([-1,1];\R^M), \gamma(-1)=a, \gamma(1)=b \right\}.
\end{equation}
Such a quantity represents the minimal energy needed to transition from one well to the other.
An optimal profile, namely a solution to the above minimization problem, approximates the behavior of a solution to \eqref{eq:min_pb_intro}.
Note that the minimization problem is a degenerate geodesic problem, where the degeneracy comes from the fact that the weight $\sqrt{W}$ vanishes at the endpoints of the curves.
In scalar case $M=1$, the change of variables $s=\gamma(t)$ gives that
\begin{equation}\label{eq:sigma_scalar_intro}
\sigma = \int_a^b 2\sqrt{W(s)}\dd s,
\end{equation}
for all admissible curves.

Since it is of particular interest for the goal of this paper, we remark that in \cite{FonTar} the authors assumed that the potential $W$ behaves quadratically near the wells $a$ and $b$.

There are several extensions and variants of the above mentioned results.
For a more comprehensive overview of such literature, the reader can consult the introduction of the paper \cite{cristoferi2025phaseseparationmultiplyperiodic}.\\

The study of non-homogeneous Cahn--Hilliard functionals initiated with the paper \cite{Bou} by Bouchitt\`{e}, where the author studied the fully coupled case with two space-dependent scalar wells.
Namely, for $u\in H^1(\Omega)$, Bouchitt\`{e} considered the functional
\[
G_\varepsilon(u) \coloneqq \frac{1}{\varepsilon}\int_\Omega g(x, u(x), \varepsilon\nabla u(x))\dd x,
\]
where $g:\Omega\times\R\times\R^N\to[0,\infty)$ is a continuous function with $g(x,u,0)=0$ if and only if $u\in\{a(x), b(x)\}$ for all $x\in\Omega$, and $p\mapsto g(x,u,p)$ is a convex function achieving its strict minimum at $p=0$, where it is also assumed to be differentiable.
The space-dependent wells $a,b:\Omega\to\R$ are assumed to be Lipschitz continuous.
The author identified the limiting functional in the sense of $\Gamma$-convergence. Since we do not want to enter into the technicalities needed in order to treat the general case, we consider here the special case
\[
g(x,u,p) = W(x,u) + |p|^2.
\]
In such a special case, the limiting functional writes as
\[
G_0(u)\coloneqq \int_{J_u\cap\Omega} \sigma(x; a(x), b(x))\dhno(x),
\]
for functions $u\in BV(\Omega)$ such that $u(x)\in\{a(x),b(x)\}$ for a.e. $x\in\Omega$, and $+\infty$ else in $L^1(\Omega)$.
Here, $J_u$ denotes the jump set of the function $u$ (see Section \ref{sec:BV}) for all $x\in\Omega$, and $p,q\in\R^M$, we set
\begin{equation}\label{eq:sigma_moving_scalar}
\sigma(x;p,q)\coloneqq \int_p^q 2\sqrt{W(x,s)}\dd s.
\end{equation}
Roughly speaking, the limiting energy density is space-dependent, and is given, at any fixed point $x\in \Omega$ by the same degenerate geodesic problem \eqref{eq:sigma_scalar_intro} where the space variable is frozen.

Then, in \cite{CriGra} the first author together with Gravina considered the case of multiple vector-valued wells for the uncoupled functional.
Namely, the authors considered the functional
\[
H_\varepsilon(u) \coloneqq \int_\Omega \left[ \frac{1}{\varepsilon}W(x,u(x)) + \varepsilon|\nabla u(x)|^2 \right] \dd x,
\]
for $u\in H^1(\Omega;\R^M)$, where $W(x,u)=0$ if and only if $u\in\{z_1(x),\dots,z_k(x)\}$ for all $x\in\Omega$.
Being the first result in that direction for vector-valued wells, strong assumptions were needed.
Indeed, the wells $z_1,\dots,z_k\colon\Omega\to\R^M$ were required to be Lipschitz continuous, and the potential $W\colon\Omega\times\R^M\to[0,\infty)$ to be Lipschitz continuous in the first variable, and of class $C^2$ in the second. Moreover, it was assumed that the potential was \emph{exactly} quadratic near the wells.
In such a case, the authors proved a compactness result, and a $\Gamma$-convergence result with respect to the strong $L^1(\Omega;\R^M)$ topology. In particular, it holds that the limiting functional is given by
\begin{equation}\label{eq:functional_H0}
H_0(u)\coloneqq \int_{J_u\cap\Omega} \sigma(x; a(x), b(x))\dhno(x),
\end{equation}
where
\begin{equation}\label{eq:sigma_intro_moving}
\sigma(x;p,q)\coloneqq \inf\left\{\int_{-1}^1 2\sqrt{W(x,\gamma(t))}|\gamma'(t)|\dd t : \gamma\in W^{1,1}([-1,1];\R^M), \gamma(-1)=p, \gamma(1)=q \right\}.
\end{equation}
Note that this is the same problem as in \eqref{eq:sigma_intro} where the space variable is frozen, similarly to what happens with \eqref{eq:sigma_moving_scalar} and \eqref{eq:sigma_scalar_intro}. 
The latter assumption has been weakened in the paper \cite{CriFonGan_moving_sub} by the first author, Ganedi and Fonseca, by requiring the potential to be \emph{controlled} from above and from below by two quadratic functions near the wells.
Moreover, the authors were able to also consider the case where the  wells and the potential are discontinuous on the boundary of a polyhedral partition of the domain $\Omega$.\\

We would like to remark that, in the case of scalar-valued wells, like in \cite{Bou}, there is no need to assume any particular behavior of the potential near the wells. The technical reason is that the optimal profile at a point $x$ will be a suitable reparametrization of the interval $[a(x), b(x)]$. In particular, using the continuity of the wells, the length of all of these intervals is uniformly bounded from above.
On the other hand, in the vector-valued case, the optimal profile is a \emph{curve} solving the minimization problem in \eqref{eq:sigma_intro}.
In such a case, the behavior of the potential $W$ around the wells is essential to ensure that such a curve has an \emph{Euclidean length} that is uniformly bounded from above. This requirement is needed in order to prove the $\Gamma$-convergence result.
This is the main difference from the technical point of view between the scalar and the vector-valued case.\\

The goal of the paper is to prove that it is possible to significantly weaken the requirements on the regularity of the potential and on its behavior close to the wells and still be able to obtain compactness and the $\Gamma$-limit result.

In particular, we consider the functional
\[
\mathcal{F}_\varepsilon(u) \coloneqq \int_\Omega \left[ \frac{1}{\varepsilon}W(x,u(x)) + \varepsilon|\nabla u(x)|^2 \right] \dd x,
\]
where, for each $x\in\Omega$, the function $u\mapsto W(x,u)$ vanishes at $\{a(x), b(x)\}$ for functions $a,b:\Omega\to\R^M$.
We prove that compactness holds under extremely weak assumptions (see \ref{H1}-\ref{H5}): we consider a Carath\'{e}odory potential with mild bounds close to the wells and at least linear growth at infinity, and bounded wells of bounded variation (see Theorem \ref{thm:compactness}).
We are able to obtain such a result by exploiting the theory of Young measures, and by carefully using the chain rule for functions of bounded variation.

The $\Gamma$-convergence result (see Theorem \ref{thm:Gamma_conv}) requires strengthening these assumptions.
Indeed, we consider wells that are allowed to jump on a piecewise $C^2$ partition of the domain $\Omega$, and that are of class $H^1\cap C^0$ inside each set of such a partition (see \ref{H6}).
The potential is assumed to also satisfies a relative continuity condition (see \ref{H7}) and a more quantitative control close to the wells (see \ref{H8}).
This set of assumptions is extremely weak compared to the one used in the literature.

Under these assumptions, we are able to prove that the limiting energy is given by \eqref{eq:functional_H0}, where the limiting surface energy density is given by \eqref{eq:sigma_intro_moving}.
The proof of the $\Gamma$-convergence result requires to face non-trivial technical challenges. On the one hand, we need to prove a uniform bound on the Euclidean length of geodesics of a class of degenerate metrics under very weak assumptions on the metric (see Theorem \ref{existence-of-geodesics-with-bounded-path-length}).
Moreover, we are able to prove that such a bound on the Euclidean length of geodesics fails if our assumptions are not satisfied (see Example \ref{ex:geodesics_infinite_Euclidean_length}).
Finally, in the liminf inequality, we need to \emph{freeze} the first variable of the potential.
This is done by introducing a suitable adjustment of the potential (see \eqref{eq:adjustment}), which allows us to work \emph{as if} the wells were not space-dependent, but only the potential is.
Finally, the construction of the optimal profile requires multiple careful estimates of energies related to changing the optimal profile from point to point.

We also show that the strategies we implement are solid enough to also cover the case where a mass constraint is imposed (see Theorem \ref{thm:Gamma_conv_mass_constraint}).\\

It is the goal of several forthcoming papers to investigate the case where the wells enjoy even lower regularity assumptions.



%%%%%%%%%%%%%%%%%%%%%%%%%%%%%%%%%%%%%%%%%%%%
%%%%%%%%%%%%%%%%%%%%%%%%%%%%%%%%%%%%%%%%%%%%
%%%%%%%%%%%%%%%%%%%%%%%%%%%%%%%%%%%%%%%%%%%%
%%%%%%%%%%%%%%%%%%%%%%%%%%%%%%%%%%%%%%%%%%%%



\section{Assumptions and main results}

In this section we collect all the several assumptions we used in establishing the results of the manuscript.
We will comment each of them, highlighting the reason why we require it, the comparison with similar assumptions in previous works in the literature, and the possibility to weaken it.\\

Let $\Omega\subset\R^N$, where $N\geq 1$, be an open bounded set with Lipschitz boundary.
Consider a function $W \colon \Omega\times\R^M\to[0,\infty)$, with $M\geq 1$, satisfying the following assumptions:
 \begin{enumerate}[label=(H\arabic*), ref=(H\arabic*)]
            \item \label{H1} $W$ is a Carath\'{e}odory function;
            \item \label{H2} There exist functions $a,b\in BV(\Omega;\R^M)\cap L^\infty(\Omega;\R^M)$ such that
            \[
                W(x,u) = 0 \,\,\,\text{ if and only if }\,\,\, u\in\{a(x), b(x)\},
            \]
            for a.e. $x\in\Omega$;
            \item \label{H3} There exists $\delta>0$ such that
            \[
                |a(x) - b(x)| \geq \delta,
            \]
            for a.e. $x\in\Omega$;
            \item \label{H4} There exist a continuous non-decreasing function $f\colon[0,\infty)\to[0,\infty)$ with $f(t)=0$ if and only if $t=0$, and $C_1>0$ such that
            \[
            \frac{1}{C_1} f\left(\min\{ |u-a(x)|, |u-b(x)| \}\right)
            \leq W(x,u)
                \leq C_1 f\left(\min\{ |u-a(x)|, |u-b(x)| \}\right),
            \]
            for a.e. $x\in\Omega$, and all $u\in \R^M$;
            \item \label{H5} There exists $C_2>0$ such that
            \[
            W(x,u) \geq \frac{1}{C_2}|u|,
            \]
            for a.e. $x \in \Omega$ and all $u\in\R^M$ with $|u|\geq C_2$.\\
    \end{enumerate}

We refer to the above assumptions as the \emph{structural} assumptions on the potential $W$.

\begin{remark}
The continuity of the potential $W$ in the second variable is needed in order to make the composition $x\mapsto W(x,u(x))$ measurable for functions $u\in H^1(\Omega;\R^M)$.
It can be relaxed to Borel measurable in the second variable.
Nevertheless, note that for the Gamma-convergence result (see Theorem \ref{thm:Gamma_conv}) we will need the potential to be piecewise continuous.
\end{remark}

\begin{remark}
We decided to work with only two wells just for simplicity of notation and in order to focus on the main novel ideas of the manuscript.
The extension to multiple wells requires a careful definition of the adjustment (see \eqref{eq:adjustment}) and a strong approximation result.
These extensions will be treated in a forthcoming paper.
\end{remark}
 
\begin{remark}
The assumption that the wells are in $L^\infty(\Omega;\R^M)$ is required in order to have \ref{H5}.
Indeed, consider the case where $a,b\in BV(\Omega;\R^M)\setminus L^\infty(\Omega;\R^M)$, and the potential is given by 
\[
W(x,u) \coloneqq \min\{ |u-a(x)|^2, |u-b(x)|^2 \}.
\]
Then, it is not possible to find $C_2>0$ such that \ref{H5} holds uniformly for a.e. $x \in \Omega$.
\end{remark}

\begin{remark}
Assumption \ref{H3} on the separation of wells is used for technical convenience. Indeed, by using techniques similar to those of \cite{Bou} and \cite{CriGra}, it is possible to extend our results, \emph{mutatis mutandis}, to the case where \ref{H3} is dropped.
Note that, in our case, we also consider the case where the wells are allowed to jump. Therefore, a bit of care will be needed in order to characterize the space where the limiting energy is finite.
\end{remark}


\begin{remark}
Condition \ref{H4} is a mild assumption that ensures non-degeneracy of the potential $W$. 
The lower bound is essentially needed in order to ensure compactness (see Theorem \ref{thm:compactness}).
The bounds are needed in order to ensure that geodesics of the optimal profile problem have uniform Euclidean length (see Lemma \ref{lemma:almost_minimizers_have_locality_prop}).
This result is essential in the construction of the recovery sequence, but it is not used in the liminf inequality.
\end{remark}

\begin{remark}
The linear growth at infinity, namely \ref{H5}, is a standard condition in the literature, and it is required to get compactness as well as a uniform $L^\infty$ bound on the optimal profiles. It goes back to the work \cite{FonTar} by Fonseca and Tartar.
In case a mass constraint is imposed, Leoni showed in \cite{Leo} that this growth at infinity can be relaxed to a non-degeneracy of the potential at infinity.
\end{remark}

The above set of assumptions will be the ones that we will need in order to establish the compactness result (see Theorem \ref{thm:compactness}).
We now introduce what we will refer to as the \emph{regularity} assumptions on the wells and on the potential. These will be needed in establishing the $\Gamma$-convergence result (see Theorem \ref{thm:Gamma_conv}).
First, we restrict our attention to a specific subclass of wells and potentials, where their jumps as well as their regularity are controlled.


Here, we use the terminology that the boundary of a set is \emph{piecewise $C^2$} to mean that it is contained in a finite union of $C^2$ hypersurfaces.
\begin{enumerate}[resume, label=(H\arabic*), ref=(H\arabic*)]
            \item \label{H6} There exists an open partition $\Omega_1,\ldots,\Omega_k$ of $\Omega$, where each $\partial\Omega_i \cap \Omega$ is piecewise $C^2$, such that
            \[
            W(x,u) = \sum_{i=1}^k W_i(x,u)\ca_{\Omega_i}(x),
            \]
            \[
            a(x) = \sum_{i=1}^k a_i(x)\ca_{\Omega_i}(x),\quad\quad\quad\quad
            b(x) = \sum_{i=1}^k b_i(x)\ca_{\Omega_i}(x),
            \]
            where $W_i\in C^0(\overline{\Omega}\times\R^M)$ and $a_i, b_i\in W^{1,2}(\Omega;\R^M) \cap C^0(\overline{\Omega};\R^M)$.\\
\end{enumerate}

Next, we introduce a relative continuity condition on the potential.
We define the \emph{adjustment}
$T \colon \Omega\times\R^M\to\R^M$ as
\begin{equation}\label{eq:adjustment}
T(x,w) \coloneqq  \frac{a(x)+b(x)}{2} + \left(b(x)-a(x), v(x)\right)\cdot w,
\end{equation}
where $v(x)\in \R^M\times\R^{M-1}$ is such that
\begin{itemize}
\item[(i)] For each $x\in \Omega$, the matrix
\[
\frac{1}{|b(x)-a(x)|}\left(b(x)-a(x), v(x)\right)
\]
is orthonormal;
\item[(ii)] The map $x\mapsto v(x)$ has the same regularity of the wells $a$ and $b$.
\end{itemize}
The \emph{raison d'\^{e}tre} of the map $T$ is the following: the \emph{adjusted potential}
\[
(x,u) \mapsto W(x, T(x,u))
\]
is such that $W(x, T(x,u))=0$ if and only if $u\in\{\pm e_1\}$ for all $x\in\Omega$.
This follows directly from the fact that $T(x,-e_1)=a(x)$ and $T(x,e_1)=b(x)$.
Namely, the adjustment $T$ allows us to reduce to the case of fixed wells.

We are now in a position to state our next assumption.

\begin{enumerate}[resume, label=(H\arabic*), ref=(H\arabic*)]
            \item \label{H7} There exists a non-decreasing modulus of continuity $\omega\colon[0,\infty)\to[0,\infty)$ such that
            \[
            | W(x, T(x,u)) - W(y, T(y,u)) | \leq \omega(|x-y|)  W(x, T(x,u)),
            \]
            for all $x,y \in \Omega$ and $u\in\R^M$.
\end{enumerate}

\begin{remark}\label{rem:T_symmetric}
In the rest of the paper, we will assume, for simplicity, that
\[
a(x) = -b(x)
\]
for a.e. $x\in\Omega$.
Note that, in this case, assumption \ref{H3} writes as
\[
    |a(x)|\geq \frac{\delta}{2},
\]
for a.e. $x\in\Omega$. For the simplicity of notation, we will write the right-hand side as $\delta$.
Moreover, denoting $a^\perp = v/2$ (cf.\ \eqref{eq:adjustment}) we will use the simplified adjustment
\[
    T_a(w)\coloneqq T(x,w) = ( a(x), a(x)^\perp ) \cdot w,
\]
for each $w\in\R^M$. Note that we can write
\[
    T_a(w) = |a(x)| R_a(x),
\]
where, for each $x\in\Omega$, the map $R_a(x):\R^M\to\R^M$ is orthonormal. In this case, since $a_i \in W^{1,2}(\Omega;\R^M) \cap C^0(\overline{\Omega};\R^M)$, therefore bounded, and $R_a$ is orthonormal, it follows that 
\begin{equation}\label{eq:T_a_properties}
    \| \nabla T_{a_i} \|_{L^2} \leq C_1 \| \nabla a_i \|_{L_2} \leq C_2, \qquad \| T_{a_i} \|_\infty = \| a_i \|_\infty \leq C, \qquad \| T_{a_i}^{-1} \|_\infty = \| a_i \|_\infty^{-1} \leq \frac{1}{\delta}.
\end{equation}
\end{remark}

\begin{remark}
Condition \ref{H7} has been used in similar investigations on models for phase transition (see, for instance \cite{FonPopo}).
Its main role is to allow us to \emph{freeze} the potential at a given point, with an error proportional to the potential itself.
We note that in our case we need to compose the potential with the map $T$ because we consider the case of \emph{moving} wells.
Otherwise, the condition
\[
| W(x, u) - W(y, u) | \leq \omega(|x-y|)  W(x, u),
\]
for all $x,y\in\Omega$ and $u\in\R^M$, would imply that the functions $a$ and $b$ are constant.
\end{remark}

Despite its technical nature, we want to show that assumption \ref{H7} is satisfied by several classes of potentials of interest. Here, we present some of them.

\begin{example}
The first example that we consider is the prototypical case of the potential for the Cahn-Hillard functional with moving wells.
Let $W \colon \Omega\times\R\to[0,\infty)$ be defined as
\[
    W(x,u) \coloneqq (u^2-a^2(x))^2,
\]
where $a\in C^0([-1,1])$, with $a(x)\geq \delta$ for all $x\in [-1,1]$. Then, it holds that
\[
T(x,u) = a(x)u.
\]
Therefore, we have that
\[
W(x,T(x,u)) = a^4(x)(u^2-1)^2.
\]
We now check that \ref{H7} holds:
\begin{align*}
| W(x, T(x,u)) - W(y, T(y,u)) | &= |a^4(x) - a^4(y)| (u^2-1) \\
    &\leq \widetilde{\omega}(|x-y|) (u^2-1)^2 \\
    &\leq \frac{\widetilde{\omega}(|x-y|)}{\delta^4} \min\{a^4(x), a^4(y)\} (u^2-1)^2 \\
    &= \frac{\widetilde{\omega}(|x-y|)}{\delta^4} \min\{ W(x, T(x,u)), W(y, T(y,u)) \},
\end{align*}
where $\widetilde{\omega} \colon [0,\infty)\to[0,\infty)$ is the modulus of continuity of the function $x\mapsto a^4(x)$.
\end{example}

Next, we generalize the above to a similar class of potentials.

\begin{example}
Let $q\in[1,\infty)$.
Define the potential $W\colon\Omega\times\R^M\to[0,\infty)$ as
    \[
        W(x, u) = \min \{ |u - a(x)|^q, |u + a(x)|^q \},
    \]
    where $a \in C^0(\Omega;\R^M)$ with $|a(x)|\geq \delta$, for all $x\in\Omega$.
    Then, we have that
    \[
        W(x, T(x,u)) = |a(x)|^q\min \{ |u - e_1|^q, |u + e_1|^q \}.
    \]
    Therefore, using similar computations as above, we have that
    \[
        |W(x, T(x,u)) - W(y, T(y,u))|
        \leq \frac{2^q}{\delta^p}\widetilde{\omega}(|x - y|)\min \{W(x, T(x,u)), W(y, T(y,u)) \},
    \]
    where $\widetilde{\omega}\colon[0,\infty)\to[0,\infty)$ is the modulus of continuity of the function $x\mapsto |a|^q(x)$.
\end{example}

Next, we show that assumption \ref{H7} is verified even by potentials with very low regularity.

\begin{example}
Let $V\colon\R^M\to[0,\infty)$ be a Borel function, and let $a\colon\Omega\to\R^M$ be a Borel function. Set $b\coloneqq -a$.
Define $W\colon\Omega\times\R^M\to[0,\infty)$ as
\[
W(x,u)\coloneqq V(T(x,u)^{-1}),
\]
where for each $x \in \Omega$, the map $T(x, \cdot)^{-1}$ is the inverse of the map defined in \eqref{eq:adjustment}.
Then, $W$ satisfies assumption \ref{H7}, since
\[
|W(x, T(x,u)) - W(y, T(y,u))|=0
\]
for all $x,y\in\Omega$ and all $u\in\R^M$.
Note that, in order for $W$ to satisfy assumption \ref{H1}, we need $V$ to be continuous.
\end{example}

Finally, we introduce a very mild condition on the function controlling the growth around the wells. This will be needed in order to ensure that the geodesics of the limiting distance function have uniformly bounded Euclidean length (see Theorem \ref{existence-of-geodesics-with-bounded-path-length}).

\begin{enumerate}[resume, label=(H\arabic*), ref=(H\arabic*)]
            \item \label{H8} There exists $C_3>0$ such that
            \[
            f(2t) \leq C_3 f(t),
            \]
            for all $t>0$, where $f\colon[0,\infty)\to[0,\infty)$ is the function given by \ref{H4}.
\end{enumerate}

\begin{remark}
We would like to stress that this assumption is very mild.
Indeed, in previous works, stronger conditions were imposed in order to have a uniform bound on the Euclidean length of certain geodesic problems.
In \cite{CriGra} it was required that the potential behaves as a quadratic function close to the wells, while in \cite{CriFonGan_moving_sub}, that it is controlled by quadratic functions from above and from below close to the wells.

In particular, note that our assumption also covers the case where the potential $W$ is \emph{concave near the wells}, that was excluded by the assumptions in the literature.
Indeed, for any $\alpha>0$, the function $f(t)\coloneqq t^\alpha$ satisfies assumption \ref{H8}.
\end{remark}

Finally, we introduce the family of functionals that we consider.

\begin{definition}
For $\varepsilon>0$, we define the functional $\mathcal{F}_\varepsilon\colon L^1(\Omega;\R^M)\to[0,+\infty]$ as
\[
\mathcal{F}_\varepsilon(u) \coloneqq \int_\Omega \left[ \frac{1}{\varepsilon} W(x,u(x))
    + \varepsilon|\nabla u(x)|^2 \right] \dd x,
\]
for $u\in H^1(\Omega;\R^M)$, and $+\infty$ else in $L^1(\Omega;\R^M)$.
\end{definition}

\begin{definition}
    We define the space $BV(\Omega;\{a,b\})$ as the space of functions $u\in BV(\Omega;\R^M)$ such that $u(x)\in\{ a(x), b(x) \}$ for a.e. $x\in\Omega$.
\end{definition}
	
\begin{definition}
    We denote the \emph{geodesic distance}
    \[
        \dd_W(x;p,q) \coloneqq \inf\left\{ \int_{-1}^1 2\sqrt{W(x,\gamma(t))}|\gamma'(t)| \,:\, \gamma\in W^{1,1}([-1,1];\R^M), \gamma(-1)=p, \gamma(1)=q  \right\},
    \]
    for all $x\in\Omega$ and all $p,q\in\R^M$. Moreover, we set the functional $\mathcal{F}_\infty: L^1(\Omega;\R^M)\to[0,+\infty]$ as
    \[
        \mathcal{F}_\infty(u) \coloneqq \int_{J_u\cap \Omega} \dd_W(x; u^-(x), u^+(x)) \dhno(x),
    \]
    for $u\in BV(\Omega;\{a,b\})$, and $+\infty$ else in $L^1(\Omega;\R^M)$.
    Here, $J_u$ denotes the jump set of the function $u$, $u^-$ and $u^+$ the two traces of $u$ (see Section \ref{sec:BV} for more details).
\end{definition}

We are now in position to state the first two main results of the manuscript. We start with the compactness result, that holds under very mild assumptions on the potential and on the wells.

\begin{theorem}\label{thm:compactness}
Assume \ref{H1}-\ref{H5} hold.
Let $\{\e_n\}_n$ be an infinitesimal sequence.
Let $\{u_n\}_n \subset H^1(\Omega;\R^M)$ be such that
\[
\sup_{n\in\N} \mathcal{F}_{\varepsilon_n}(u_n) < +\infty.
\]
Then, there exists a subsequence (not relabeled) such that $u_{n}\to u$ strongly in $L^1(\Omega;\R^M)$, for some $u\in BV(\Omega; \{a,b\})$.
\end{theorem}
    
We are now in position to state the $\Gamma$-convergence result.

\begin{theorem}\label{thm:Gamma_conv}
Assume \ref{H1}-\ref{H8} hold.
Let $\{\e_n\}_n$ be an infinitesimal sequence.
Then, the sequence of functionals $\{\mathcal{F}_{\varepsilon_n}\}_n$ $\Gamma$-converges to $\mathcal{F}_\infty$ with respect to the $L^1(\Omega;\R^M)$ topology.
\end{theorem}

In particular, the liminf and the limsup inequalities will be given by Theorem \ref{thm:liminf} and Theorem \ref{thm:limsup}, respectively.
We would like to remark here that we are able to prove the liminf inequality under a weaker assumption on the regularity of the partition of $\Omega$.\\

The construction of the recovery sequence we will provide in Theorem \ref{thm:limsup} can be easily adapted to also cover the case where a mass constraint is imposed.

\begin{definition}
For $\varepsilon>0$ and $m\in\R^M$, we define the functional $\mathcal{F}_\varepsilon^m\colon L^1(\Omega;\R^M)\to[0,+\infty]$ as
\begin{gather*}
    \mathcal{F}_\varepsilon^m(u) \coloneqq \begin{cases}
        \displaystyle\int_\Omega \left[ \frac{1}{\varepsilon} W(x,u(x)) + \varepsilon|\nabla u(x)|^2 \right] \dd x \qquad &u \in H^1 (\Omega; \R^M), \; \displaystyle\int_\Omega u \dd x = m,\\
        + \infty \qquad &\text{otherwise}.
    \end{cases}
\end{gather*}
\end{definition}

\begin{definition}
For $m\in\R^M$, we define the functional $\mathcal{F}_\infty^m\colon L^1(\Omega;\R^M)\to[0,+\infty]$ as
\begin{gather*}
    \mathcal{F}_\infty^m (u) \coloneqq \begin{cases}
        \displaystyle\int_{J_u\cap \Omega} \dd_W(x; u^-(x), u^+(x)) \dhno(x) \qquad &u \in BV(\Omega; \{a,b\}), \; \displaystyle\int_\Omega u \dd x = m,\\
        + \infty \qquad &\text{otherwise.}
    \end{cases}
\end{gather*}
\end{definition}

In this case, our main results are stated as follows.

\begin{theorem}\label{thm:Gamma_conv_mass_constraint}
Assume \ref{H1}-\ref{H8} hold, with the function $f$ from $\ref{H5}$ satisfying $f(t) \leq |t|^\alpha$ with $\alpha > 1$ close to the wells.
Let $\{ \e_n\}_n$ be an infinitesimal sequence.
Then, the following hold:
\begin{enumerate}
    \item If $\{ u_n \}_n \subset H^1 (\Omega; \R^M)$ is such that
    $$ \sup_n \mathcal{F}_{\varepsilon_n}^m (u_n) < \infty, $$
    then, up to a subsequence, we have that $u_n \to u$ in $L^1$, with $ u \in BV(\Omega; \{ a, b\})$.
    \item The sequence of functionals $\{ \mathcal{F}_{\varepsilon_n}^m \}_n $ $\Gamma$-converges to $\mathcal{F}_\infty^m$ with respect to the $L^1$ topology.
\end{enumerate}
\end{theorem}

The proof of the recovery sequence for the mass constrained functional requires a slight modification of the proof of Theorem \ref{thm:limsup}, by using a similar strategy to that employed in the proof of \cite[Theorem 15]{cristoferi2025phaseseparationmultiplyperiodic} and will be discussed at the end of Section \ref{sec:limsup}.



%%%%%%%%%%%%%%%%%%%%%%%%%%%%%%%%%%%%%%%%%%%%
%%%%%%%%%%%%%%%%%%%%%%%%%%%%%%%%%%%%%%%%%%%%
%%%%%%%%%%%%%%%%%%%%%%%%%%%%%%%%%%%%%%%%%%%%
%%%%%%%%%%%%%%%%%%%%%%%%%%%%%%%%%%%%%%%%%%%%



\section{Preliminaries}

In this section, we collect some definitions and tools used throughout this paper.

\subsection{Functions of bounded variation and sets of finite perimeter}\label{sec:BV}

We recall some basic definitions and facts about functions of bounded variation and sets of finite perimeter. We refer to \cite{AFP} for more details.

\begin{definition}[Functions of bounded variation] Let $u \in L^1 (\Omega; \R^M)$. We say that $u$ has bounded variation in $\Omega$, denoted by $u \in BV(\Omega; \R^M)$, if its distributional derivative $Du$ is a matrix-valued Radon measure on $\Omega$. In particular, its variation in $\Omega$, denoted by $|Du|(\Omega)$, is equal to 
\begin{equation}
    |Du|(\Omega) \coloneqq \sup \left\{ \sum_{i=1}^M \int_\Omega u_i \, \mathrm{div} \varphi_i \, \dd x : \varphi \in C_c^\infty(\Omega; \R^{M \times N}), \| \varphi \|_\infty \leq 1 \right\}.
\end{equation}
This can also be denoted by $V(u, \Omega)$.
\end{definition}

\begin{definition}[Jump set]
    Let $u \in L^1 (\Omega; \R^M)$. We define the jump set of $u$, denoted by $J_u$, as the set of points where there exist finite one-sided Lebesgue limits. In particular, for each $x \in J_u$ there exists distinct vectors $a,b \in \R^M$ and a normal $\nu \in \mathbb{S}^{N-1}$ such that
    \begin{equation}
        \lim_{\rho \to 0^+} \frac{1}{\rho^N} \int_{B^+ (x, \rho, \nu)} | u(y) - a | \, \dd y = 0, \qquad \lim_{\rho \to 0^+} \frac{1}{\rho^N} \int_{B^- (x, \rho, \nu)} | u(y) - b | \, \dd y = 0,
    \end{equation}
    where $B^+ (x, \rho, \nu)$ and $B^- (x, \rho, \nu)$ are defined as
    \begin{equation*}
        B^+ (x, \rho, \nu) = \{ y \in B(x,\rho) : (y - x) \cdot \nu \geq 0 \}, \quad B^- (x, \rho, \nu) = \{ y \in B(x,\rho) : (y - x) \cdot \nu \leq 0 \}.
    \end{equation*}
    If $x \in J_u$, we denote $(a,b,\nu)$ as $(u^+ (x), u^- (x), \nu_u (x))$.
\end{definition}

\begin{theorem}[Decomposition of distributional derivative]
    The jump set $J_u$ of a function $u \in BV(\Omega; \R^M)$ is countably $\mathcal{H}^{N-1}$-rectifiable. Moreover, 
    \begin{equation}
        Du = \nabla u\, \mathcal{L}^N + ( u^+ - u^- ) \otimes\nu_u \mathcal{H}^{N-1} \llcorner J_u + D^c u,
    \end{equation}
    where $D^c u$ denotes the Cantor part of the distributional derivative.
\end{theorem}

\begin{theorem}[Chain rule in $BV$]
    Let $u \in BV(\Omega; \R^M)$, and let $f : \R^M \to \R^k$ be a Lipschitz continuous function. Then $v \coloneqq f \circ u$ belongs to $BV(\Omega; \R^k)$ and 
    \begin{equation*}
        D v = \nabla f(u) \nabla u\, \mathcal{L}^N + (f(u^+) - f(u^-)) \otimes \nu_u \mathcal{H}^{N-1} \llcorner J_u + \nabla f (\widetilde{u}) D^c u,
    \end{equation*}
    where $\widetilde{u}$ is the precise representative of $u$, defined everywhere except on $J_u$.
\end{theorem}

\begin{definition}[Variation along a direction]
    Let $u \in L^1 (\Omega; \R^M)$ and $\nu \in \mathbb{S}^{N-1}$. We define the \emph{total variation} of $u$ along the direction $\nu$ as 
    \begin{equation*}
        V_\nu (u, \Omega) = \sup \left\{ \sum_{i=1}^M \int_\Omega u_i \frac{\partial \varphi_i}{\partial \nu} \, \dd x : \varphi \in C_c^\infty (\Omega; \R^{M \times N}), \| \varphi \|_\infty \leq 1 \right\}.
    \end{equation*}
\end{definition}

Given now $\nu \in \mathbb{S}^{N-1}$, we denote by $\pi_\nu$ the hyperplane orthogonal to $\nu$, and by $\Omega^\nu$ the projection of $\Omega$ onto $\pi_\nu$. Now, for any $y \in \Omega^\nu$, we define a slice of $\Omega$ as
\begin{equation*}
    \Omega^\nu_y \coloneqq \left\{ t \in \R : y + t \nu \in \Omega \right\}.
\end{equation*}
\begin{figure}[ht]
        \includegraphics[width=0.4\linewidth]{slicing}
    \caption{Example of slicing.}
    \label{fig:slicing}
\end{figure}

Analogously, for $u \colon \Omega \to \R^M$, for any $y \in \Omega^\nu$ we define a slice of $u$ as 
\begin{equation*}
    u_y^\nu \colon \Omega_y^\nu \to \R^M, \qquad u_y^\nu (t) \coloneqq u(y + t \nu).
\end{equation*}

\begin{theorem}[Characterization of $BV$ by sections]
    Let $u \in L^1 (\Omega; \R^M)$ and $\nu \in \mathbb{S}^{N-1}$. Then
    \begin{equation}
        V_\nu (u, \Omega) = \int_{\Omega^\nu} V( u_y^\nu, \Omega_y^\nu) \, \dd y.
    \end{equation}
    Moreover, $u \in BV(\Omega; \R^M)$ if and only if there exist $N$ linearly independent unit vectors $\nu_1, .., \nu_N$ such that $u_y^{\nu_j} \in BV(\Omega_y^{\nu_j}; 
    \R^M)$ for $\mathcal{H}^{N-1}$-a.e. $y \in \Omega^{\nu_j}$, and
    \begin{equation}
        \int_{\Omega^{\nu_j}} V (u_y^{\nu_j}, \Omega_y^{\nu_j}) \, \dd y < \infty \qquad \forall j = 1, \ldots, N.
    \end{equation}
\end{theorem}

We focus now on the functions of bounded variations which are characteristic functions of sets.

\begin{definition}[Finite perimeter set]
    Let $E \subseteq \Omega$. We say that $E$ has \emph{finite perimeter in $\Omega$} if its characteristic function $\mathbbm{1}_E \colon \Omega \to \{0,1\}$ has bounded variation on $\Omega$. Equivalently, we can say $E$ has finite perimeter in $\Omega$ if
    \begin{equation*}
        \mathrm{Per}(E; \Omega) \coloneqq \sup \left\{  \int_E \mathrm{div} \varphi \, \dd x : \varphi \in C_c^\infty (\Omega; \R^N), \| \varphi \|_\infty \leq 1 \right\} < +\infty.
    \end{equation*}
\end{definition}

\begin{definition}[Reduced boundary]
    Let $E \subseteq \Omega$ be a set of finite perimeter. We define its \emph{reduced boundary} $\partial^* \!E$ as the subset of $\partial E$ for which the limit 
    \begin{equation*}
        \nu_E (x) \coloneqq - \lim_{r \to 0^+} \frac{D\mathbbm{1}_E (B(x,r))}{| D \mathbbm{1}_E | (B(x,r))},
    \end{equation*}
    exists and it is such that $| \nu_E (x) | = 1$. The vector $\nu_E (x)$ is called \emph{measure theoretic exterior normal} to $E$.
\end{definition}

\begin{theorem}[De Giorgi's structure theorem]
    Let $E \subseteq \Omega $ be a set of finite perimeter. Then $\partial^* \! E$ is countably $\mathcal{H}^{N-1}$-rectifiable and
    \begin{equation}
        \mathrm{Per}(E; \Omega) = \mathcal{H}^{N-1}(\partial^* \! E \cap \Omega).
    \end{equation}
    Moreover, for $x_0 \in \partial^* \! E$ and $r > 0$, by defining the blow-up set $E_r$ and the half-plane $H$ as
    \begin{equation*}
        E_r \coloneqq \frac{E - x_0}{r}, \qquad H \coloneqq \{x \in \R^N : x \cdot \nu_E (x) \geq 0 \},
    \end{equation*}
    we get that
    $$ E_r \to H \quad \text{in } L^1_\text{loc} (\Omega), \qquad \lim_{r \to 0^+} \frac{\mathcal{H}^{N-1}( \partial^* \! E \cap Q(x_0,r))}{r^{N-1}} = 1. $$
\end{theorem}

We conclude this subsection by stating the coarea formula for $(N-1)$-rectifiable sets in $\R^N$, suitably modified for our setting. We refer to \cite[Theorem 2.93]{AFP} for the general result.
\begin{theorem}[Coarea]
    Let $E \subset \R^N$ be a countably $\mathcal{H}^{N-1}$-rectifiable set. Let $f \colon E \to \R$ be a Lipschitz continuous function, and let $g : E \to [0,+\infty]$ be a Borel function. Then
    \begin{equation}
        \int_E g(x) | \nabla^E f(x) | \, \dd \mathcal{H}^{N-1} (x) = \int_{\R} \left( \int_{E \cap \{ f = t \}} g(y) \, \dd \mathcal{H}^{N-2} (y) \right) \dd t,
    \end{equation}
    where $\nabla^E f(x)$ represents the tangential gradient of $f$ with respect to the approximate tangent space $T_x E$.
\end{theorem}

Finally, we recall that sets of finite perimeter can be approximated strongly thanks to a result proved by de Gromard in \cite{DeGrom}.

\begin{theorem}\label{thm:strong_BV}
Let $A\subset \R^M$ be an open set, and let $E\subset A$ be a set of finite perimeter in $A$.
Then, for each $\varepsilon>0$ there exist a set $F\subset A$ of finite perimeter in $A$, and a compact set $C\subset A$ such that the following holds:
\begin{itemize}
\item[(i)] $\partial F\cap A$ is contained in a finite union of $C^1$ hypersurfaces;
\item[(ii)] $\| \ca_{E} - \ca_{F} \|_{BV(A)}<\varepsilon$;
\item[(iii)] $\hno(\partial F\cap A\setminus \partial^* E)<\varepsilon$;
\item[(iv)] $F\subset E+B(0,\varepsilon)$, and
    $D\setminus F\subset (A\setminus E)+B(0,\varepsilon)$;
\item[(v)] $C \subset A\cap \partial^* E \cap \partial F$;
\item[(vi)] $\nu_E(x) = \nu_F(x)$ for all $x\in C$;
\item[(vii)] $|D\ca_E|(D\setminus C)<\varepsilon$.
\end{itemize}
\end{theorem}



%%%%%%%%%%%%%%%%%%%%%%%%%%%%%%%%%%%%%%%%%%%%
%%%%%%%%%%%%%%%%%%%%%%%%%%%%%%%%%%%%%%%%%%%%
%%%%%%%%%%%%%%%%%%%%%%%%%%%%%%%%%%%%%%%%%%%%
%%%%%%%%%%%%%%%%%%%%%%%%%%%%%%%%%%%%%%%%%%%%



\subsection{$\Gamma$-convergence}

    This section provides a brief overview of the definition and fundamental properties of $\Gamma$-convergence. Given the setting of this paper, we define $\Gamma$-convergence using sequences. For a comprehensive treatment of $\Gamma$-convergence in general topological spaces, we refer the reader to \cite{Dalmasobook} (see also \cite{Braides}).

    \begin{definition}
        Let $(X, \mathrm{d})$ be a metric space. A sequence of functionals $F_n: X \to [-\infty, +\infty]$ is said to \emph{$\Gamma$-converge} to a functional $F: X \to [-\infty, +\infty]$ \emph{with respect to the metric $\mathrm{d}$}, denoted by $F_n \overset{\Gamma-d}{\longrightarrow} F$, if the following two conditions are satisfied:
        \begin{itemize}
            \item[(i)] (Liminf inequality) For every $x \in X$ and for every sequence $\{x_n\}_n \subset X$ converging to $x$, the following inequality holds:
            \[
                F(x) \leq \liminf_{n\to\infty} F_n(x_n).
            \]
            \item[(ii)] (Limsup inequality) For every $x \in X$, there exists a sequence $\{x_n\}_n \subset X$ such that $x_n \to x$ and
            \[
                \limsup_{n\to\infty} F_n(x_n) \leq F(x).
            \]
        \end{itemize}
    \end{definition}

    The concept of $\Gamma$-convergence was specifically developed to provide a variational characterization of the asymptotic behavior of minimization problems. It ensures the convergence of both global minimizers and minimum values (refer to \cite[Corollary 7.20]{Dalmasobook}).

    \begin{theorem}
        Let $(X, \mathrm{d})$ be a metric space. Consider a sequence of functionals $F_n: X \to \mathbb{R} \cup \{\infty\}$ that $\Gamma$-converges to a functional $F: X \to \mathbb{R} \cup \{\infty\}$.
        Suppose that for each $n \in \mathbb{N}$, $x_n \in X$ is a minimizer of $F_n$.
        Then, any cluster point $x \in X$ of the sequence $\{x_n\}_n$ is a minimizer of $F$, and satisfies
        \[
            F(x) = \limsup_{n\to\infty} F_n(x_n).
        \]
        Moreover, if the sequence $\{x_n\}_n$ converges to $x$, the limsup above is a full limit.
    \end{theorem}



%%%%%%%%%%%%%%%%%%%%%%%%%%%%%%%%%%%%%%%%%%%%
%%%%%%%%%%%%%%%%%%%%%%%%%%%%%%%%%%%%%%%%%%%%
%%%%%%%%%%%%%%%%%%%%%%%%%%%%%%%%%%%%%%%%%%%%
%%%%%%%%%%%%%%%%%%%%%%%%%%%%%%%%%%%%%%%%%%%%


\subsection{Young measures}
Young measures are an important tool in the theory of calculus of variations. Here, we recap the definitions and the compactness property that we use. For a more general introduction to Young measures, we refer to \cite{RindlerCalcVarBook}.
\begin{definition}
    Let $\Omega \subset \R^N$ open. We call a family of probability measures $\{\nu_x\}_{x \in \Omega} \subset \mathcal P(\R^M)$ a \emph{Young measure} if
    \begin{itemize}
        \item $(x \mapsto \nu_x)$ is weakly$^*$ measurable,
        \item we have 
        \[
            \int_\Omega  \int_{\R^M} |p| \, \nu_x(p) \, \mathrm d x < \infty.
        \]
    \end{itemize}
\end{definition}
Here, we are interested in Young measures generated by $L^1$ functions.
\begin{definition}
    Let $\Omega \subset \R^N$ open. We say that a Young measure $\{\nu_x\}_{x \in \Omega}$ is \emph{generated} by a sequence of $\{u_n \}_n \subset L^1(\Omega; \R^M)$ functions if $\{\delta_{u(x)}\}_{x \in \Omega}$ converges weakly$^*$ to $\{\nu_x\}_{x \in \Omega}$ in $L^\infty_{w^*}(\Omega; \M_b(\R^M))$.
\end{definition}
One of the features of Young measures are the compactness properties. More specifically, we will use the following result (cf.\ \cite[Lemma 4.3]{RindlerCalcVarBook}:
\begin{theorem}\label{thm:compactness_YM}
    Let $\Omega \subset \R^N$ open. Let $\{u_n\}$ be an bounded sequence in $L^1(\Omega; \R^M)$. Then, there exists a Young measure $\nu$ such that $u_n$ generates $\nu$ (after extracting a subsequence). Furthermore, if $\{f(\cdot , u_n(\cdot))\}_{n \in \N} \subset L^1(\Omega; \R)$ is a bounded and equiintegrable sequence in $L^1(\Omega; \R)$ for a Carath\'eodory function $f: \Omega \times \R^M \to \R$ then 
    \[
        \int_\Omega f(x, u_n(x)) \, \mathrm d x \to \int_\Omega \int_{\R^M} f(x, p) \, \mathrm d \nu_x(p) \, \mathrm d x.
    \]
\end{theorem}


%%%%%%%%%%%%%%%%%%%%%%%%%%%%%%%%%%%%%%%%%%%%
%%%%%%%%%%%%%%%%%%%%%%%%%%%%%%%%%%%%%%%%%%%%
%%%%%%%%%%%%%%%%%%%%%%%%%%%%%%%%%%%%%%%%%%%%
%%%%%%%%%%%%%%%%%%%%%%%%%%%%%%%%%%%%%%%%%%%%



\section{Technical results}

In this section we collect the technical results needed to prove the main result of the manuscript. We decided to present them in here because they are also of a separate interest for the reader.


%%%%%%%%%%%%%%%%%%%%%%%%%%%%%%%%%%%%%%%%%%%%
%%%%%%%%%%%%%%%%%%%%%%%%%%%%%%%%%%%%%%%%%%%%
%%%%%%%%%%%%%%%%%%%%%%%%%%%%%%%%%%%%%%%%%%%%
%%%%%%%%%%%%%%%%%%%%%%%%%%%%%%%%%%%%%%%%%%%%


\subsection{Properties of the geodesic distance}

In this subsection, we discuss properties of the geodesic distance. More precisely, given a non-negative Borel function $W_0: \R^M \to [0, \infty)$ we associate to $W_0$ a geodesic distance via 
\[
    \dd_{W_0}(p,q) := \inf\left\{ \int_{-1}^1 W_0(\varphi)|\varphi'| \, \mathrm{d}s :\, \varphi \in W^{1,1}\left([-1,1]; \R^M \right) \right\},
\]
for all $p,q \in \R^M$.
The typical choice of $W_0$ in the case of phase transitions is $W_0 = 2\sqrt W$ where $W$ is the double-well potential of the Cahn--Hilliard functional. If $W_0$ is such that $W_0 \geq \delta >0$, lower semi-continuous and coercive, then it is easy to see that the infimum is in fact a minimum and that the corresponding geodesics have bounded path length. If we allow $W_0$ to be $0$ at certain points this does not hold in general. Even if there exist geodesics, their path length does not necessarily have to be bounded (see Example \ref{ex:geodesics_infinite_Euclidean_length}). The main result of this subsection provides a general condition for this to hold true. The setting we assume throughout the section is the existence of a continuous function $f\colon [0,\infty) \to [0, \infty)$ with $f(0) = 0$ and $f > 0$ on $(0,\infty)$ such that certain growth conditions hold true. In what follows, $a,b \in \R^M$ are always fixed single pointed wells of our potential, i.e., we assume 
$$
    \{ W_0 = 0\} = \{a,b\}.
$$
We introduce a few different growth conditions:   

\begin{enumerate}[label=(G\arabic*), ref=(G\arabic*)]
    \item\label{G1} \emph{Controlled upper growth}: There exists a constant $C_G > 0$ such that 
    \[
        W_0(u) \leq C_Gf(\min \{|u - a|, |u - b|\}),  
    \]
    for all $u \in \R^M$.
    \item\label{G2} \emph{Summable growth around the wells}: There exist $\alpha, \, R, \, C_G > 0$ such that $f$ fulfils 
    \begin{enumerate}
        \item $f$ is non-decreasing in $[0,R]$;
        \item $f(2t) \leq C_Gf(t)$ for all $t \in [0, \infty)$;
        \item for all $u \in B_R(a)\cup B_R(b)$ it holds that
        \[
            \frac{1}{C_G}f(\min \{|u - a|, |u - b|\}) \leq W_0(u) \leq C_Gf(\min \{|u - a|, |u - b|\}).
        \]
    \end{enumerate}
    \item\label{G3} \emph{Non-summable growth at $\infty$}: We have 
    \[
        \int_0^\infty f(t) \, dt = \infty,
    \]
    and there exists a constant $C_G > 0$ such that 
    \[
        \frac{1}{C_G}f(|u|) \leq W_0(u),
    \]
    for all $u \in \R^M$.
    \item\label{G4} \emph{Strong non-summable growth at $\infty$}: There exist $\alpha, R > 0$ such that 
    \[
        f(t) \geq \alpha
    \]
    for $t \geq R$,
    and there exists a constant $C_G > 0$ such that 
    \[
        \frac{1}{C_G}f(|u|) - C_G \leq W_0(u),
    \]
    for all $u \in \R^M$.
\end{enumerate}
To simplify the notation in the proofs, we assume 
$$ f(1) < 1/2,\qquad \alpha = R = 1,\qquad  |a - b| = 2.$$
We start the discussion with the following elementary lemma.

\begin{lemma}\label{arc-length-estimate}
    Let $p, q \in \R^2$, with polar coordinates $p = r_pe^{i\psi_p}, q = r_qe^{i\psi_q}$ with $\psi_p, \psi_q \in [0,2\pi)$ and $p_q, r_q \geq 0$. Consider $r(s) := sr_q + (1-s)r_p$ and $\psi(s) := s\psi_q + (1-s)\psi_p$ and set
    \[
        \gamma(s) := r(s)e^{i\psi(s)}.
    \]
    Then,
    \[
        \int_0^1|\gamma'|\, \mathrm{d}s \leq |r_p - r_q| + 2\pi\max\{r_p, r_q\}.
    \]
\end{lemma}

\begin{proof}
    Let us consider the situation graphically.
    \begin{figure}[H]
        \includegraphics[width=0.25\linewidth]{polar}
    \end{figure}
    \noindent We directly compute:
    \[
        \int_0^1 | \gamma ' | \, \mathrm d s \leq \int_0^1 |r_p - r_q| + | r(s) (\psi_p - \psi_q) | \, \mathrm d s.
    \]
    From which we get the required bound.
\end{proof}

In \cite{ZunSter} it has been proved that if $W_0$ is continuous, then the minimization problem defining $\dd_{W_0} (p,q)$ admits a solution for all $p,q \in \R^M$.
Since minimizers might not exist if $W_0$ is only Borel, we use the concept of $\eps$-minimizers.

\begin{definition}
    Let $\varepsilon > 0$. A curve $\varphi \in W^{1,1}([-1,1]; \R^M)$ is called an \emph{$\eps$-minimizer} with respect to the geodesic distance $\dd_W(\varphi(-1), \varphi(1))$ if 
    \begin{align}\label{def:eps-minimizer}
        \int_{-1}^1 W_0(\varphi)|\varphi'|\,\mathrm{d}s \leq \dd_{W_0}(\varphi(-1), \varphi(1)) + \varepsilon.
    \end{align}
    Since the left-hand side is invariant under reparameterization we will also call any curve $\psi \in W^{1,1}([c,d]; \R^M)$ an $\eps$-minimizer if there exists a reparameterization to a curve $\varphi \in W^{1,1}([-1,1]; \R^M)$ such that \eqref{def:eps-minimizer} holds.
\end{definition}

Before stating the main result of this subsection, we want to recap some well-known properties of the geodesic distance and its $\eps$-minimizers. We give some proofs for convenience of the reader.

\begin{lemma}\label{lemma:almost_minimizers_have_locality_prop}
    Let $W_0: \R^M \to [0,\infty)$ be a locally bounded Borel function, and let $\eps > 0$. Suppose that $\varphi \in W^{1,1}([-1,1], \R^M)$ is such that
    \[
        \int_{-1}^1 W_0(\varphi)|\varphi'|\,\mathrm{d}s \leq d_{W_0}(\varphi(-1), \varphi(1)) + \eps.
    \]
    Then,  
    \[
        \int_{t_1}^{t_2} W_0(\varphi)|\varphi'|\,\mathrm{d}s \leq d_{W_0}(\varphi(t_1), \varphi(t_2)) + \eps
    \]
    for all $t_1 \leq t_2 \in [-1,1]$. 
\end{lemma}
\begin{proof}
    Suppose that there exist $t_1, t_2 \in [-1,1]$ with $t_1 \leq t_2$ such that
    \[
        \int_{t_1}^{t_2} W_0(\varphi)|\varphi'|\,\mathrm{d}s > \dd_{W_0}(\varphi(t_1), \varphi(t_2)) + \eps.
    \]
    This directly implies 
    \begin{align*}
        \int_{-1}^1 W_0(\varphi)|\varphi'|\,\mathrm{d}s &= \int_{-1}^{t_1} W_0(\varphi)|\varphi'|\,\mathrm{d}s + \int_{t_1}^{t_2} W_0(\varphi)|\varphi'|\,\mathrm{d}s + \int_{t_2}^{1} W_0(\varphi)|\varphi'|\,\mathrm{d}s \\
        & > \dd_{W_0}(\varphi(-1), \varphi(t_1))+ (\dd_{W_0}(\varphi(t_1), \varphi(t_2)) + \eps) + \dd_{W_0}(\varphi(t_2), \varphi(1)) \\
        & \geq \dd_{W_0}(\varphi(-1), \varphi(1)) + \eps.
    \end{align*}
    This would contradict that $\varphi$ was an $\eps$-minimizer with respect to $\dd_W(\varphi(-1), \varphi(1))$.
\end{proof}

The next lemma states that even if the function $W_0$ is not continuous, the corresponding geodesic distance is still Lipschitz continuous.

\begin{lemma}\label{lemma:geodesic_distance_is_locally_lipschitz}
    Suppose that $W_0: \R^M \to [0, \infty)$ is a locally bounded Borel function. Then, $\dd_{W_0}$ is locally Lipschitz. \\
    Moreover, for any bounded convex set $K \subset \R^M$ the Lipschitz constant $L_K$ is bounded by $\|W_0\|_{L^\infty(K)}$. \\
    In particular, if \ref{G1} holds there exists a constant $C > 0$ (only dependent on $|a|$, $|b|$, $f$ and $C_G$) such that for any ball $B_r(0)$ we have 
    \[
        L_{\overline{B_r(0)}} \leq C(1 + f(r)).
    \]
\end{lemma}
The proof can be found, for instance, in \cite[Lemma 2.7]{DeutschSolidSolid}. It well-known that $\eps$-minimizers with respect to the geodesic distance are uniformly bounded if the potential is coercive. Here, we show that, in fact, the weaker condition of a non-summable growth condition at $\infty$ suffices.
\begin{lemma}\label{lemma:eps_minimizers_are_bounded}
    Suppose that $W_0: \R^M \to [0, \infty)$ is a locally bounded Borel function, and $\ref{G3}$ holds. Moreover, let $\eps \in (0,1)$ and $r > 0$. Then, there exists a constant $K > 0$ (only dependent on $|a|$, $|b|$, $r$, $f$ and $C_G$) such that for any $p, q \in B_r(0)$ and any $\eps$-minimizer $\varphi$ with respect to $\dd_{W_0}(p,q)$ we have 
    \[
        \| \varphi \|_{L^\infty([-1,1];\R^M)} \leq K. 
    \]
\end{lemma}
\begin{proof}
    By \ref{G3} there exist a partition $r \leq t_0 < ... < t_n$ such that 
    \begin{align}\label{eq:correct_partition_choice}
        \sum_{k = 1}^n \min_{t \in [t_{i - 1}, t_i]}f(t) (t_i - t_{i - 1}) > 2C_G\left(\max_{p,q \in B_r(0)}\{\dd_{W_0}(p,q)\} + 1 \right)
    \end{align}
    We argue by contradiction: Suppose that there exist $\tilde p, \tilde q \in B_r(0)$ such that $$\| \varphi \|_{L^\infty([-1,1];\,  \R^N)} > t_n$$ for an $\eps$-minimizer $\varphi$ of $\dd_{W_0}(\tilde p, \tilde q)$. Then, there exist $s_0 < ... < s_n$ such that $f(s_i) = t_i$ for $i = 0,.., n$. In particular, we have
    \begin{align*}
        \int_{-1}^1 W_0(\varphi)|\varphi'|\,\mathrm{d}s &\geq \sum_{k = 1}^n \int_{s_{i-1}}^{s_i} W_0(\varphi)|\varphi'|\,\mathrm{d}s \\ 
        &\geq \frac{1}{C_G}\sum_{k = 1}^n \min_{t \in [t_{i - 1}, t_i]}f(t) (t_i - t_{i - 1}) \\
        &\geq 2\left(\max_{p,q \in B_r(0)}\{\dd_{W_0}(\tilde p, \tilde q)\} + 1 \right) \\
        &> \dd_{W_0}(\tilde p, \tilde q) + \eps.
    \end{align*}
    This contradicts $\varphi$ being an $\eps$-minimizer.
\end{proof}

\begin{remark}\label{remark:geodesic_parameter_discussion_1}
    The constant in Lemma \ref{lemma:geodesic_distance_is_locally_lipschitz} can be made explicit in the following sense. Since $f$ is positive in $(0,\infty)$, the function 
    $$F(t) = \int_r^T f(t) \, \mathrm \dd t$$
    is non-decreasing and admits a continuous, non-decreasing left inverse $G\colon(0,\infty) \to (r, \infty)$. In particular, we have 
    \[
        F(K) = 3\left(\max_{p,q \in B_r(0)}\{\dd_{W_0}(\tilde p, \tilde q)\} + 1 \right),
    \]
    which we can express via the left inverse as
    \begin{align}\label{eq:expression_of_constant}
        K = G\left(3\left(\max_{p,q \in B_r(0)}\{\dd_{W_0}(\tilde p, \tilde q)\} + 1 \right)\right).
    \end{align}
    This has the following consequence: if \ref{G1} and \ref{G4} hold, repeating the proof of Lemma \ref{lemma:eps_minimizers_are_bounded} with $\tilde f = \chi_{[R, \infty)}\alpha$ we obtain that the resulting $K$ is then just an affine combination of $R$ (from \ref{G4}) and $f(r)$. Indeed, this can also be derived by choosing the partition $t_0 = R$ and $t_1 = K := R + C(1 + f(r))$ in \eqref{eq:correct_partition_choice} for a suitable large constant $C > 0$. In particular, we have
    \[
        K \leq C(1 + R + f(r)) 
    \]
    where $C$ is only dependent on $|a|$, $|b|$, $C_G$ and $\alpha$. 
\end{remark}

\begin{remark}\label{remark:geodesic_parameter_discussion_2}
    A consequence of Remark \ref{remark:geodesic_parameter_discussion_1} is the following. Let $f, g \in C([0,\infty), [0,\infty))$ with $f \leq g$ both satisfying $\ref{G3}$, in the sense that they fulfill the non-summability condition at $\infty$. Then, Lemma \ref{lemma:eps_minimizers_are_bounded} now gives constants $K_f$ and $K_g$. By \eqref{eq:expression_of_constant}, we can directly see that we can choose the constants such that $K_f \leq K_g$ holds.
\end{remark}

Next, we will observe that there exists a minimising sequence of curves $\varphi_m$ for the geodesic distance which has bounded path-length.

\begin{theorem}\label{existence-of-geodesics-with-bounded-path-length}
    Suppose that $W_0$ fulfils \ref{G1}, \ref{G2} and \ref{G3}. Let $r > 0$ and let $p, q \in B_r(0)$. Then, there exists a sequence $\{ \varphi_m \}_m \subset W^{1,1}([-1,1]; \R^M)$ with $\varphi(-1) = p$, $\varphi(1) = q$ and 
    \[
        \int_{-1}^1 W_0(\varphi_m)|\varphi_m'| \,\mathrm{d}s \leq \dd_{W_0}(p, q) + C\frac{1}{2^m}
    \]
    such that
    \begin{equation}\label{eq:geodesic_bounded_Euclidean}
        \int_{-1}^1 |\varphi_m'| \,\mathrm{d}s \leq C(1 + f(|p|)|p| + f(|q|)|q|),
    \end{equation}
    with a constant $C >0 $ only depending on $|a|$, $|b|$, $r$ and $C_G$. If, in addition, \ref{G4} holds, then $C$ is only dependent on $|a|$, $|b|$, $\alpha$, $R$ and $C_G$, but independent of $r$.
    
    As a consequence, the infimization problem defining $\dd_{W_0}(p, q)$ admits a solution if $W_0$ is lower-semicontinuous.
\end{theorem}

\begin{figure}
    \includegraphics[width=0.8\linewidth]{geodesic}
    \caption{The substitution of $\tilde \varphi_m$ in Theorem \ref{existence-of-geodesics-with-bounded-path-length}.}
    \label{fig:geodesic}
\end{figure}

\begin{proof} We divide the proof in steps.\\[0.5em]
    \textbf{Step 0: The setup.}
    Let $\{ \tilde \varphi_m \}_m$ be a sequence such that
    \[
       \int_{-1}^1 W_0(\tilde \varphi_m)|\tilde \varphi_m'| \,\mathrm{d}s \leq \dd_{W_0}(p, q) + \frac{d_{W_0}(a,b)}{2^{2m}}, 
    \]
    for every $m \in \N$. Define the following sets for $k \in \N$ and $\delta > 0$:
    \[
        \T^{1,k} := \left\{ u \in \R^M: \frac{\delta}{2^{k+1}} \leq |u-a| \leq \frac{\delta}{2^{k}}\right\},
    \]
    \[
        \T^{2,k} := \left\{ u \in \R^M: \frac{\delta}{2^{k+1}} \leq |u-b| \leq \frac{\delta}{2^{k}}\right\}.
    \]
    For $l = 1,2$ set 
    \[
        t_0^l := \inf \, (\tilde \varphi_m)^{-1} \left(\bigcup_{k \geq m}\T^{l,k}\right) \quad \text{ and } \quad
        t_1^l := \sup \, (\tilde \varphi_m)^{-1} \left(\bigcup_{k \geq m}\T^{l,k}\right).
    \]
    Here we point out that these values need not be well-defined if $\tilde \varphi_m$ does not pass through the corresponding regions. In what follows, without loss of generality, we assume that those values are well-defined since the analysis in the other cases is carried out analogously. By the local optimality of $\eps$-minimizers for the geodesic distance (cf.\ Lemma \ref{lemma:almost_minimizers_have_locality_prop}), we know that for any $r, t \in [-1,1]$ with $\varphi_m(r) \in B_\delta(a)$ and $\varphi_m(t) \in B_\delta(b)$
    \[
        \left|\int_{r}^tW_0(\tilde \varphi_m)|\tilde \varphi_m'|\, \mathrm ds \right|  \leq d_{W_0}(a,b) + \omega(\delta) + \frac{d_W(a,b)}{2^{2m}}
    \]
    for some modulus of continuity $\omega$. In particular, for small $\delta > 0$, we get 
    \[
        \left|\int_{s}^tW_0(\tilde \varphi_m)|\tilde \varphi_m'|\, \mathrm ds \right| < \frac{5}{2}d_W(a,b) < 3d_W(a,b).
    \]
    This implies 
    \[
        \sup \, (\tilde \varphi_m)^{-1} (B_\delta(a)) < \inf \, (\tilde \varphi_m)^{-1} (B_\delta(b))
    \]
    or 
    \[
        \sup \, (\tilde \varphi_m)^{-1} (B_\delta(b)) < \inf \, (\tilde \varphi_m)^{-1} (B_\delta(a)).
    \]
    We assume that we are in the first case. The idea of our strategy is the following. In this case, we immediately infer $t_1^1 \leq t_0^2$ for a fixed small $\delta  > 0$ (independent of $m$) (cf.\ \cite[Lemma 2.8]{DeutschSolidSolid}). Now, in what follows, we will substitute $\tilde \varphi_m$ with a straight path in $\bigcup_{k \geq m}\T^{l,k}$. This substitution increases the path length and the length with respect to $W_0$ only by a negligible amount. Then, we estimate the path length of $\tilde \varphi_m$ in each $\T^{l,k}$ separately by finding a suitable competitor for the geodesic distance in this set. Lastly, we estimate the path length of the remaining curve. \\[-0.5em]
    
    \textbf{Step 1: Substituting $\tilde \varphi_m$ around the wells.} Consider the curve 
    \[
        \varphi_m := \begin{cases}
            \theta_{1,m} & \text{ in } \text{conv}\left((\tilde \varphi_m)^{-1}(\bigcup_{k \geq m}\T^{1,k})\right), \\
            \theta_{2,m} & \text{ in } \text{conv}\left((\tilde \varphi_m)^{-1}(\bigcup_{k \geq m}\T^{2,k})\right), \\
            \tilde \varphi_m & \text{ otherwise},
        \end{cases}
    \]
    where $\text{conv}(A)$ denotes the convex hull of a set $A \subset \R^M$, and $\theta_{1,m}$, $\theta_{2,m}$ are defined as follows.
    Write $\tilde \varphi_m(t_0^1) = P_1 + a$ and $\tilde \varphi_m(t_1^1) = P_2 + a$ with $P_1, P_2 \in \R^M$. Now, define $\theta_{1,m}$ as the straight line from $P_1$ to $P_2$
    \[
        \theta_{1,m}(t) := \frac{t_1^1 - t}{t_1^1 - t_0^1}P_1 + \frac{t - t_0^1}{t_1^1 - t_0^1}P_2 + a.
    \]
    Analogously, we define $\theta_{2,m}$.
    Notice now that by \ref{G2}
    \begin{align*}
        \int_{t_0^1}^{t_1^1} {W_0(\theta_{1,m})}|\theta_{1,m}'| \,\mathrm{d}s &\leq Cf\left( \frac{\delta}{2^m} \right)|P_1 - P_2|
        \leq \frac{C}{2^m},
    \end{align*}
    where we used $|P_1 - P_2| \leq (1/2)^{m-1}$ and $f(1) < 1/2$ in the last line. Similarly, we infer 
    \begin{align*}
        \int_{t_0^1}^{t_1^1} {W_0(\theta_{2,m})}|\theta_{2,m}'| \,\mathrm{d}s &\leq \frac{C}{2^m}.
    \end{align*}
    In particular, this substitution only increases the length with respect to $W_0$ by a small amount, which means that for large $m$ these are $1/2$-minimizers of the geodesic distance. Applying Lemma \ref{lemma:eps_minimizers_are_bounded}, we get a constant $K > 0$  such that 
    \begin{align}\label{eq:uniform_estimate_sequence}
        \|\varphi_m\|_{L^\infty} \leq K.
    \end{align}
    We now want to estimate the path length of $\varphi_m$ in the interval $[-1, t^1_0]$. Assume $t^1_0 > -1$, otherwise there is nothing to estimate. Now, we define inductively a decreasing sequence $s_0, s_1, .., s_{m}$ the following way: Set
    \[
        s_0 := t^1_0,
    \]
    and, for $k = 1, .., m$
    \[
        s_{k} := 
            \inf_m (\tilde \varphi_m)^{-1}(\T^{1, m - k}) \cap [-1, s_{k - 1}].
    \]
    Notice that these values need not be well defined for all $k$. Indeed, we denote the last index where it is well defined with $k_0$. Notice that $s_{k_0} = -1$ has to hold if $k_0 < m$. \\[-0.5em]
    
    \textbf{Step 2: Estimating the path length of $\varphi_m|_{[s_{k+1}, s_k]}$.} 
    Notice first, that by construction and by \ref{G2}
    \[
        W_0(\tilde \varphi_m)|_{[s_{k + 1}, s_{k}]} \geq f\left(\frac{\delta}{2^{m - k}}\right).
    \]
    Now, we use the local optimality property of Lemma \ref{lemma:almost_minimizers_have_locality_prop}, i.e., we have for every $k \in \{0,.., k_0 - 1\}$ 
    \[
        \int_{s_{k + 1}}^{s_{k}} W_0(\tilde \varphi_m)|\tilde \varphi_m'| \,\mathrm{d}s \leq \dd_{W_0}\left(\tilde \varphi_m(s_{k + 1}), \tilde \varphi_m(s_{k })\right) + \frac{C}{2^{2m}}.
    \]
    Consider the circular arc $\zeta^k_m$ connecting $\tilde \varphi_m(s_{k+1})$ and $\tilde \varphi_m(s_k)$. More precisely, write $\tilde \varphi_m(s_{k+1}) = P_1 + a$ and $\tilde \varphi_m(s_{k}) = P_2 + a$ with corresponding $P_1, P_2 \in \R^M$. Then, let $\gamma$ be the circular arc from Lemma \ref{arc-length-estimate} connecting $P_1$ and $P_2$ in the plane spanned by $\{0, P_1, P_2\}$ which we reparameterise to the interval $[s_{k+1}, s_k]$. Now, set 
    \[
        \zeta^k_m := \gamma + a.
    \]
    We have 
    \begin{align*}
        f\left(\frac{\delta}{2^{m - k}}\right) \int_{s_{k + 1}}^{s_{k}}|(\tilde \varphi_m)'|\, \mathrm{d}s 
        &\leq \int_{s_{k + 1}}^{s_{k}} W_0(\tilde \varphi_m)|(\tilde \varphi_m)'|\, \mathrm{d}s \\
        &\leq \dd_{W_0}\left(\tilde \varphi_m(s_{k + 1}), \tilde \varphi_m( s_{k })\right) + \frac{1}{2^{2m}} \\
        &\leq \int_{s_{k + 1}}^{s_{k}}W_0(\zeta_m^{k})|(\zeta_m^{k})'|\, \mathrm{d}s + \frac{1}{2^{2(m - k)}}.
    \end{align*}
    Now, observe that we can get the following estimate by using Lemma \ref{arc-length-estimate} and the doubling condition on $f$:
    \begin{align*}
        \int_{s_{k + 1}}^{s_{k}} W_0(\zeta^{k}_m)|(\zeta^{k}_m)'|\, \mathrm{d}s & \leq c_2\int_{s_{k + 1}}^{s_{k}}f(|\zeta^k_m-a|)|( \zeta^{k}_m)'|\, \mathrm{d}s \\
        & \leq C f\left(\frac{\delta}{2^{m - k - 1}}\right)\int_{s_{k + 1}}^{s_{k}}|( \zeta_m^{k})'|\, \mathrm{d}s \\
        &\leq Cf\left(\frac{\delta}{2^{m-k}} \right) \frac{1}{2^{m - k}}.
    \end{align*}
    Putting these inequalities together, we derive
    \[
        \int_{s_{k+1}}^{s_k}|(\tilde \varphi_m)'|\, \mathrm{d}s \leq \left(\frac{C}{2^{m - k}} \right) .
    \]
    Summing over $k$ we infer
    \begin{align}\label{eq:geodesic_estimate_1}
        \int_{s_{k_0}}^{s_{0}} |(\tilde \varphi_m)'| \, \mathrm{d}s &\leq C.
    \end{align}

    \textbf{Step 3: Estimating the path length of $\varphi_m|_{[-1, s_{k_0}]}$.} 
    If $k_0 < m$ we are done since then $s_{k_0} = -1$ and we have nothing to estimate. In the other case, we observe that \eqref{eq:uniform_estimate_sequence} implies $W_0(\varphi_m|_{(-1, s_m)}) \geq \tilde \alpha$ where 
    $$\tilde \alpha := \min_{B_K(0)\setminus(B_\delta(a)\cup B_\delta(b))}W_0(u)$$
    and $K > 0$ is from $\eqref{eq:uniform_estimate_sequence}$.
    Note that if the stronger assumption \ref{G4} holds, we can choose 
    \[
        \tilde \alpha = \min\left\{ \alpha, \min_{B_R(0)\setminus(B_\delta(a)\cup B_\delta(b))}W_0(u) \right\}
    \]
    where $\alpha$ and $R$ are given by \ref{G4}.
    Using this, we estimate with \ref{G1}: 
    \begin{align}\label{eq:geodesic_estimate_2}
        \tilde \alpha \int_{-1}^{s_{m}} |(\tilde \varphi_m)'| \, \mathrm{d}s &\leq \int_{-1}^{s_m} W_0(\tilde \varphi_m)|(\tilde \varphi_m)'|\, \mathrm{d}s \nonumber \\
        &\leq \dd_{W_0}\left(\tilde \varphi_m(-1), \tilde \varphi_m( s_{m} )\right) + \frac{1}{2^{2m}} \nonumber \\
        &\leq \left(\sup_{u \in B_{1 + |p|}(a)} W_0(u)\right) |p - \tilde \varphi_m(s_m)| + \frac{1}{2^{2m}} \nonumber \\
        &\leq C \left( f(1 + |p|)|p - a| \right) + \frac{1}{2^{2m}}
    \end{align}
    Notice now that the doubling condition of $f$ (from \ref{G2}) implies $f(1 + |p|) \leq C(1 + f(|p|))$ so we can use \eqref{eq:geodesic_estimate_1} and \eqref{eq:geodesic_estimate_2} to obtain
    \[
        \int_{-1}^{t_0^1} |(\varphi_{m})'| \, \mathrm{d}s \leq C(1 + f(|p|)|p|).
    \]
    In an analogous way, one estimates the path lengths in the remaining intervals $[t_1^1, t_0^2]$ and $[t_0^2, t_1^2]$ to derive 
    \[
        \int_{t_1^1}^{t_0^2} |(\varphi_{m})'| \, \mathrm{d}s \leq C(1 + f(|q|)|q| + f(|p|)|p|).
    \]
    and 
    \[
        \int_{t_0^2}^1 |(\varphi_{m})'| \, \mathrm{d}s \leq C(1 + f(|q|)|q|).
    \]

    \textbf{Step 4: Existence of a solution for $\dd_{W_0}(p,q)$.} 
    This follows by a standard argument based on the Ascoli-Arzel\`{a} Theorem. For instance, see \cite[Lemma 3.1]{CriGra} and note that the continuity of the potential can be weaken to lower semi-continuity.
    % Thanks to \eqref{eq:geodesic_bounded_Euclidean}, this follows directly by a standard argument using the Ascoli-Arzel\`{a} Theorem.
\end{proof}

\begin{remark}\label{remark:zuniga_sternberg}
Note that we are able to prove the existence of a geodesics for $\dd_{W_0}(p,q)$ even if the potential $W_0$ is not continuous, but lower semi-continuous and satisfies certain growth conditions.
Our result is in the same spirit as that obtained by Zuniga and Sternberg in \cite{ZunSter}, where they require the potential $W_0$ to be continuous, but without assuming any growth condition near the wells.
In their case, though, they cannot ensure that geodesics have finite Euclidean length. 
Indeed, as Example \ref{ex:geodesics_infinite_Euclidean_length} shows, there exist potentials satisfying the assumptions of the main result by Zuniga and Sternberg whose corresponding geodesics have infinite Euclidean length.
\end{remark}


\begin{figure}
    \includegraphics[width=0.5\linewidth]{maze}
    \caption{The idea for the construction of potentials $W$ in Example \ref{ex:geodesics_infinite_Euclidean_length}. Heuristically, we create obstacles in the energy landscape such that geodesics connecting points in $C_1$ to points in $C_{1/2}$ cannot pass through.}
\end{figure}

\begin{example}\label{ex:geodesics_infinite_Euclidean_length}
We sketch a construction of a potential for which geodesics with respect to the induced distance exist, but necessarily have unbounded Euclidean path length. \\
Let 
$$ \frac{1}{2} < r_4 < r_3  < \frac{3}{4} < r_2 < r_1 < 1.$$
For $\delta < \pi/16$ consider the annular sectors
\[
    A_1 \coloneqq \left\{ x \in \R^2 : r_2 < |x| < r_1 \, -\frac{\pi}{2} - \delta < \arg(x) < -\frac{\pi}{2} + \delta \right\},
\]
and
\[
    A_2 \coloneqq \left\{ x \in \R^2 : r_4 < |x| < r_3, \, \frac{\pi}{2} - \delta < \arg(x) < \frac{\pi}{2} + \delta \right\}
\]
where $\arg: \R^2 \to [-\pi, \pi)$ denotes the polar angle.
Let 
\[
    C_r := \{ x \in \R^2 : |x| = r \}
\]
be the circle lines with radius $r > 0$, and 
\[
    S_1 = \{0\}\times(-1, -3/4) \qquad \text{ and } \qquad S_2 = \{0\} \times (1/2, 3/4).
\]
For $M_1 = 1$ and $M_1/2 > \eps_1 > \eps_2$ define
\[
    \tilde W(p) := \begin{cases}
        \eps_1 & p \in C_1 \cup S_1, \\
        M_1 & p \in A_1, \\
        \eps_1 & p \in C_{3/4}, \\
        M_1 & p \in A_2, \\
        \eps_2 & p \in C_{1/2}, \\
        (1-t) \eps_2 + t \eps_1 & p = \left(0, \frac{t+1/2}{2}\right) \in S_2
    \end{cases}
\]
and let $W$ be the natural piecewise linear extension to $\overline{B_1/B_{1/2}}$. If $\eps_1$ is chosen small enough, any geodesic $\varphi$ with respect to $\dd_W$ connecting two points $p_1 \in C_1$ and $p_2 \in C_{1/2}$ will not pass through $A_1, A_2$. Indeed, suppose that $\varphi(s) \in A_1 \cup A_2$ and $\varphi(t) \in C_{3/4}$ for some $s,t \in [-1,1]$. Then,
\[
    \left| \int_{s}^t W_0(\varphi)|\varphi'| \, \dd s \right| \geq C \int_{\eps}^{M_1} r \dd r \geq CM_1^2.
\]
But taking an injective competitor $\tilde \varphi$ that is contained in $C_1 \cup S_1 \cup C_{3/4} \cup S_2 \cup C_{1/2}$ we derive 
\[
    \dd_{W}(p_1, p_2) \leq C\eps_1
\]
which yields a contradiction for $\eps_1 $ chosen small enough since $\varphi$ is a geodesic with respect to $\dd_{W}(p_1, p_2)$. This has the implication that the geodesic does not pass through $A_1$ and $A_2$, and we have that the path length of $\varphi$ is bounded from below by $1$.  \\
To extent $W$ to $\overline{B_{1/n}\setminus B_{1/(n+1)}}$, we repeat this construction with suitable $M_n > 0$ and $\eps_n > 0$ such that
$$  
    M_n \xrightarrow{n \to \infty} 0,
$$ 
and that any geodesic $\psi$ with respect to $d_W$ connecting two points in $C_{1/n}$ and $C_{1/(n+1)}$ fulfils
\begin{align}\label{eq:piecewise_lower_bound_pathlength}
    \int_{-1}^1 |\psi'| \, \dd r \geq \frac{C}{n}.
\end{align}
By extending $W$ with $M_1$ outside of the closed unit ball, we get a continuous potential defined on $\R^M$. In particular, geodesics between two points exist (cf.\ \cite{ZunSter}). 
Next, let $\varphi$ be a geodesic connecting $p \in C_1$ with $0$. Since $\varphi([-1,1]) \cap C_{1/n} \neq \emptyset$, we can take an increasing sequence $s_n \in \varphi^{-1}(C_{1/n})$. Now, we can compute a lower bound on the path length by 
\begin{align*}
    \int_{-1}^1 |\varphi'| \dd r \geq \sum_{n \in \N} \int_{s_n}^{s_{n+1}} |\varphi'| \, \dd r \geq  
    \sum_{n \in \N} \frac 1 n = +\infty.
\end{align*}
We have shown that any geodesic with respect to this constructed potential $W$ always has unbounded path length.
\end{example}
In our setting, we will encounter potentials that jump. For this reason, we introduce an adapted version of a geodesic distance function.
\begin{definition}
    Let $W_0, W_1: \R^M \to [0, \infty)$ be non-negative Borel functions. We define the adapted geodesic distance by
    \[
        \dd_W(p,q) := \inf_{r \in \R^M} (\dd_{W_0}(p, r) + \dd_{W_1}(r, q)).
    \]
\end{definition}

An analogous result as in Lemma \ref{lemma:geodesic_distance_is_locally_lipschitz} holds also for the adapted distance. More precisely, we have the following statement.
\begin{lemma}\label{lemma:adapted_geodesic_distance_is_locally_lipschitz}
    Suppose that $W_0, W_1: \R^M \to [0,\infty)$ are locally bounded Borel functions. Then, $\dd_W$ is locally Lipschitz, and the local Lipschitz constants of $\dd_W$ are lower than or equal to the Lipschitz constants of $\dd_{W_0}$ and $\dd_{W_1}$.
\end{lemma}

If \ref{G3} holds then we also can prove an analogous result as in Lemma \ref{lemma:eps_minimizers_are_bounded}, whose proof works in exactly the same way. 
\begin{proposition}\label{prop:infimum_attained_in_ball}
    Suppose that $W_0, W_1: \R^M \to [0,\infty)$ are locally bounded Borel functions that fulfill \ref{G3}, and let $r > 0$. Then, there exists $K_0 > r$ only depending on $r$,  $|a|$, $|b|$, $f$ and $C_G$ (given in \ref{G3}) such that 
    \[
        \dd_W(p,q) = \inf_{s \in B_{K_0}(0)} (\dd_{W_0}(p, s) + \dd_{W_1}(s, q)),
    \]
    for all $p, q \in B_R(0)$. 
\end{proposition}


The next corollary is a direct application of Proposition \ref{prop:infimum_attained_in_ball} and Lemma \ref{lemma:eps_minimizers_are_bounded}.
\begin{corollary}\label{corol:truncation_constant}
    Suppose that $W_0, W_1$ are locally bounded Borel functions fulfiling \ref{G1}, \ref{G3}, and let $r > 0$. Then, there exists $M_0 = M_0(r, a, b, f, C_G) > 0$ such that for all $M > M_0$ we have that 
    \[
        \dd_{W \land M}(p, q) = \dd_{W}(p, q),
    \]
    for all $p, q \in B_R(0)$, where 
    \[
        \dd_{W \land M}(p, q) := \inf_{r \in \R^N} ( \dd_{W_0 \land M}(p, r) + \dd_{W_1 \land M}(r, q) ).
    \]
\end{corollary}
\begin{proof}
    We have
    \[
        \dd_{W \land M}(p, q) \leq \dd_{W}(p, q),
    \]
    for all $p,q \in B_R (0)$.
    Now, let $K_0 > 0$ be the constant given in Proposition \ref{prop:infimum_attained_in_ball} and $K_1$ be the constant of Lemma \ref{lemma:eps_minimizers_are_bounded} with respect to the parameters $K_0$, $|a|$, $|b|$, $f$ and $C_G$. Set 
    $$ M_0 \coloneqq \max \left\{\|W_0 \|_{L^\infty(B_{K_1}(0))}, \|W_1 \|_{L^\infty(B_{K_1}(0))} \right\}.$$ 
    Let $M > M_0$. We observe first that Proposition \ref{prop:infimum_attained_in_ball} coupled with Remark \ref{remark:geodesic_parameter_discussion_2} applied to $\tilde f \coloneqq f \land M$ gives $r_0 \in B_{K_0}(0)$ such that  
    \[
         \dd_W(p,q) = \dd_{W_0\land M}(p, r_0) + \dd_{W_1 \land M}(r_0, q).
    \]    
    By Lemma \ref{lemma:eps_minimizers_are_bounded} together with Remark \ref{remark:geodesic_parameter_discussion_2} and applied to $\tilde f \coloneqq f\land M$, for any $\eps \in (0,1)$ we can find curves $\varphi_0, \, \varphi_1 \in W^{1,1}([-1,1], \R^M)$ which are $\eps$-minimizers of $\dd_{W_0 \land M}(p, r_0)$ (resp. $\dd_{W_1 \land M}(r_0, q)$) such that they are bounded in $L^\infty$ by $K_1$.  In particular, we infer 
    \begin{align*}
        \dd_W(p,q) &\leq \int_{-1}^1 W_0(\varphi_0)|\varphi_0'| \, ds + \int_{-1}^1 W_1(\varphi_1)|\varphi_1'| \, ds  \\
        &=  \int_{-1}^1 (W_0\land M)(\varphi_0)|\varphi_0'| \, ds + \int_{-1}^1 (W_1\land M)(\varphi_1)|\varphi_1'| \, ds \\
        &= \dd_{W_0}(p, r_0) + \dd_{W_1}(r_0, q) \\
        &= \dd_{W \land M}(p, q),
    \end{align*}
    which concludes the proof.
\end{proof}


%%%%%%%%%%%%%%%%%%%%%%%%%%%%%%%%%%%%%%%%%%%%
%%%%%%%%%%%%%%%%%%%%%%%%%%%%%%%%%%%%%%%%%%%%
%%%%%%%%%%%%%%%%%%%%%%%%%%%%%%%%%%%%%%%%%%%%
%%%%%%%%%%%%%%%%%%%%%%%%%%%%%%%%%%%%%%%%%%%%


\subsection{An approximation result for sets of finite perimeter}

In this section, we present an approximation result for sets of finite perimeter with piecewise $C^2$ sets, both in configuration and in energy.
The peculiarity of such an approximation is that the energy is lower semi-continuous, and jumps on a piecewise $C^2$ surface.
Therefore, the smooth approximating sequence needs to preserve the parts of the original set on such boundaries, as much as possible.
A similar result was obtained in \cite[Lemma 5.25]{CriFonGan_moving_sub} in the case the jumps of the energy where on sets with polyhedral boundary.
Here, we reproduce a similar proof by adding more details.

\begin{proposition}\label{prop:approx_sets}
Let $\Omega_1\dots,\Omega_k$ be a partition of $\Omega$, where each $\partial\Omega_i$ is piecewise $C^2$.
Let $g:\Omega\to\R$ be a bounded lower semi-continuous function such that its restriction to $\Omega\setminus\cup_{i=1}^k \partial\Omega_i$ is continuous.
Let $E\subset\Omega$ be a set with finite perimeter in $\Omega$.
Then, there exists a sequence of sets $\{F_n\}_n$ with piecewise $C^2$ boundary such that
\[
\lim_{n\to\infty} |E\triangle F_n|=0,\quad\quad\quad
\lim_{n\to\infty} \left| \int_{\partial^*E} g(x) \dhno(x) 
    - \int_{\partial^*F_n} g(x) \dhno(x) \right| =0.
\]
Moreover, for each $i\in\{1,\dots,k\}$, it holds that $\partial F_n\cap \partial\Omega_i$ is relatively open in $\partial\Omega_i$.
\end{proposition}

\begin{proof}
Fix $\lambda>0$. We show that there exists a set $F\subset \Omega$ with piecewise $C^2$ boundary such that
\[
|E\triangle F|\leq\lambda,\quad\quad\quad
\left| \int_{\partial^*E} g(x) \dhno(x)
    - \int_{\partial^*F} g(x) \dhno(x) \right| \leq \lambda,
\]
and such that for each $i\in\{1,\dots,k\}$ it holds that $\partial F\cap \partial\Omega_i$ is relatively open in $\partial\Omega_i$.\\

\textbf{Step 1: Piecewise $C^1$ approximation.}
Fix $\mu>0$.
We claim that there exists a set $L\subset\Omega$ with boundary contained in a finite union of $C^1$ hypersurfaces, such that
\[
|E\triangle L|\leq\mu,\quad\quad\quad
\left| \int_{\partial^*E} g(x) \dhno(x)
    - \int_{\partial^*L} g(x) \dhno(x) \right| \leq \mu.
\]
Indeed, fixed $\varepsilon>0$, let $L\subset\Omega$ be the set given by the strong approximation result sets of finite perimeter by de Gromard (see Theorem \ref{thm:strong_BV}).
Then, it holds that
\[
|E\triangle L|\leq\varepsilon.
\]
Moreover, using the other properties of $L$ that Theorem \ref{thm:strong_BV} ensured, we get that
\begin{align*}
\Bigl| \int_{\partial^*E} g(x) \dhno(x)
    &- \int_{\partial^* L} g(x) \dhno(x) \Bigr| \\
&\leq
    \int_{\partial^*E\setminus \partial^* L} g(x)\, \dhno(x) 
    +\int_{\partial^* L \setminus \partial^*E} g(x)\, \dhno(x) \\
&\leq C_1\left[\, \hno\left(\partial^*E\setminus \partial^* L \right)
    + \hno\left( Q\cap \partial^*L \setminus \partial^*E \right) \,\right] \\
&\leq C_1\left[\, |D\ca_E|(\Omega\setminus C )
    + \hno\left(\partial^* L \setminus \partial^*E \right) \,\right] \\
&\leq 2 C_1\varepsilon,
\end{align*}
where $C\subset\Omega$ is the set provided by Theorem \ref{thm:strong_BV} , and we used the fact that the function $g$ is bounded.
Therefore, choosing $\varepsilon>0$ small enough, we obtain the desired claim.\\

\textbf{Step 2: Isolation of singularities.}
Fix $\mu>0$.
We claim that there exists a set $S\subset\Omega$ such that
\begin{itemize}
\item[(i)] $S$ is piecewise $C^2$;
\item[(ii)] $\partial L\setminus S$ is a pairwise disjoint union of finitely many $C^1$ surfaces compactly contained in $\Omega$;
\item[(iii)] $\partial S$ is orthogonal to $\partial L$, in the sense that $\nu_{\partial S}(x)\cdot\nu_{\partial L}(x)=0$ for all $x\in\partial L \cap \partial S$;
\item[(iv)] $\hno(\partial L\cap \overline{S})\leq \mu$.
\end{itemize}
Indeed, let $\Sigma$ denote the singular set of $\partial L$ union the set $\partial L \cap\partial\Omega$.
Then, $\Sigma$ is compact, being the intersection of finitely many compact $C^1$-surfaces.
Thus, for any given $r>0$, it is possible to cover $\Sigma$ with the finite union of balls of radius $r$.
We denote such a finite union by $S$.
Then, requirements (i) and (ii) above hold directly, while (iv) is in force provided $r>0$ is chosen small enough.
Finally, a local modification of $\partial S$ ensures that also the requirement in (iii) is satisfied.\\

\textbf{Step 3: Approximation of the contact set.}
Fix $\mu>0$.
We claim that there exists a set $C\subset \overline{L}\setminus S$ with the following properties:
\begin{itemize}
\item[(1)] $C$ is piecewise $C^2$;
\item[(2)] It holds that
\[
\hno\left( \left[\partial L \cap \bigcup_{i=1}^k \partial\Omega_i \right] \triangle \left[ \partial L \cap C \right] \right) \leq\mu;
\]
\item[(3)] $C\cap \partial \Omega_i$ is relatively open in each $\partial \Omega_i$.
\end{itemize}
In order to prove our claim, we reason as follows.
Fix $i\in\{1,\dots,k\}$.
We first notice that the sets
\[
C_i^+\coloneqq \{ x\in \partial L \cap \partial\Omega_i\setminus S : \nu_L(x)=\nu_{\partial\Omega_i}(x) \},\quad
C_i^-\coloneqq \{ x\in \partial L \cap \partial\Omega_i\setminus S : \nu_L(x)=-\nu_{\partial\Omega_i}(x) \}
\]
are compact, because $\partial L$ and $\partial\Omega_i$ are, and $\nu_L$ and $\nu_{\partial\Omega_i}$ are continuous outside of $S$.
Let 
\[
d \coloneqq \min_{i=1,\dots,k} \mathrm{dist}(C^+_i, C^-_i).
\]
Note that $d>0$.
Using this fact, together with the outer regularity of the Radon measure $\hno\restr\partial\Omega_i$, it is possible to find disjoint sets $P_i, N_i\subset\partial\Omega_i$ relatively open in $\partial\Omega_i$ such that
\begin{equation}\label{eq:C_close}
\hno(C^+_i \triangle P_i) + \hno(C^-_i \triangle N_i) < \mu.
\end{equation}
Without loss of generality, using an approximation argument, we can also ensure that $\partial P_i$ and $\partial N_i$ are sets with $C^1$ boundary with
\begin{equation}\label{eq:boundary_Pi_Ci}
\mathcal{H}^{N-2}(\partial P_i)<\infty, \quad\quad\quad
\mathcal{H}^{N-2}(\partial N_i)<\infty.
\end{equation}
We define
\[
C\coloneqq \bigcup_{i=1}^k P_i\cup N_i.
\]
By its very definition and by using \eqref{eq:C_close}, we see that the set $C$ satisfies the required properties.\\

\textbf{Step 4: The approximating set.}
Fix $\mu>0$.
First of all, we note that a standard result for the trace of functions of bounded variation (see \cite[Theorem 2.11]{Giusti}) yields the existence of infinitesimal sequences $\{s^i_n\}_n$ and $\{t^i_n\}_n$ with the following properties:
\begin{equation}\label{eq:hno_Pi}
\left|\, \lim_{n\to\infty} \hno\left( \{ x - s^i_n\nu_L(x) : x\in P_i \} \cap L \right) -  \hno(P_i) \,\right| \leq\mu,
\end{equation}
\begin{equation}\label{eq:hno_Ni}
\left|\, \lim_{n\to\infty} \hno\left( \{ x + s^i_n\nu_L(x) : x\in N_i \} \cap L \right) - \hno(N_i) \,\right| \leq\mu.
\end{equation}
Using the fact that, for each $i=1,\dots, k$, the sets $\partial\Omega_i\setminus S$ is made of finitely many closed surfaces of class $C^2$, it is possible to find $\tau>0$ such that their $\tau$-normal tubular neighborhoods are pairwise disjoint and compactly contained in $\Omega$.
Without loss of generality, we can assume that
\begin{equation}\label{eq:tau}
s_n^i, t_n^i < \tau,
\end{equation}
for all $i=1,\dots, k$, and all $n\in\N$.
We fix an index $n\in\N$ that will be chosen later.
Let $\delta>0$ be such that
\[
\delta \leq \frac{1}{3} \min\left\{ d,  \min\{ s^i_n, t^i_n : i=1,\dots,k \} \right\}.
\]
Let $L_\delta$ be a smooth approximation of $L$ such that its Hausdorff distance from $L$ is less than $\delta$.
Without loss of generality, we can also assume that
\[
\hno\left(\partial L_\delta \cap \bigcup_{i=1}^k \partial\Omega_i\right) = 0.
\]
Define the set
\[
F \coloneqq L_\delta\cap\Omega \, \cup \,
    \bigcup_{i=1}^k P_i(s^i_n)
    \cup \,
    \bigcup_{i=1}^k N_i(t^i_n)
    \cup S,
\]
where, for each $i\in\{1,\dots,k\}$, we set
\[
P_i(s^i_n)\coloneqq \left\{ x - s\nu_L(x) : x\in P_i, s\in[0, s^i_n]  \right\},\quad\quad
N_i(t^i_n)\coloneqq \left\{ x + s\nu_L(x) : x\in N_i, s\in[0, t^i_n]  \right\}
\]
Note that, by construction, the set $F$ is piecewise of class $C^2$.
Moreover, it holds that
\begin{equation}\label{eq:boundary_F}
\hno\left( (\partial F\cap \partial\Omega_i) \triangle (P_i\cup N_i) \right)=0.
\end{equation}
for all $i=1,\dots,k$.\\

\textbf{Step 5: Estimates.}
First of all, we notice that, by taking $\tau$ and $\delta$ sufficiently small, we can ensure that
\[
|E\triangle F|\leq \mu.
\]
In order to prove the energy estimate, we reason as follows.
We write
\[
\partial^* F = F_r \cup F_s \cup F_a,
\]
as a pairwise disjoint union, where
\[
B_s\coloneqq \partial^* F \cap \overline{S},
\]
\[
B_r\coloneqq \partial^*F\cap \overline{R},\quad\quad\quad
    R\coloneqq \bigcup_{i=1}^k P_i(s^i_n)
    \cup \,
    \bigcup_{i=1}^k N_i(t^i_n)
    \setminus \overline{S},
\]
and
\[
B_o\coloneqq \partial^*F \setminus(L_a\cup L_s).
\]
First of all, using the continuity of the function $g$ in $\Omega\setminus \cup_{i=1}^k \Omega_i$, we get that
\[
\left| \int_{B_o} g(x)\dd\hno(x)
    -  \int_{\partial L\setminus(\overline{S}\cup\overline{R})} g(x)\dd\hno(x) \right| \leq \frac{\lambda}{3},
\]
provided $\delta>0$ is chosen small enough.
Now, using (iv) of Step 2 together with the boundness of the function $g$ and \eqref{eq:boundary_F}, we get that
\[
\left| \int_{B_s} g(x)\dd\hno(x)
    -  \int_{\partial L\cap \overline{S}} g(x)\dd\hno(x) \right| \leq \frac{\lambda}{3},
\]
since it is possible to choose $\delta>0$ small enough so that
\[
\hno(\partial^* F\cap \overline{S})\leq \mu.
\]
Finally, we see that
\begin{align}\label{eq:est_rectangles}
\Bigl| \int_{B_r} g(x)\dd\hno(x)
   & -  \int_{\partial L\cap \overline{R}} g(x)\dd\hno(x) \Bigr|
    \leq
        \sum_{i=1}^k\left| \int_{(B_r\cap\partial\Omega_i) \triangle (\partial L\cap\partial\Omega_i)} g(x)\dd\hno(x)  \right| \nonumber \\
&\hspace{3cm} + C \sum_{i=1}^k \left( s_n^i \mathcal{H}^{N-2}(\partial P_i) \right)
    + C \sum_{i=1}^k \left( t_n^i \mathcal{H}^{N-2}(\partial N_i) \right)
 \nonumber \\
&\hspace{3cm} + C \hno(\partial^*E \cap R) + C\hno(B_r\cap R),
\end{align}
where we used the boundness of the function $g$.
Now, using \eqref{eq:C_close}, \eqref{eq:hno_Pi}, and \eqref{eq:hno_Ni}, we can choose $\mu>0$ small enough so that
\[
\sum_{i=1}^k\left| \int_{(B_r\cap\partial\Omega_i) \triangle (\partial L\cap\partial\Omega_i)} g(x)\dd\hno(x)  \right| \leq \frac{\lambda}{18}.
\]
Moreover, noting that once $\mu>0$ is chosen, then the sets $P_i$ and $N_i$ are also chose, using \eqref{eq:boundary_Pi_Ci} and \eqref{eq:tau}, it is possible to choose $\tau>0$ small enough so that
\[
C \sum_{i=1}^k \left( s_n^i \mathcal{H}^{N-2}(\partial P_i) \right)
    + C \sum_{i=1}^k \left( t_n^i \mathcal{H}^{N-2}(\partial N_i) \right)
    \leq \frac{\lambda}{18}.
\]
Finally, we note that thanks to \eqref{eq:hno_Pi}, and \eqref{eq:hno_Ni}, we have that, up to choosing $\tau>0$ small enough, we can ensure that
\[
C \hno(\partial^*E \cap R) + C \hno(B_r\cap R) \leq\frac{\lambda}{18}.
\]
This proves the desired estimate and concludes the proof.
\end{proof}



%%%%%%%%%%%%%%%%%%%%%%%%%%%%%%%%%%%%%%%%%%%%
%%%%%%%%%%%%%%%%%%%%%%%%%%%%%%%%%%%%%%%%%%%%
%%%%%%%%%%%%%%%%%%%%%%%%%%%%%%%%%%%%%%%%%%%%
%%%%%%%%%%%%%%%%%%%%%%%%%%%%%%%%%%%%%%%%%%%%



\section{Compactness}\label{sec:compactness}

This section is devoted to investigate the compactness of sequences with uniformly bounded energy.

\begin{proof}[Proof of Theorem \ref{thm:compactness}]
    The main idea of this proof is to apply a transformation to the wells such that they become fixed, then obtain the result for this case, and then apply the inverse transformation.
    Recalling Remark \ref{rem:T_symmetric}, we have that $b(x) = - a(x)$, and we choose $a(x)^\perp \in \R^{M \times (M-1)}$ such that 
    \[
        R_a(x) := \frac{1}{|a(x)|}(a(x), a(x)^\perp).
    \]
    is a frame and in $BV(\Omega;\R^{M \times M})$. Note that this is always well-defined thanks to assumption \ref{H3}.\\
    Next, we recall the definition of the adjustment $T_a(x)$,
    \[
        T_a(x) := |a(x)|R_a.
    \]
    Note that, since $a \in BV (\Omega; \R^M) \cap L^\infty (\Omega; \R^M)$ and $x\mapsto |a(x)|$ is bounded away from zero, it holds that
    \begin{equation}\label{eq:gradient_Ta}
        | \nabla T_a | \leq C | \nabla a|, \qquad | \nabla T_a^{-1} | \leq C | \nabla a|.
    \end{equation}
    We start with showing that $u_n$ is uniformly bounded in $L^1$ and equiintegrable. Indeed, using \ref{H6}, we choose a constant $C_1 > C_2$ such that 
    \[
        \int_{|u_n| > C_1} |u_n| \, \mathrm{d}x \leq C_2 \int_\Omega W(x, u_n) \, \mathrm{d}x \leq C_2C\eps_n,
    \]
    which implies equiintegrability. Using again \ref{H6} we also have 
    \begin{align*}
        \int_\Omega |u_n|\, \dd x &= \int_{|u_n|\leq C_1}|u_n| \, \dd x + \int_{|u_n| > C_1}|u_n| \, \dd x \\
        &\leq C_1|\Omega| + C_2 \int_\Omega W(x, u_n) \, \dd x\\
        &\leq C_1 | \Omega| + C_2 C \eps_n,
    \end{align*}
    which implies uniform boundness.\\
    Since $a \in L^\infty(\Omega; \R^M)$, we deduce $T_a, T_a^{-1} \in L^\infty(\Omega;\R^M)$. This, in particular, implies that 
    $$ v_n \coloneqq T_a^{-1} u_n, $$
    is uniformly bounded in $L^1$ and equiintegrable, since $u_n$ is uniformly bounded in $L^1$ and equiintegrable. It therefore generates a suitable Young measure $\nu_x$ (cf.\ Theorem \ref{thm:compactness_YM}), and $v_n \rightharpoonup v_0$ in $L^1(\Omega; \R^M)$.
    
    We now notice that 
    \[
        CW(x,u_n) \geq Cf(|a(x)|\min \{|v_n - e_1|, \, |v_n + e_1| \}) \geq f(\delta \min \{|v_n - e_1|, \, |v_n + e_1| \}) =: \widetilde W(v_n).
    \]
    Since $\widetilde{W}$ is continuous, we infer 
    \[
        \spt \nu_x \subset \{\pm e_1 \},
    \]
    which implies 
    \[
        \nu_x = \theta(x) \delta_{e_1} + (1-\theta(x)) \delta_{-e_1},
    \]
    for a measurable function $\theta: \Omega \to [0,1]$.
    
    Now, we consider the geodesic distance function
    \[
        \dd_{\tilde W}(p,q) \coloneqq \inf \left\{ \int_{-1}^1 2 \min\left\{\sqrt{\widetilde W(\varphi)}, K \right\}|\varphi'| \, \mathrm{d}s : \varphi \in W^{1,1}([-1,1], \R^M), \, \varphi(-1) = p,\, \varphi(1) = q \right\}.
    \]
    for a suitably large constant $K > 0$. Thanks to Lemma \ref{lemma:geodesic_distance_is_locally_lipschitz}, this distance is globally Lipschitz continuous and its Lipschitz constant is bounded by 
    \[
        |\nabla \dd_{\tilde W}| \leq C(1 + 2K).
    \]
    We define, for $L>0$, the component-wise truncation map $\mathcal{T}_L:\R^M\to\R^M$, and set
    \[
    w_n(x) \coloneqq \mathcal{T}_L v_n(x).
    \]
    Note that
    \begin{equation}\label{eq:w_n_restricted_to_Q_L}
        |D w_n| \leq |D v_n|\restr Q_L,
    \end{equation}
    where $Q_L\subset\R^M$ is the cube centered at the origin and with side length $L$, and therefore is a $BV$ function.
    Moreover, by choosing $L$ large enough, $w_n$ and $v_n$ generate the same Young measure.
    
    Let us define 
    $$ z_n(x) \coloneqq \dd_{\tilde W}(-e_1, w_n(x)), $$
    as the geodesic distance from a well in the transformed space.
    Since $w_n \in BV(\Omega; \R^M)$ and the geodesic distance belongs to $\mathrm{Lip}(\R^M \times \R^M; \R)$, by \cite[Proposition 3.69(c)]{AFP} we get that $z_n \in BV(\Omega; \R)$.
    Also, since $w_n$ is given by
    $$ w_n = \mathcal{T}_L T_a^{-1} u_n, $$
    it follows that 
    \begin{equation}\label{eq:estimate_jump_sets_Cantor}
        J_{w_n} \subseteq J_{T_a^{-1}} = J_a, \qquad | D^c w_n | \leq C |D^c a|.
    \end{equation}
    We now want to estimate $| D z_n | (\Omega)$. Therefore, using the chain rule in $BV$ (see \cite[Theorem 3.101]{AFP}), together with \eqref{eq:gradient_Ta} and \eqref{eq:estimate_jump_sets_Cantor} we get
    \begin{align*}
        |D z_n|(\Omega) &\leq \int_\Omega |\nabla \dd_{\tilde W}(-e_1, w_n)||\nabla w_n| \, \mathrm{d}x + C \H^{d-1}(J_{w_n} \cap \Omega) + C |D^c w_n|(\Omega) \\
        &\leq C \left( \int_\Omega 2\sqrt{\widetilde{W} (w_n)}|\nabla w_n| \, \mathrm{d}x + \H^{N-1}(J_a \cap \Omega) + |D^c a|(\Omega) \right).
    \end{align*}
    Since $\widetilde{W}(v_n)$ is bounded above by $C W(x,u_n)$, and $w_n$ is the component-wise truncation of $v_n$, using \eqref{eq:w_n_restricted_to_Q_L} we get
    \begin{align*}
        \int_\Omega 2 \sqrt{\widetilde{W}(w_n)} |\nabla w_n| \, \dd x &\leq \int_\Omega 2 \sqrt{W(x,u_n)} |T_a^{-1}| | \nabla u_n| \, \mathrm{d} x + \int_{Q_L} 2 \sqrt{\widetilde{W}(w_n)} | u_n | | \nabla T_a^{-1}| \, \mathrm{d} x \\
        &\leq \frac{1}{\delta} \mathcal{F}_{\eps_n} (u_n) + C_L \| a \|_\infty | D a|(\Omega).
    \end{align*}
    In particular, this implies that for sequences $\{ u_n \}_n$ uniformly bounded in energy, the total variation of $z_n$ is uniformly bounded. Then, up to a subsequence (without relabelling), we have $z_n \overset{*}{\rightharpoonup} z$ in $BV(\Omega; \R)$. By the Rellich-Kondrachov $BV$ compactness Theorem we know that, up to a subsequence, the convergence is also strongly in $L^1$.\\
    
    Since $w_n$ is equiintegrable and $\dd_{\tilde W}(-e_1, \cdot)$ is a Lipschitz function, we infer that $z_n$ is also equiintegrable. Therefore, it generates a Young measure $\eta_x$ and, by the strong convergence in $L^1$, we derive $\eta_x = \delta_{z(x)}$. 
    
    Now, notice that we can test with $\varphi \in C_0(\Omega)$ and $\psi \in C_0(\R^N)$ to derive 
    \begin{align*}
        \int_{\Omega} \int_{\R} \varphi(x)\psi(y) \, \dd\eta_x(y) \, \mathrm{d}x &= \int_{\Omega} \varphi(x) \psi(z(x)) \, \mathrm{d}x \\
        &= \lim_{n \to \infty}\int_\Omega \varphi(x)\psi(z_n(x)) \, \mathrm{d}x \\
        &= \int_\Omega \int_{\R^N} \varphi(x) \psi(\dd_{\tilde W}(-e_1, p)) \dd \nu_x(p) \dd x \\
        &= \int_\Omega \varphi(x) (\theta(x) \psi(\dd_{\tilde W}(-e_1, e_1)) + (1-\theta(x))\psi(0)) \, \dd x.
    \end{align*}
    From the arbitrariness of $\varphi$ we infer 
    \begin{align*}
        \langle \eta_x, \psi \rangle = \int_\Omega \psi(z(x)) \, \mathrm{d}x &= \int_\Omega \theta(x) \psi(\dd_{\tilde W}(-e_1, e_1)) + (1-\theta(x))\psi(0) \, \dd x \\
        &= \left\langle \theta(x)\delta_{\dd_{\tilde W}(-e_1, e_1)} + (1 - \theta(x)\delta_0), \psi \right\rangle,
    \end{align*}
    which holds pointwise for almost every $x \in \Omega$. This in turn implies, by duality, that
    \[
        \delta_{z(x)} = \eta_x = \theta(x)\delta_{\dd_{\tilde W}(-e_1, e_1)} + (1 - \theta(x)\delta_0),
    \]
    for almost every $x \in \Omega$. 
    We infer $\theta(x) \in \{0, 1\}$. Since $z \in BV(\Omega; \R)$ we derive that $\theta(x) = \chi_E$ for a set of finite perimeter $E$. 
    
    Now, taking an arbitrary $\varphi \in C(\Omega)$, we can derive 
    \begin{align*}
        \int_\Omega w(x) \varphi(x) \, \mathrm{d}x&= \lim_{n \to \infty} \int_\Omega w_n(x) \varphi(x) \, \mathrm{d}x  \\
        &= \int_\Omega \varphi(x) \int_{\R^N} p \, \dd \nu_x(p) \, \dd x \\
        &= \int_\Omega \varphi(x) (\chi_E(x)e_1 + (1 - \chi_E(x) (-e_1)) \, \dd x.
    \end{align*}
    Consequently, 
    \[
        w(x)= \chi_E(x)e_1 + (1 - \chi_E(x)) (-e_1).
    \]
    Since the same argument also applies for $v$, we infer $v = w$.
    Furthermore, we have for some $q \in (1, \infty)$ that
    \[
        \lim_{n \to \infty} \int_\Omega |w_n|^q \, \mathrm{d}x = \int_{\Omega} \int_{\R^N}|p|^q \, \dd\nu_x(p) \, \mathrm{d}x = |\Omega|.
    \]
    Since $w_n$ is uniformly bounded in $L^\infty$, it also converges weakly to some $w \in L^q$. Thanks to the reflexivity of $L^q$, since we have convergence of the norms and weak convergence of $w_n$ to $w$, we know that $w_n$ converges strongly to $w$ in $L^q$, and therefore in $L^1$. 
    
    Since $v_n$ is equiintegrable and $w_n$ is its truncation at a suitable high value, we know that $v_n$ also converges strongly. Now, since $T_a \in L^\infty(\Omega;\R^{M \times M})$ we can finally derive that
    \[
        \lim_{n \to \infty} u_n = \lim_{n \to \infty} T_a v_n = T_a v =  \chi_E(x)a(x) + (1 - \chi_E(x)) (-a(x)),
    \]
    in $L^1$, and we conclude.
\end{proof}



%%%%%%%%%%%%%%%%%%%%%%%%%%%%%%%%%%%%%%%%%%%%
%%%%%%%%%%%%%%%%%%%%%%%%%%%%%%%%%%%%%%%%%%%%
%%%%%%%%%%%%%%%%%%%%%%%%%%%%%%%%%%%%%%%%%%%%
%%%%%%%%%%%%%%%%%%%%%%%%%%%%%%%%%%%%%%%%%%%%



\section{Liminf inequality}\label{sec:liminf}

The aim of this section is to prove the $\liminf$ inequality of our $\Gamma$-limit result (see Theorem \ref{thm:Gamma_conv}).
We are able to prove it under a weaker assumption on the regularity of the partition of $\Omega$, compared to \ref{H6}.
    
\begin{theorem}\label{thm:liminf}
    Assume that \ref{H1}-\ref{H5}, and that \ref{H7}-\ref{H8} hold.
    Moreover, assume that there exists a Caccioppoli partition $\Omega_1,\dots,\Omega_k$ of $\Omega$, such that
            \[
            W(x,u) = \sum_{i=1}^k W_i(x,u)\ca_{\Omega_i}(x),
            \]
            \[
            a(x) = \sum_{i=1}^k a_i(x)\ca_{\Omega_i}(x),\quad\quad\quad\quad
            b(x) = \sum_{i=1}^k b_i(x)\ca_{\Omega_i}(x),
            \]
            where $W_i\in C^0(\overline{\Omega}\times\R^M)$ and $a_i, b_i\in W^{1,2}(\Omega;\R^M) \cap C^0(\overline{\Omega};\R^M)$.

    Let $u \in L^1(\Omega;\R^M)$ and $\{u_n\}_n \subset H^1(\Omega;\R^M)$ such that $u_n \to u$ in $L^1(\Omega;\R^M)$. Then, we have
    \[
        \mathcal{F}(u) \leq \liminf_{n \to \infty} \mathcal{F}_{\eps_n}(u_n),
    \]
    for every sequence $\eps_n \to 0$.
\end{theorem}

\begin{proof}
We divide the proof in steps.\\[0.5em]
    \textbf{Step 0: The setup.}
    Let $\delta > 0$ be arbitrary. By passing to a subsequence (not relabeled), we may assume 
    \begin{align}\label{eq:bounded_radon}
        \lim_{n\to \infty} \mathcal{F}_{\eps_n}(u_n) = \liminf_{n \to \infty} \mathcal{F}_{\eps_n}(u_n) < \infty.
    \end{align}
    We apply the blow-up method of Fonseca and M\"uller \cite{fonsecamueller1999}. Define the Radon measures 
    \[
        \mu_n(B) := \mathcal{F}_{\eps_n}(u_n, B)
    \]
    for measurable sets $B \subset \Omega$.
    By \eqref{eq:bounded_radon}, $\{\mu_n\}_n$ is uniformly bounded in $\mathcal{M}_b(\Omega)$, the space of Radon measures in $\Omega$. Thus, there exists $\mu \in \mathcal{M}_b(\Omega)$ such that $\mu_n \rightharpoonup^* \mu$ in $\mathcal{M}_b(\Omega)$ (after extracting a subsequence). Consider now $\lambda := \H^{N-1}|_{J_u}$ and observe that, for $\H^{N-1}$-a.e. $x_0 \in J_u$,
    \[
        \frac{d\mu}{d\lambda}(x_0) < \infty \qquad \text{and} \qquad \frac{d\H^{N-1}}{d\lambda}(x_0) = 1.
    \]
    By the Besicovitch differentiation theorem, we also know that for $\H^{N-1}$-a.e. $x_0 \in J_u$ we have 
    \[
        \frac{d\mu}{d\lambda}(x_0) = \lim_{\rho \to 0^+} \frac{\mu(Q(x_0, \rho))}{\lambda(Q(x_0, \rho))} = \lim_{\rho \to 0^+} \frac{\mu(Q(x_0, \rho))}{\rho^{N-1}},
    \]
    where $Q(x_0, \rho)$ are cubes with side length $\rho > 0$ with two faces orthogonal to $\nu(x_0)$. Since the set of $\rho$ such that $\mu(\partial Q(x_0, \rho)) > 0$ is at most countable, it follows that
    \[
        \frac{\mu(Q(x_0, \rho))}{\rho^{N-1}} = \lim_{n \to \infty} \frac{\mu_n(Q(x_0, \rho))}{\rho^{N-1}},
    \]
    for $\rho \in E$ where $\H^{1}((0,\infty)\setminus E) = 0$. Without loss of generality, we can assume $x_0 = 0$, $\nu(x_0) = e_N$. For $\rho > 0$, we set
    \begin{align*}
        &Q(\rho) := (-\rho/2, \rho/2)^N, \quad Q^+(\rho):= (0, \rho/2)^N,   \\
        &Q^-(\rho) := (-\rho/2, 0)^N, \quad \text{and} \quad Q'(\rho) := (-\rho/2, \rho/2)^{N-1}.
    \end{align*}
    We also set $Q := Q(1)$ and $Q' := Q'(1)$. Now, by applying Lemma \ref{lemma:technical_approx_limits}, we choose a sequence $\rho_m \to 0$ such that 
    \[
        u_n \xrightarrow{n \to \infty} u
    \]
    in $L^1(\partial Q(\rho_m); \R^M)$ holds, and that for some $i,j = 1,.., k$ (with possibly $i = j$) we have
    \begin{align}\label{eq:conv_1}
        u(\rho_m \cdot , \rho_m/2) \xrightarrow{m \to \infty} u^+(x_0),
    \end{align}   
    and
    \begin{align}\label{eq:conv_2}
        u(\rho_m \cdot, -\rho_m/2) \xrightarrow{m \to \infty} u^-(x_0),
    \end{align}
    in $L^1(Q'; \R^M)$.
    Without loss of generality, we can assume that
    \[
    u^+(x_0)=a_i(x_0),\quad\quad\quad
    u^-(x_0)=a_j(x_0),
    \]
    with $i\neq j$.
    The case $u^+(x_0)=a_i(x_0)$ and $u^-(x_0)=b_j(x_0)$, for some $i,j\in\{1,\dots,k\}$, follows by using a similar argument.
    Furthermore, by the slicing theory for $BV$-functions (cf.\ section 3.11 in \cite{AFP}) we infer that for $\H^{N-1}$-a.e. $x' \in Q'$ we have 
    \[
        u_{x'} := u(\rho_m x', \cdot) \in BV((-\rho_m/2,\rho_m/2), \R^M)
    \]
    with the additional property that 
    $$
        J_{u_{x'}} = P_N(J_u \cap (\{\rho_m x'\} \times (-\rho_m/2,\rho_m/2)))
    $$ where $P_N(x) := x_N$. \\[0.5em]
    \textbf{Step 1: Extraction of a nice set}\\[0.5em]
    We will now argue the existence of a large subset of $Q'(1)$ such that $u_{x'}$ is nicely behaved. By the application of the co-area formula \cite[Theorem 2.93]{AFP} with $E = J_u$ and $f(x) = x'$, we derive
    \[
        \frac{1}{\rho_m^{N-1}} \int_{Q'(\rho_m)} \H^0(J_{u_{x'}}) \, \mathrm{d}x' = \frac{1}{\rho_m^{N-1}} \int_{J_u} |\nu \cdot e_n| \, \mathrm{d}x' \leq \frac{\H^{N-1}(J_u \cap Q(\rho_m))}{\rho_m^{N-1}} \xrightarrow{m \to \infty} 1.
    \]
    Moreover, we notice that if $\H^0(J_{u_{x'}}) = 0$ holds for some $x'$ we have $u_{x'} = a_l(\rho_m x', \cdot)$ for some $l=1,..,k$ by \ref{H4}. We obtain
    \[
        \rho_m \leq C\int_{(0, \rho_m)}| u - a_i(x_0)|\, \mathrm{d}s + \int_{(-\rho_m, 0)}| u - a_j(x_0)|\, \mathrm{d}s
    \]
    where $C$ is not dependent on $x'$. Using the definition of approximate limits, we have for large $m$ 
    \begin{align*}
        \rho_m\H^{N-1}(\{x' \in Q'(\rho_m) : \H^0(J_{u_{x'}}) = 0 \}) & \leq C\left( \int_{Q^+(\rho_m)}| u - a_i(x_0)|\, \mathrm{d}x + \int_{Q^-(\rho_m)}| u - a_j(x_0)|\, \mathrm{d}x \right) \\
        &\leq \delta \rho_m^{N}.
    \end{align*}
    In particular, for large $m \in \N$ there exists a measurable set $M_m \subset Q'(\rho_m)$ such that for all $x' \in M_m$ we have 
    \[
        \H^0(J_{u_{x'}}) = 1
    \]
    and 
    \[
        \H^{N-1}(M_m) \geq \left(1 - \frac{\delta}{2} \right)\rho_m^{N-1}.
    \]
    Since the trace of the slices coincides with the standard trace for $\H^{N-1}$-a.e. $x' \in Q'(\rho_m)$, by \eqref{eq:conv_1} and \eqref{eq:conv_2} we can further assume that
    \begin{align*}
        u_{x'}(\rho_m) \to a_i(x_0) \quad \text{ and } \quad  u_{x'}(-\rho_m) \to a_j(x_0)
    \end{align*}
    for all $x' \in M_m$ (after possibly removing a null set contained in $M_m$). Furthermore, we denote the jump point in $J_{u_{x'}}$ by $(x', h(x'))$. We claim that there exists $M_m'\subset M_m$ with 
    $$\H^{N-1}(M_m') \leq \frac{\delta\rho_m^{N-1}}{2}$$ 
    such that for $x' \in M_m\setminus M_m'$ 
    \begin{align}\label{eq:cond_set_1}
        \H^1(\{ (x', t): t > h(x') \} \cap \Omega_l \}) = 0
    \end{align}
    for all $l \neq i$ and 
    \begin{align}\label{eq:cond_set_2}
        \H^1(\{ (x', t): t < h(x') \} \cap \Omega_l \}) = 0
    \end{align}
    for all $l \neq j$. 
    
    Indeed, suppose there exists a subsequence $(m_\alpha)$ for which you cannot extract such a set $N_\alpha \subset M_{m_\alpha}$ with $|N_\alpha|  \geq \delta \rho_{m_\alpha}$. By rescaling, we can infer the existence of sets $\tilde N_\alpha \subset Q'$ with $\H^{N-1}(\tilde N_\alpha) \geq \delta$ such that for all $x' \in \tilde N_\alpha$
    \begin{align}\label{eq:first_case}
        \H^1\left( \{ (x', t): t > h(x') \} \cap \bigcup_{l \neq i} \Omega_l \right) > 0
    \end{align}
    or 
    \[
        \H^1\left( \{ (x', t): t < h(x') \} \cap \bigcup_{l \neq j} \Omega_l \right) > 0.
    \]
    Without loss of generality, let us assume the first case \eqref{eq:first_case} holds for all $x' \in \tilde N_\alpha$. This, in turn, implies that for some constant $C > 0$ and large $m$ we have $|u_{x'}(\rho_m) - a_i(x_0)| > C$ for all $x' \in N_\alpha$ since $u_{x'}$ only has one jump point, $a_l$ are continuous and $a_l(x_0) \neq a_i(x_0)$ for $l \neq i$. This directly contradicts the $L^1$ convergence in \eqref{eq:conv_1} since 
    \begin{align*}
        u_{x'}(\rho_m) = u(\rho_m/2 x', \rho_m/2) 
    \end{align*}
    for $H^{N-1}$-a.e. $x' \in Q'(\rho_m)$ holds. With abuse of notation, we will denote $M_m\setminus M_m'$ by $M_m$ again and note that
    \begin{align}\label{eq:relative_large_good_set}
        \H^{N-1}(M_m) \geq \left(1 - \delta \right)\rho_m^{N-1}
    \end{align}
    holds. \\[0.5em]
    \textbf{Step 2: The blow-up.} After this technical setup, we now perform the blow-up. We remind ourselves of
    \begin{align}\label{eq:blow_up_property}
        \frac{d\mu}{d\lambda}(x_0) = \lim_{m \to \infty} \lim_{n \to \infty} \frac{\mu_n(Q(x_0, \rho_m))}{\rho_m^{N-1}} =  \lim_{m \to \infty} \lim_{n \to \infty} \frac{1}{\rho_m }\int_{Q(\rho_m)} \frac{1}{\eps_n}W(x, u_n) + \eps_n |\nabla u_n|^2 \, \mathrm{d}x.
    \end{align}
    Now, let $M_0 = M_0(\|a\|_\infty, f, C_G)$ be the constant given by Corollary \ref{corol:truncation_constant}. In the following, the truncation of a vector $p \in \R^M$ by a scalar $L > 0$ is to be understood component-wise, i.e., we set $(p \land L)_i := \max\{ \min\{ p_i, L \}, - L \} $. Using \ref{H6}, we infer the existence of $M > M_0$ and $L > \| a\|_\infty$ such that the relation 
    \begin{align*}
        W(p) \land M = W(p \land L) \land M
    \end{align*}
    holds for all $p \in \R^N$.  Now, we define 
    \begin{align*}
        \widetilde W := W \land L, \, \widetilde W_i := W_i \land L\, \text{ and }\tilde u_n := u_n \land M.
    \end{align*}
    We first observe the straightforward estimate
    \begin{align}\label{eq:liminf_partial_inequality_1}
        &\lim_{m \to \infty} \lim_{n \to \infty} \frac{1}{\rho_m^{N-1} }\int_{Q(\rho_m)} \frac{1}{\eps_n}W(x, u_n) + \eps_n |\nabla u_n|^2 \, \mathrm{d}x  \nonumber \\
        & \hspace{1cm} \geq \liminf_{m \to \infty} \liminf_{n \to \infty} \frac{1}{\rho_m^{N-1} }\int_{Q(\rho_m)} \frac{1}{\eps_n}\widetilde W(x, \tilde u_n) + \eps_n |\nabla \tilde u_n|^2 \, \mathrm{d}x.
    \end{align}
    
    Let $T_a$ be defined as in \eqref{eq:adjustment} with inverse $S_a = T_a^{-1}$. Set $v_n := S_a \tilde u_n$. We observe that \ref{H8} implies
    \begin{align}\label{eq:relative_continuity_cond_bounded}
        |\widetilde W_l(x, T_{a_l}(x)v_n) - \widetilde W_l(x_0, T_{a_l}(x_0) v_n)| \leq \omega(\rho_m)W_l(x, T_{a_l}(x)v_n),
    \end{align}
        
    for all $x \in \Omega$ and $l = 1, .., k$.
    Now, we set $$Q_l(\rho_m) := \Omega_l \cap Q(\rho_m).$$ From \eqref{eq:relative_continuity_cond_bounded}, we infer
    \begin{align}\label{eq:liminf_partial_inequality_2}
        &\liminf_{m \to \infty} \liminf_{n \to \infty} \frac{1}{\rho_m^{N-1}} \int_{Q(\rho_m)} \frac{1}{\eps_n}\widetilde W(x, \tilde u_n(x)) + \eps_n |\nabla \tilde u_n|^2 \, \mathrm{d}x  \nonumber \\
        &\hspace{1 cm} = \liminf_{m \to \infty}  \liminf_{n \to \infty} \frac{1}{\rho_m^{N-1}} \sum_{l = 1}^k \int_{Q_l(\rho_m)} \frac{1}{\eps_n}\widetilde W_l(x, T_{a_l}(x)v_n) + \eps_n |\nabla \tilde u_n|^2 \, \mathrm{d}x \nonumber\\
        &\hspace{1 cm} \geq \liminf_{m \to \infty} \liminf_{n \to \infty} \frac{1}{\rho_m^{N-1}}\sum_{l = 1}^k \int_{Q_l(\rho_m)} \frac{1}{\eps_n}\widetilde W_l(x_0, T_{a_l}(x_0) v_n) + \eps_n |\nabla \tilde u_n|^2 \, \nonumber \\ 
        &\hspace{9 cm}- \omega(\rho_m)\frac{1}{\eps_n} \widetilde W_l(x, T_{a_l}(x)v_n)  \, \mathrm{d}x.
    \end{align}
    Now, we notice that we can estimate the last term  
    \begin{align*}
        \frac{1}{\rho_m^{N-1}}\omega(\rho_m)\frac{1}{\eps_n} \int_{Q_l(\rho_m)} \widetilde W_l(x, T_{a_l}(x)v_n)  \, \mathrm{d}x \leq \frac{\omega(\rho_m)\mu_n(Q(\rho_m))}{\rho_m^{N-1}}.
    \end{align*}
    Therefore, by \eqref{eq:blow_up_property} the term vanishes. Next, we infer 
    \begin{align}\label{eq:liminf_partial_inequality_3}
        &\liminf_{m \to \infty} \liminf_{n \to \infty} \frac{1}{\rho_m^{N-1}} \int_{Q(\rho_m)} \frac{1}{\eps_n}\widetilde W(x, \tilde u_n(x)) + \eps_n |\nabla \tilde u_n|^2 \, \mathrm{d}x  \nonumber \\
        & \hspace{1cm} =  \liminf_{m \to \infty} \liminf_{n \to \infty} \frac{1}{\rho_m^{N-1}}\sum_{l = 1}^k \int_{Q_l(\rho_m)} \frac{1}{\eps_n} \widetilde W_l(x_0, T_{a_l}(x_0) v_n) + \eps_n |\nabla \tilde u_n|^2 \, \mathrm{d}x \nonumber \\
        & \hspace{1cm} \geq  \liminf_{m \to \infty} \liminf_{n \to \infty} \frac{1}{\rho_m^{N-1}}\sum_{l \in \{i,j\}} \int_{Q_l(\rho_m)} \frac{1}{\eps_n} \widetilde W_l(x_0, T_{a_l}(x_0) v_n) + \eps_n |\nabla \tilde u_n|^2 \, \mathrm{d}x.
    \end{align}
    Since $\tilde u_n \in L^{\infty}(\Omega; \R^{M})$ by construction, and $S_a \in L^{\infty}(\Omega; \R^{M \times M})$ by \ref{H4}, we can apply the product rule for the approximate gradients of $BV$-functions. We derive $\nabla v_n = \nabla S_a \tilde u_n + S_a \nabla \tilde u_n \in L^2(\Omega; \R^{M \times N})$. Furthermore, using 
    \[
        |T_a Z| = |a||Z|,
    \]
    for $Z \in \R^{M \times N}$ and Young's inequality ($c,d,\delta > 0$)
    \[
        cd \leq \delta c^2 + \frac 1\delta d^2,
    \]
    we obtain
    \begin{align*}
        \eps_n |\nabla \tilde u_n|^2 &= \eps_n |a|^2|\nabla v_n - \nabla S_a \tilde u_n|^2 \\
        &\geq \eps_n(|a|^2(1 - \delta)|\nabla v_n|^2 - C_\delta|\nabla S_a |^2),
    \end{align*} 
    where $C_\delta > 0$ asymptotically behaves like $\frac1\delta$. We apply this with $w_n := T_a(x_0)v_n$ to infer, for $l \in \{i,j\}$, that 
    \begin{align}\label{eq:liminf_partial_inequality_4}
        &  \frac{1}{\rho_m^{N-1}}\int_{Q_l(\rho_m)} \frac{1}{\eps_n}{\widetilde W}_l(x_0, w_n) + \eps_n |\nabla \tilde u_n|^2 \, \mathrm{d}x  \nonumber \\
        & \hspace{1cm}\geq ~  (1-\delta) \frac{1}{\rho_m^{N-1}}\int_{Q_l(\rho_m)} \frac{1}{\eps_n}{{\widetilde W}_l}(x_0, w_n) + \eps_n |a_l(x)|^2|\nabla v_n|^2  - \eps_n C_\delta |S_a|^2 \, \mathrm{d}x.
    \end{align}
    Here, the last term vanishes as $n \to \infty$. We further estimate that
    \begin{align}\label{eq:liminf_partial_inequality_5}
        &\liminf_{m \to \infty} \liminf_{n \to \infty}\frac{1}{\rho_m^{N-1}}\int_{Q_l(\rho_m)} \frac{1}{\eps_n}{{\widetilde W}_l}(x_0, w_n) + \eps_n |a_l(x)|^2|\nabla v_n|^2 \, \mathrm{d}x \nonumber \\
        & \hspace{1cm} = ~ \liminf_{m \to \infty} \liminf_{n \to \infty} \frac{1}{\rho_m^{N-1}}\int_{Q_l(\rho_m)} \frac{1}{\eps_n}{\widetilde W}_l(x_0, w_n ) + \eps_n \left( \frac{|a_l(x)|}{|a_l(x_0)|} \right)^2|\nabla w_n|^2 \, \mathrm{d}x \nonumber\\
        & \hspace{1cm} \geq ~  \liminf_{m \to \infty} \liminf_{n \to \infty} \frac{1}{\rho_m^{N-1}}\int_{Q_l(\rho_m)}\frac{|a_l(x)|}{|a_l(x_0)|} 2\sqrt{{\widetilde W}_l(x_0, w_n)}|\nabla w_n| \, \mathrm{d}x  \nonumber \\
        & \hspace{1cm} \geq ~ \liminf_{m \to \infty} \liminf_{n \to \infty} \frac{1-\tilde \omega_l(\rho_m)}{\rho_m^{N-1}}\int_{Q_l(\rho_m)} 2\sqrt{{\widetilde W}_l(x_0, w_n)}|\nabla w_n| \, \mathrm{d}x \nonumber \\
        & \hspace{1cm} \geq ~ (1-\delta) \liminf_{m \to \infty} \liminf_{n \to \infty} \frac{1}{\rho_m^{N-1}} \int_{Q_l(\rho_m)} 2\sqrt{{\widetilde W}_l(x_0, w_n)}|\nabla w_n| \, \mathrm{d}x, 
    \end{align}
    where $\tilde \omega_l$ is the modulus of continuity of $x \mapsto |a_l(x)/a_l(x_0)|$ at $x_0$. To summarize \eqref{eq:blow_up_property}, \eqref{eq:liminf_partial_inequality_1}, \eqref{eq:liminf_partial_inequality_2}, \eqref{eq:liminf_partial_inequality_3}, \eqref{eq:liminf_partial_inequality_4} and \eqref{eq:liminf_partial_inequality_5}, at this point we have 
    \begin{align}\label{eq:liminf_partial_result}
        \frac{d\mu}{d\lambda}(x_0) \geq (1-\delta)^2 \liminf_{m \to \infty} \liminf_{n \to \infty} \frac{1}{\rho_m^{N-1}} \sum_{l \in \{i,j\}} \int_{Q_l(\rho_m)} 2\sqrt{{\widetilde W}_l(x_0, w_n)}|\nabla w_n| \, \mathrm{d}x.
    \end{align}
    Since 
    $$|(M_m \times (-\rho_m/2, \rho_m/2)) \setminus (Q_i(\rho_m) \cup Q_j(\rho_m))| = 0$$ 
    by \eqref{eq:cond_set_1} and \eqref{eq:cond_set_2}, we can derive
    \begin{align}\label{eq:liminf_partial_inequality_6}
        &\sum_{l \in \{i,j\}} \int_{Q_l(\rho_m)} \frac{1}{\eps_n} \widetilde W_l(x_0, T_{a_l}(x_0) v_n) + \eps_n |\nabla \tilde u_n|^2 \, \mathrm{d}x \nonumber \\
        & \hspace{1 cm} \geq \int_{M_m} \int_{-1/2}^{h(x')} 2\sqrt{{\widetilde W}_j(x_0, w_n)}|\nabla w_n| \, \mathrm{d}t + \int_{h(x')}^{1/2} 2\sqrt{{\widetilde W}_i(x_0, w_n)}|\nabla w_n| \, \mathrm{d}t \, \mathrm{d}x' \nonumber \\
        & \hspace{1 cm} \geq \int_{M_m} d_{\widetilde W_j}(x_0, w_n(x', -\rho_m/2), w_n^-(x', h(x'))) + d_{\widetilde W_i}(x_0, w_n^+(x', h(x')), w_n(x', \rho_m/2)) \, \mathrm{d}x',
    \end{align}
    where $w_n^+, w_n^-$ are the one-sided trace values at $(x', h(x'))$.
    Now, we observe that 
    \begin{align*}
        w_n(\cdot , -\rho_m/2) = T_{a_j}(x_0)S_{a_j}(\cdot , -\rho_m/2)\tilde u_n(\cdot, -\rho_m/2) \xrightarrow{n \to \infty} T_{a_j}(x_0)S_{a_j}(\cdot , -\rho_m/2)u(\cdot, -\rho_m/2)
    \end{align*}
    in $L^1(Q'(\rho_m))$ (analogously for $w_n(\cdot , \rho_m/2)$). Moreover, we have 
    \begin{align*}
        w_n^-(x', h(x')) = T_{a_j}(x_0)S_{a_j}((x', h(x')) u_n(x', h(x'))
    \end{align*}
    and 
    \begin{align*}
        w_n^+(x', h(x')) = T_{a_i}(x_0)S_{a_i}((x', h(x')) u_n(x', h(x')).
    \end{align*}
    Now, we set 
    \[
        z^-_m(x') := T_{a_j}(x_0)S_{a_j}(x' , -\rho_m/2)u(x', -\rho_m/2),
    \]
    \[
        z^+_m(x') := T_{a_i}(x_0)S_{a_i}(x' , \rho_m/2)u(x', \rho_m/2),
    \]
    \[
        L^{j,m}_{x'} := T_{a_j}(x_0)S_{a_j}((x', h(x')),
    \]
    and 
    \[
        L^{i,m}_{x'} := T_{a_i}(x_0)S_{a_i}((x', h(x')).
    \]
    Since  $d_{W_l}$ is Lipschitz on compact sets $K \subset \R^M$ (cf.\ \ref{lemma:geodesic_distance_is_locally_lipschitz}) we obtain
    \begin{align}\label{eq:liminf_technical_equation_1}
        \inf_{r \in \R^M} d_{\widetilde W_j}(x_0, p, L^{j,m}_{x'}r) + d_{\widetilde W_i}(x_0, L^{i,m}_{x'}r, q) \leq C_K(1 + |p| + |q|). 
    \end{align}
    for any $p, q \in K$. Since the sequence $\{w_n\}$ is uniformly bounded in $L^\infty$ we can apply \eqref{eq:liminf_technical_equation_1} in conjunction with the Dominated Convergence Theorem to infer 
    \begin{align}\label{eq:liminf_partial_inequality_7}
        &\lim_{n \to \infty}\int_{M_m} d_{\widetilde W_j}(x_0, w_n(x', -\rho_m/2), w_n^-(x', h(x'))) + d_{\widetilde W_i}(x_0, w_n^+(x', h(x')), w_n(x', \rho_m/2)) \, \mathrm{d}x' \nonumber \\
        & \hspace{1cm }\geq \lim_{n \to \infty} \int_{M_m} \inf_{r \in \R^M} \left( d_{\widetilde W_j}(x_0, w_n(x', -\rho_m/2), L^{j,m}_{x'}r) + d_{\widetilde W_i}(x_0, L^{i,m}_{x'}r, w(x', \rho_m/2)) \right) \, \mathrm{d}x' \nonumber \\
        & \hspace{1cm} = \int_{M_m} \inf_{r \in \R^M} \left( d_{\widetilde W_j}(x_0, z_m^-(x'), L^{j,m}_{x'}r) + d_{\widetilde W_i}(x_0, L^{i,m}_{x'}r, z_m^+(x')) \right) \, \mathrm{d}x'.
    \end{align}
    Again, using the uniform boundedness of the integrand in $L^\infty$ with \eqref{eq:relative_large_good_set}, we obtain
    \begin{align}\label{eq:liminf_partial_inequality_8}
        &\int_{M_m} \inf_{r \in \R^M} \left( d_{\widetilde W_j}(x_0, z_m^-(x'), L^{j,m}_{x'}r) + d_{\widetilde W_i}(x_0, L^{i,m}_{x'}r, z_m^+(x')) \right) \, \mathrm{d}x' \nonumber \\
        & \hspace{1cm} \geq \int_{Q'(\rho_m)} \inf_{r \in \R^M} \left( d_{\widetilde W_j}(x_0, z_m^-(x'), L^{j,m}_{x'}r) + d_{\widetilde W_i}(x_0, L^{i,m}_{x'}r, z_m^+(x')) \right) \, \mathrm{d}x' - C\delta\rho_m^{N-1}.
    \end{align}
    If we rescale now the integral by $1/\rho_m^{N-1}$, we derive 
    \begin{align}\label{eq:liminf_partial_inequality_9}
        &\frac{1}{\rho_m^{N-1}}\int_{Q'(\rho_m)} \inf_{r \in \R^M} \left( d_{\widetilde W_j}(x_0, z_m^-(x'), L^{j,m}_{x'}r) + d_{\widetilde W_i}(x_0, L^{i,m}_{x'}r, z_m^+(x')) \right) \, \mathrm{d}x \nonumber \\
        & \hspace{1cm} = \int_{Q'} \inf_{r \in \R^M} \left( d_{\widetilde W_j}(x_0, w_m^-(\rho_m x'), L^{j,m}_{\rho_m x'}r) + d_{\widetilde W_i}(x_0, L^{i,m}_{\rho_m x'}r, w_m^+(\rho_m x')) \right) \, \mathrm{d}x. 
    \end{align}

    Now again, since $w_m^\pm$ are uniformly bounded in $L^\infty$, also the $\inf$ is uniformly attained in a large ball (cf.\ Proposition \ref{prop:infimum_attained_in_ball}). Therefore, we notice that the integrand is uniformly bounded in $L^\infty(Q')$. By \eqref{eq:conv_1} and \eqref{eq:conv_2}, and the continuity of $S_{a_l}$ we also obtain
    \[
        w_m^+(\rho_m \cdot ) \xrightarrow{m \to \infty} a_i(x_0) \qquad \text{ and } \qquad  w_m^-(\rho_m \cdot ) \xrightarrow{m \to \infty} a_j(x_0)
    \]
    in $L^1(Q', \R^M)$. In particular, we have pointwise convergence almost everywhere. Now, we also observe that for every $x' \in Q'$ we have 
    \[
        L^{i,m}_{\rho_m x'} \xrightarrow{m \to \infty} Id_{M} \qquad \text{ and } \qquad L^{j,m}_{\rho_m x'} \xrightarrow{m \to \infty} Id_{M}.
    \]
    By applying Lemma \ref{lemma:technical}, for $\H^{N-1}$-a.e. $x' \in Q'$, it follows that
    \begin{align*}
        &\inf_{r \in \R^M} \left( d_{\widetilde W_j}(x_0, w_m^-(\rho_m x'), L^{j,m}_{\rho_m x'}r) + d_{\widetilde W_i}(x_0, L^{i,m}_{\rho_m x'}r, w_m^+(\rho_m x')) \right)  \\
        & \hspace{3cm} \xrightarrow{m \to \infty}  \inf_{r \in \R^M} \left( d_{\widetilde W_j}(x_0, a_j(x_0), r) + d_{\widetilde W_i}(x_0, r, a_i(x_0)) \right).
    \end{align*}
    Applying the dominating convergence theorem, we infer 
    \begin{align}\label{eq:liminf_partial_inequality_10}
        & \lim_{m \to \infty} \int_{Q'} \inf_{r \in \R^M} \left( d_{\widetilde W_j}(x_0, w_m^-(\rho_m x'), L^{j,m}_{\rho_m x'}r) + d_{\widetilde W_i}(x_0, L^{i,m}_{\rho_m x'}r, w_m^+(\rho_m x')) \right) \, \mathrm{d}x' \nonumber \\
        & = \inf_{r \in \R^M} \left(  d_{\widetilde W_j}(x_0, a_j(x_0), r) + d_{\widetilde W_i}(x_0, r, a_i(x_0))\right) \nonumber \\
        & = d_{\widetilde W}(x_0, u^-(x_0), u^+(x_0)). 
    \end{align}
    Gathering \eqref{eq:liminf_partial_result}, \eqref{eq:liminf_partial_inequality_6}, \eqref{eq:liminf_partial_inequality_7}, \eqref{eq:liminf_partial_inequality_8}, \eqref{eq:liminf_partial_inequality_9}, \eqref{eq:liminf_partial_inequality_10}, we obtain  
    \begin{align*}
        \frac{d\mu}{d\lambda}(x_0) \geq (1-\delta)^2d_{\widetilde W}(x_0, u^-(x_0), u^+(x_0)) - C\delta = (1-\delta)^2d_{ W}(x_0, u^-(x_0), u^+(x_0)) - C\delta. 
    \end{align*}
    by our choice of truncation constants (cf.\ Corollary \ref{corol:truncation_constant}) and with $C > 0$ not depending on $\delta$.  
    Since $\delta > 0$ was arbitrary, we conclude.
\end{proof}



%%%%%%%%%%%%%%%%%%%%%%%%%%%%%%%%%%%%%%%%%%%%
%%%%%%%%%%%%%%%%%%%%%%%%%%%%%%%%%%%%%%%%%%%%
%%%%%%%%%%%%%%%%%%%%%%%%%%%%%%%%%%%%%%%%%%%%
%%%%%%%%%%%%%%%%%%%%%%%%%%%%%%%%%%%%%%%%%%%%



\section{Limsup inequality}\label{sec:limsup}

This section is devoted to the proof of the limsup inequality.
we start by proving the result in the case with no mass constraint.

\begin{theorem}\label{thm:limsup}
    Let $\{ \varepsilon_n \}_n$ be an infinitesimal sequence. 
    Assume that \ref{H1}-\ref{H8} hold.
    Let $u \in L^1(\Omega; \R^M)$. Then, there exists $\{u_n\}_n \subset H^1(\Omega; \R^M)$ such that 
    \[
        \mathcal{F}_\infty(u) = \lim_{n \to \infty} \mathcal{F}_{\eps_n}(u_n).
    \]
    and $u_n \to u$ in $L^1(\Omega;\R^M)$.
\end{theorem}

To prove this, we need the following preliminary result.

\begin{lemma}\label{lemma:reparameterization}
    Fix $\lambda > 0$, $\varepsilon > 0$, $x \in \Omega$, and $p,q \in \R^M$. Let $\gamma \in C^1([-1,1]; \R^M)$ with $\gamma'(s) \neq 0$ for all $s \in (-1,1)$, $\gamma(-1) = p$, $\gamma(1) = q$. Then, there exist $\tau > 0$ and $C > 0$ with
    $$ C \varepsilon \leq \tau \leq \frac{\varepsilon}{\sqrt{\lambda}} \int_{-1}^1 | \gamma'(s)| \dd s, $$
    and $g \in C^1([0, \tau]; [-1,1])$ such that
    \begin{equation}\label{eq:reparameterization_speed}
        (g'(t))^2 = \frac{\lambda + W(x,\gamma(g(t)))}{\varepsilon^2 | \gamma'(g(t))|^2},
    \end{equation}
    for all $t \in (0,\tau)$, $g(0) = -1$, $g(\tau) = 1$, and
    \begin{align*}
        \int_0^\tau \left[ \frac{1}{\varepsilon} W(x, \gamma(g(t))) + \varepsilon |\gamma'(g(t))|^2 | g'(t)|^2 \right] \dd t \leq \int_{-1}^1 2 \sqrt{W(x, \gamma(s))} | \gamma'(s) | \dd s + 2 \sqrt{\lambda} \int_{-1}^1 | \gamma'(s) |\dd s.
    \end{align*}
\end{lemma}

A proof of this result can be found in \cite[Lemma 6.3]{CriGra} (see also \cite[Proposition 2]{modica87} and \cite[Lemma 3.2]{baldo}).

\begin{proof}[Proof of Theorem \ref{thm:limsup}]
We divide the proof into several steps.\\[0.5em]
\textbf{Step 0: Approximation with $C^2$ sets.} 
We use Proposition \ref{prop:approx_sets} to approximate the jump set of $u$ with $C^2$ sets such that they coincide with the $C^2$ jump sets of the wells as much as possible. 
We define the singular set $S$ of the approximated jump set as the set of points belonging to the approximation for which either the measure-theoretic exterior normal is not defined, or its norm is not equal to $1$. Since S has finite $\mathcal{H}^{N-2}$-measure, we can define 
$$ S_n \coloneqq \mathcal{N}_{C\varepsilon_n} (S) = \{ x \in \Omega : | x - y | < C \varepsilon_n \text{ for } y \in S \}, $$
and it follows that
$$ \mathcal{L}^N (S_n) \leq D \varepsilon_n^2, $$
where the constant $D > 0$ depends only on $\mathcal{H}^{N-2} (S)$.

Therefore, considering only $\Omega \setminus S_n$, we are effectively isolating the singularities, and we can now assume that each interface $\Sigma$ is a $C^2$ parameterized hypersurface, where the normals are defined and the distance function from the interface is $C^1$ in a tubular neighborhood around it. In the following steps, we will carry out the construction of the recovery sequence and the relative computations only for the local case of a single $C^2$ interface. The global construction can be carried out in a similar way as in \cite[Lemma 4.6]{CriGra}, by using interpolation between the local constructions built below. We decided not to detailed this standard construction, and focus on the ideas for building the recovery sequence close to the regular part of the interfaces.\\[0.5em]
\textbf{Step 1: Construction of a local recovery sequence.}
For $0 < \alpha < \beta < 1$ define
\begin{equation}
    r_n \coloneqq \varepsilon_n^\alpha, \qquad \eta_n \coloneqq \varepsilon_n^\beta.
\end{equation}
This choice allows us to have $\varepsilon_n \ll \eta_n \ll r_n$.
Thanks to the assumption $a,b \in L^\infty $, using Theorem \ref{existence-of-geodesics-with-bounded-path-length} we know that there exists a constant $\overline{L} > 0$ such that the geodesics satisfy $L(\gamma) \leq \overline{L}$, where $L(\gamma)$ is the length of the geodesic $\gamma$. We therefore also define
\begin{equation}
    \ell_n \coloneqq \overline{L} \varepsilon_n.
\end{equation}
We are assuming that $\Sigma$ is the graph of a $C^2$ function, therefore there exist an open set $D \subset \R^{N-1}$, and a $C^2$ function $\psi \colon D \to \R$ such that $ \psi(D) = \Sigma $. \\
For $y \in D$, we define the normal to the interface at $\psi(y)$ as $\nu(y)$. With this, we define the function $\Psi \colon D \times [-\ell_n, \ell_n] \to \R^N$ as 
$$ \Psi(y,t) \coloneqq \psi(y) + t \nu(y). $$
For $n$ large enough, the tubular neighborhood theorem guarantees that this map is a $C^1$ diffeomorphism, and that every point in it has a unique projection on the interface, meaning that $\Psi(y,t)$ get projected to $\psi(y)$. \\
Now, take a lattice of points $\{y_i\}_{i \in I_n} \subset D$ with spacing $2r_n$, and define the set
\begin{equation}
    U_n^i \coloneqq \{ x \in \Omega : \exists y \in Q'(y_i,r_n), \exists t \in [-\ell_n, \ell_n] : \Psi(y,t) = x \},
\end{equation}
as the local tubular neighborhood around $\psi(y_i)$.
\begin{figure}[htbp]
        \includegraphics[width=0.9\linewidth]{tubular_neighborhood}
    \caption{Partition of the tubular neighborhood.}
    \label{fig:tub_neighborhood}
\end{figure}

We observe that the, for the index set $I_n$, it holds that
\begin{equation}\label{eq:estimate_index_set}
    \#| I_n | \leq C \frac{1}{r_n^{N-1}},
\end{equation}
where the constant $C$ does not depend on $n$.

Take also $\{\varphi_n^i\}_{n, i \in I_n} \subset C^\infty (D)$ a partition of unity such that
\begin{equation}
    0 \leq \varphi_n^i \leq 1, \qquad \varphi_n^i \equiv 1 \in Q'(y_i,r_n-\eta_n), \qquad \varphi_n^i \equiv 0 \in D \setminus Q'(y_i,r_n+\eta_n)
\end{equation}
\begin{equation}
    \sum_{i \in I_n} \varphi_n^i(y) \equiv 1 \;\;\text{ in }\; D, \qquad | \nabla \varphi_n^i | \leq \frac{C}{2\eta_n}.
\end{equation}
Let $\gamma_i$ be the geodesic given by Theorem \ref{existence-of-geodesics-with-bounded-path-length} relative to $\psi(y_i)$, namely a curve in $C^1([-1,1]; \R^M)$ with $\gamma_i (-1) = u^- (\psi(y_i))$, $\gamma_i (1) = u^+ (\psi(y_i))$ and $\gamma_i'(s) \neq 0$, such that
\begin{equation}
    L(\gamma_i) \leq \overline{L}, \qquad \int_{-1}^1 2 \sqrt{W(\psi(y_i), \gamma_i(s))} |\gamma_i'(s)|\dd s = \dd_W(\psi(y_i), u^+(\psi(y_i)), u^-(\psi(y_i))).
\end{equation}
Now, we apply Lemma \ref{lemma:reparameterization} with
$$ \lambda = \varepsilon_n^2, \qquad \varepsilon = \varepsilon_n, \qquad \gamma = \gamma_i $$
to find $\tau_n^i < \ell_n$ and $g_n^i \in C^1(-\tau_n^i, \tau_n^i)$ such that
\begin{equation}\label{eq:reparameterization_limsup}
    ((g_n^i)'(t))^2 = \frac{\varepsilon_n^2 + W(x,\gamma_i(g_n^i(t)))}{\varepsilon^2 | \gamma_i'(g_n^i(t))|^2},
\end{equation}
and
\begin{equation*}
    \int_0^{\tau_n^i} \! \left[ \frac{1}{\varepsilon} W(x, \gamma_i(g_n^i(t))) + \varepsilon |\gamma_i'(g_n^i(t))|^2 | (g_n^i)'(t)|^2 \right] \! \dd t\leq \int_{-1}^1 2 \sqrt{W(x, \gamma_i(s))} | \gamma_i'(s) |  \dd s + 2 \overline{L} \varepsilon_n. %\label{eq:reparameterization_limsup}
\end{equation*}
This result, together with Theorem \ref{existence-of-geodesics-with-bounded-path-length}, also implies that the reparametrized geodesics are bounded
\begin{equation}\label{eq:limsup_bounded_geodesics}
    \| \gamma_i \circ g_n^i \|_\infty \leq C < + \infty.
\end{equation}
This in turn, together with \ref{H4} and Lemma \ref{lemma:reparameterization}, gives us
\begin{equation}\label{eq:limsup_geodesics_1overepssquared}
    | \gamma_i'(g_n^i(t)) | \, |(g_n^i)'(t) | \leq \frac{C}{\varepsilon_n^2}.
\end{equation}
Extend $g_n^i$ to the whole real line by setting $g_n^i(s) \equiv 1$ for $s \geq \tau_n^i$ and $g_n^i(s) \equiv -1$ for $s \leq -\tau_n^j$.\\
We define $v_n^i$ in $ D \times [-\ell_n, \ell_n] $ using the maps $T_a$ defined in Remark \ref{rem:T_symmetric}, as
\begin{equation}
    v_n^i(y,t) \coloneqq T_a (\Psi(y,t)) T_a^{-1} (\psi(y_i)) \gamma_i(g_n^i(t)).
\end{equation}
The reason for this choice is that this respects the boundary values, since 
\begin{align}\label{eq:vni_positive_boundary}
    v_n^i(y,\ell_n) &= T_a (\Psi(y,\ell_n)) T_a^{-1} (\psi(y_i)) \gamma_i(g_n^i(\ell_n)) \nonumber\\
    &= T_a (\Psi(y,\ell_n)) T_a^{-1} (\psi(y_i)) a ( \psi (y_i)) \nonumber\\
    &= T_a (\Psi(y,\ell_n)) e_1 \nonumber\\
    &= a(\Psi(y, \ell_n)),
\end{align}
and analogously 
\begin{equation}\label{eq:vni_negative_boundary}
    v_n^i(y, -\ell_n) = -a(\Psi(y,-\ell_n)).
\end{equation} Then, for $y \in D$ and $t \in [-\ell_n, \ell_n]$ we define
\begin{equation}
    v_n(y,t) \coloneqq \sum_{i \in I_n} \varphi_n^i(y) v_n^i (y,t).
\end{equation}
This sequence still respects the boundary values, since
\begin{align*}
    v_n (y, \ell_n) &= \sum_{i \in I_n} \varphi_n^i(y) v_n^i (y,\ell_n) \\
    &= \sum_{i \in I_n} \varphi_n^i(y) a (\Psi(y, \ell_n)) \\
    &= a (\Psi(y, \ell_n)) \sum_i \varphi_n^i(y) \\
    &= a (\Psi(y, \ell_n)),
\end{align*}
where in the last step we used \eqref{eq:vni_positive_boundary}. Analogously it holds
$$ v_n (y, - \ell_n) = - a (\Psi(y, -\ell_n)).$$
Also, thanks to $a \in L^\infty$, it is clear that 
\begin{equation}\label{eq:vn_bounded}
    \sup_{n \in \N} \| v_n \|_\infty \leq C < + \infty.
\end{equation}
Since $\Psi$ is a diffeomorphism, we can define the local recovery sequence $u_n$ around $\Sigma$ as
\begin{gather}\label{eq:recovery_sequence}
    u_n(x) \coloneqq \begin{cases}
        a(x) &\qquad \mathrm{dist}_\Sigma(x) > \ell_n,\\
        v_n(\Psi^{-1}(x)) &\qquad |\mathrm{dist}_\Sigma(x)| \leq \ell_n,\\
        -a(x) &\qquad \mathrm{dist}_\Sigma(x) < - \ell_n.
    \end{cases}
\end{gather}
Here, $\mathrm{dist}_\Sigma$ denotes the signed distance from $\Sigma$.
Note that this construction implies $u_n \in L^1(\Omega; \R^M)$.
Finally, we denote the tubular neighborhood by $\Sigma_n \coloneqq \Psi(D \times [-\ell_n, \ell_n])$ (see Figure \ref{fig:tub_neighborhood}).\\[0.5em]
\textbf{Step 2: Localization in strip.}
Since we define the recovery sequence to be equal to the wells outside of the $\ell_n$-tubular neighborhood of the interface, and the wells are in $W^{1,2}$, we get that it is only necessary to study what happens inside of the tubular neighborhood.\\
Indeed,
\begin{align*}
    \mathcal{F}_{\varepsilon_n}(u_n) &= \int_\Omega \left[ \frac{1}{\varepsilon_n} W(x, u_n) + \varepsilon_n | \nabla u_n|^2 \right] \dd x \\
    &= \int_{\Sigma_n} \left[ \frac{1}{\varepsilon_n} W(x, u_n) + \varepsilon_n | \nabla u_n|^2 \right] \dd x + \int_{\{ |\mathrm{dist}_\Sigma(x)| > \ell_n \}}  \varepsilon_n |\nabla a|^2 \dd x.
\end{align*}
The last term of the last line can be bounded by
$$ \int_{\{ |\mathrm{dist}_\Sigma(x)| > \ell_n \}}  \varepsilon_n |\nabla a|^2 \dd x \leq \varepsilon_n \| \nabla a \|^2_{L^2(\Omega; \R^{M \times N})} $$
which vanishes in the limit since $a \in W^{1,2}(\Omega; \R^M)$. We are therefore left with estimating the integral over the tubular neighborhood.
In order to change coordinates with the diffeomorphism $\Psi$, we need to see how the gradient changes under this reparameterization. 

Thanks to the definition of $\Psi$, we have that
\begin{gather*}
    J \Psi = ( J \psi + t J \nu \; | \; \nu ).
\end{gather*}
We now want to compute the determinant $\mathrm{det}(J\Psi)$.
Since $v_n (y,t) = u_n (\Psi(y,t))$, by the chain rule we get $ \nabla_{(y,t)} v_n = (J \Psi)^T \nabla_x u_n $, therefore
$$ \nabla_x u_n = ((J \Psi)^{T})^{-1} \nabla_{(y,t)} v_n. $$
To find $| \nabla_x u_n |^2$, we compute the dot product
$$ | \nabla_x u_n |^2 = (\nabla_{(y,t)} v_n)^T ( (J\Psi)^T J\Psi)^{-1} \nabla_{(y,t)} v_n. $$
The matrix $ (J\Psi)^T J\Psi$ yields
\begin{gather*}
    (J\Psi)^T J\Psi = \begin{pmatrix}
        (J\psi + t J \nu)^T (J\psi + t J \nu) & (J\psi + t J \nu)^T \nu \\
        \nu^T (J\psi + t J \nu) & \nu^T \nu
    \end{pmatrix}.
\end{gather*}
We denote the top left block by 
\begin{equation}
    G = (J\psi + t J \nu)^T (J\psi + t J \nu).
\end{equation}
Since $\nu$ is a unit normal vector, we have $\nu^T \nu = |\nu|^2 = 1$. Differentiating this identity we also find that $(J\nu)^T \nu = 0$. Moreover, since $\nu $ is orthogonal to the surface tangent vectors $J \psi$, we have $(J\psi)^T \nu = 0$. In particular, these imply that the matrix is equal to
\begin{gather*}
    (J\Psi)^T J\Psi = \begin{pmatrix}
        G & 0 \\
        0 & 1
    \end{pmatrix} \implies ((J\Psi)^T J\Psi)^{-1} = \begin{pmatrix}
        G^{-1} & 0 \\
        0 & 1
    \end{pmatrix}.
\end{gather*}
Splitting the gradient of $v_n$ into a tangential part and a normal part, we get
\begin{gather}
    | \nabla_x u_n |^2 = \begin{pmatrix}
        ( \nabla_y v_n)^T & \partial_t v_n
    \end{pmatrix} \begin{pmatrix}
        G^{-1} & 0 \\ 0 & 1
    \end{pmatrix} \begin{pmatrix}
        \nabla_y v_n \\ 
        \partial_t v_n
    \end{pmatrix} = \nabla_y v_n \cdot G^{-1} \nabla_y v_n + | \partial_t v_n |^2.
\end{gather}
Therefore, applying the change of coordinates $x = \Psi(y,t)$, we get
\begin{align*}
    &\int_{\Sigma_n} \left[ \frac{1}{\varepsilon_n} W(x, u_n) + \varepsilon_n | \nabla u_n|^2 \right] \dd x \\
    &= \sum_{i \in I_n} \int_{U_n^i} \left[ \frac{1}{\varepsilon_n} W(x, u_n) + \varepsilon_n | \nabla u_n|^2 \right] \dd x \\
    &= \sum_{i \in I_n} \int_{Q'(y_i, r_n)} \int_{-\ell_n}^{\ell_n} \left[ \frac{1}{\varepsilon_n} W(\Psi(y,t), v_n) + \varepsilon_n | \partial_t v_n|^2 + \varepsilon_n \nabla_y v_n \cdot G^{-1} (y,t) \nabla_y v_n \right] |\mathrm{det}(J\Psi)| \, \mathrm{d} y \mathrm{d} t.
\end{align*}
For brevity, we define
\begin{equation*}
    g_n (y,t) \coloneqq \frac{1}{\varepsilon_n} W(\Psi(y,t), v_n) + \varepsilon_n | \partial_t v_n|^2 + \varepsilon_n \nabla_y v_n \cdot G^{-1} (y,t) \nabla_y v_n.
\end{equation*}

Since $\psi$ is a $C^2$ map, and $\nu$ is $C^1$, we have that $J\psi + t J\nu$, and therefore $G$, is continuous, therefore its norm is bounded on compact sets. Moreover, thanks to tubular neighborhood theorem, $G$ is also always strictly positive definite. What these two facts imply is that for all $n \in \N$ and for all $y \in D, t \in (-\ell_n, \ell_n)$, there exist $C > 0$ and $t_0 \in (-\ell_n, \ell_n)$ such that $\forall t \in (-t_0,t_0)$ it holds
\begin{equation}\label{eq:tangential_derivative}
    \frac1C | \nabla_y v_n |^2 \leq \nabla_y v_n \cdot G^{-1} \nabla_y v_n \leq C | \nabla_y v_n |^2.
\end{equation}
A proof of this can be found in \cite[Lemma 2.3]{lotoreichik2023spectralasymptoticsdiracoperator}.
Thanks to the definition of the recovery sequence and the partition of unity $\varphi_n^i$, we can split the integral as
\begin{align*}
    \sum_{i \in I_n} \int_{Q'(y_i, r_n)} \int_{-\ell_n}^{\ell_n} g_n (y,t) | \mathrm{det}(J\Psi)| \,\mathrm{d}y \mathrm{d}t = &\sum_{i \in I_n} \int_{Q'(y_i,r_n-\eta_n)} \int_{-\ell_n}^{\ell_n} g_n (y,t) | \mathrm{det}(J\Psi)| \, \mathrm{d} y \mathrm{d} t \\
    + &\sum_{i \in I_n} \int_{Q'(y_i,r_n) \setminus Q'(y_i, r_n - \eta_n)} \int_{-\ell_n}^{\ell_n} g_n (y,t) | \mathrm{det}(J\Psi)| \, \mathrm{d} y \mathrm{d} t.
\end{align*}
We denote the first term of the right-hand side by $I_1^n$ and the second term by $I_2^n$. \\[0.5em]
\textbf{Step 3: Estimate of $I_2^n$.}
We prove that the contribution of $I_2$ is neglibile in the limit.\\
Thanks to Assumption \ref{H5}, together with \eqref{eq:vn_bounded}, we know that $\forall n \in \N$ it holds $W(\Psi(y,t), v_n) \leq C$, which implies
\begin{align*}
    &\sum_{i \in I_n} \int_{Q'(y_i,r_n) \setminus Q'(y_i, r_n - \eta_n)} \int_{-\ell_n}^{\ell_n} \cfrac{1}{\varepsilon_n} W(\Psi(y,t),v_n) | \mathrm{det}(J\Psi)|\, \dd y \dd t \\
    &\qquad \qquad \qquad\qquad \qquad\qquad\leq \frac{1}{r_n^{N-1}} \left[ r_n^{N-1} - (r_n - \eta_n)^{N-1} \right] C \frac{\ell_n}{\varepsilon_n} \\
    &\qquad \qquad \qquad\qquad \qquad\qquad= \overline{L} C \left[ 1 - \left( 1 - \frac{\eta_n}{r_n} \right)^{N-1} \right] \to 0,
\end{align*}
where we used $\ell_n = \overline{L} \varepsilon_n$ and \eqref{eq:estimate_index_set}.
This term vanishes in the limit since $\eta_n \ll r_n$.

For the gradient part of $I_2$, observe that for a point $y \in Q'(y_i,r_n) \setminus Q'(y_i, r_n - \eta_n)$, there exists a finite set of indexes $J \subset I_n$ independent of $y$ such that $\varphi_n^j(y) = 0$ for all $j \in I_n \setminus J$. In particular, the cardinality of the set $J $ depends only on the dimension and is uniform in $n$. Therefore for $I_2$ we have
$$ v_n(y,t) \coloneqq \sum_{i \in J} \varphi_n^i(y) T_a (\Psi(y,t)) T_a^{-1} (\psi(y_i)) \gamma_i(g_n^i(t)). $$

We first treat the derivative with respect to $t$. For this we have
\begin{align*}
    | \partial_t v_n (y,t)|^2 &= \left| \partial_t \sum_{i \in J(y)} \varphi_n^i(y) v_n^i (y,t) \right|^2 \\
    &\leq \sum_{i \in J(y)} |\partial_t v_n^i(y,t)|^2 \\
    &\leq \sum_{i \in J(y)} C | \nabla T_a (\Psi(y,t)) \cdot n(y) \gamma_i (g_n^i(t)) + T_a (\Psi(y,t)) \gamma_i'(g_n^i(t)) (g_n^i)'(t) |^2 \\
    &\leq C | \nabla T_a (\Psi(y,t)) |^2 + \frac{C}{\varepsilon_n^2} |T_a (\Psi(y,t))|^2,
\end{align*}
where we used that $| \varphi_n^i | \leq 1$, that $| J(y) | $ is uniform in $n$, together with \eqref{eq:limsup_bounded_geodesics} and \eqref{eq:limsup_geodesics_1overepssquared}.
This implies 
\begin{align}
    &\sum_{i \in J(y)} \int_{Q'(y_i,r_n) \setminus Q'(y_i, r_n - \eta_n)} \int_{-\ell_n}^{\ell_n} \varepsilon_n | \partial_t v_n|^2 | \mathrm{det}(J \Psi) | \,\dd y \dd t \nonumber\\
    &\qquad\qquad\leq C \varepsilon_n \| \nabla T_a \|^2_{L^2(\Omega)} + C\left[ 1 - \left( 1 - \frac{\eta_n}{r_n} \right)^{N-1} \right] \frac{\ell_n}{\varepsilon_n} \| T_a\|^2_\infty \to 0,
\end{align}
where this last term vanishes since $\ell_n = \overline{L} \varepsilon_n$, together with $\mathrm{det}(J \Psi)$ being uniformly bounded, \eqref{eq:T_a_properties} and since
$$  1 - \left( 1 - \frac{\eta_n}{r_n} \right)^{N-1} \leq 2 (N-1) \frac{\eta_n}{r_n}, $$
and we chose $\eta_n, r_n$ such that $\eta_n \ll r_n$.

For the term involving the tangential derivative, using \eqref{eq:tangential_derivative} we just need to bound $|\nabla_y v_n|^2$. Therefore, using analogous reasoning as the normal derivative, we have
\begin{align*}
    | \nabla_y v_n(y,t)|^2 &= \left| \nabla_y \sum_{i \in J(y)} \varphi_n^i(y) v_n^i (y,t) \right|^2 \\
    &\leq C \sum_{i \in J(y)} | \nabla_y \varphi_n^i(y) T_a (\Psi(y,t)) + \varphi_n^i(y) (J \psi + t Jn) \nabla_y T_a (\Psi(y,t)) |^2\\
    &\leq C | \nabla T_a (\Psi(y,t))|^2 + \frac{C}{\eta_n^2} | T_a (\Psi(y,t)) |^2,
\end{align*}
where we also used $| \nabla \varphi_n^i | \leq \frac{1}{2 \eta_n}$, together with $J\psi + t J \nu$ being uniformly bounded.
This implies
\begin{align}
    &\sum_{i \in J(y)} \int_{Q'(y_i,r_n) \setminus Q'(y_i, r_n - \eta_n)} \int_{-\ell_n}^{\ell_n} \varepsilon_n \nabla_y v_n G^{-1} \nabla_y v_n | \mathrm{det}(J\Psi)|\, \dd y \dd t \nonumber\\
    &\qquad\qquad\qquad\leq \varepsilon_n \| \nabla T_a \|^2_{L^2(\Omega)} + \left[ 1 - \left( 1 - \frac{\eta_n}{r_n} \right)^{N-1} \right] \frac{\ell_n \varepsilon_n}{\eta_n^2} \| T_a \|^2_\infty \to 0,
\end{align}
where we used that $\mathrm{det}(J \Psi)$ is bounded, together with \eqref{eq:T_a_properties}.
This last term vanishes since 
$$ \left[ 1 - \left( 1 - \frac{\eta_n}{r_n} \right)^{N-1} \right] \frac{\ell_n \varepsilon_n}{\eta_n^2} \leq C \frac{\varepsilon_n^2}{\eta_n^2}, $$
and we chose $\varepsilon_n, \eta_n$ such that $\varepsilon_n \ll \eta_n$. We can therefore conclude.\\[0.5em]
\textbf{Step 4: estimate of $I_1^n$.}
We want to prove that $I_1^n$ converges to our desired limit, namely
\begin{align*}
    &\sum_{i \in I_n} \int_{Q'(y_i,r_n-\eta_n)} \int_{-\ell_n}^{\ell_n} \left[ \frac{1}{\varepsilon_n} W(\Psi(y,t), v_n) + \varepsilon_n | \partial_t v_n|^2 + \varepsilon_n \nabla_y v_n G^{-1} \nabla_y v_n \right] | \mathrm{det}(J\Psi)|\, \dd y \dd t \\
    &\to \int_D \dd_W (\psi(y), u^+(\psi(y)),u^-(\psi(y))) | \mathrm{det}(J \psi)|\, \dd \mathcal{H}^{N-1}(y)\\
    &= \int_\Sigma \dd_W (x, u^+ (x), u^- (x)) \, \dd \mathcal{H}^{N-1}(x),
\end{align*}
where $|\mathrm{det}(J\psi)|$ is the surface area element relative to the parameterization $\psi$, which for a $C^2$ graph is equal to $\sqrt{1 + |\nabla \psi |^2}$.

Thanks to Step 3, we can equivalently consider the following integral
$$ \sum_{i \in I_n} \int_{Q'(y_i,r_n)} \int_{-\ell_n}^{\ell_n} \left[ \frac{1}{\varepsilon_n} W(\Psi(y,t), v_n) + \varepsilon_n | \partial_t v_n|^2 + \varepsilon_n \nabla_y v_n G^{-1} \nabla_y v_n \right]| \mathrm{det}(J\Psi)|\, \dd y \dd t, $$
as effectively we are adding a vanishing term.\\
Using the relative continuity condition \ref{H7}, we can rewrite the potential term as
\begin{align*}
    &\sum_{i \in I_n} \int_{Q'(y_i,r_n)} \int_{-\ell_n}^{\ell_n} \frac{1}{\varepsilon_n} W(\Psi(y,t), v_n)| \mathrm{det}(J\Psi)|\, \dd y \dd t \\
    &\qquad= \sum_{i \in I_n} \int_{Q'(y_i,r_n)} \int_{-\ell_n}^{\ell_n} \frac{1}{\varepsilon_n} W(\Psi(y,t), T_a ( \Psi(y,t) ) T_a^{-1}(\Psi(y,t))v_n)| \mathrm{det}(J\Psi)|\, \dd y \dd t \\
    &\qquad\leq \sum_{i \in I_n} \int_{Q'(y_i,r_n)} \int_{-\ell_n}^{\ell_n} \frac{1}{\varepsilon_n} W(\psi(y_i), T_a(\psi(y_i)) T_a^{-1}(\Psi(y,t))v_n) | \mathrm{det}(J\Psi)|\, \dd y \dd t \\
    &\qquad\quad+ \sum_{i \in I_n} \int_{Q'(y_i,r_n)} \int_{-\ell_n}^{\ell_n} \cfrac{1}{\varepsilon_n} \omega (|\Psi(y,t) - \Psi(y_i,0)|) W(\Psi(y,t),v_n) | \mathrm{det}(J\Psi)|\, \dd y \dd t,
\end{align*}
where $\omega$ is the modulus of continuity from \ref{H7}.
We claim that the last term vanishes in the limit. Indeed, since $\Psi$ is $C^1$, and for $(y,t) \in Q'(y_i,r_n) \times [-\ell_n, \ell_n]$ we get 
$$ |(y,t)-(y_i,0)| \leq \sqrt{(N-1)r_n^2 + \ell_n^2} \eqqcolon f_n \to 0, $$
and therefore
$$ \omega (|\Psi(y,t) - \Psi(y_i,0)|) \leq \omega( \omega_\Psi (|(y,t) - (y_i,0)|)) \leq \omega ( \omega_\Psi (f_n)) \to 0. $$
This and assumption \ref{H5} imply that
\begin{align*}
    &\sum_{i \in I_n} \int_{Q'(y_i,r_n)} \int_{-\ell_n}^{\ell_n} \cfrac{1}{\varepsilon_n} \omega (|\Psi(y,t) - \Psi(y_i,0)|) W(\Psi(y,t),v_n)| \mathrm{det}(J\Psi)| \dd y \dd t \\
    &\qquad \qquad \qquad \qquad \qquad \qquad \qquad\leq C \cfrac{1}{r_n^{N-1}} r_n^{N-1} \cfrac{\ell_n}{\varepsilon_n} \omega( \omega_\Psi(f_n) ) \to 0.
\end{align*}
We are therefore left with estimating
\begin{align}
    &\sum_{i \in I_n} \int_{Q'(y_i,r_n)} \int_{-\ell_n}^{\ell_n} \Bigl[ \frac{1}{\varepsilon_n} W(\psi(y_i), T_a(\psi(y_i)) T_a^{-1}(\Psi(y,t))v_n) \nonumber \\
    &\hspace{6cm}+ \varepsilon_n | \partial_t v_n|^2 + \varepsilon_n \nabla_y v_n G^{-1} \nabla_y v_n \Bigr] | \mathrm{det}(J\Psi)| \dd y \dd t\label{eq:limsup_estimate_dontknow}
\end{align}
We define 
$$ w_n(y,t) \coloneqq T_a^{-1}(\Psi(y,t)) v_n(y,t). $$
Let us fix $d > 0$. For all $a,b \geq 0$, since $\left( \frac{a}{\sqrt{d}} - \sqrt{d} b \right)^2 \geq 0$, we get $2 a b \leq \frac1d |a|^2 + d |b|^2$. Thanks to this, we get
\begin{equation}\label{eq:minkowski_trick}
    | a + b |^2 \leq \left( 1 + \frac1d \right) |a|^2 + \left( 1 + d \right) |b|^2.
\end{equation}
Using this, together with \eqref{eq:T_a_properties}, for the normal derivative we get
\begin{align}
    \varepsilon_n | \partial_t v_n |^2 &= \varepsilon_n | \nabla T_a (\Psi(y,t)) \cdot n(y) w_n (y,t) + T_a(\Psi(y,t)) \partial_t w_n (y,t) |^2 \nonumber\\
    &\leq C \left( 1 + \frac1d \right) \varepsilon_n |\nabla T_a(\Psi(y,t))|^2 + C\left( 1 + d \right) \varepsilon_n |a(\Psi(y,t))|^2 |\partial_t w_n(y,t)|^2 \label{eq:change_normal}.
\end{align}
For the tangential derivative, we use again \eqref{eq:tangential_derivative}, together with \eqref{eq:minkowski_trick}, to get
\begin{align}
    \varepsilon_n \nabla_y v_n G^{-1} \nabla_y v_n &\leq C \varepsilon_n | \nabla_y v_n |^2 \nonumber\\
    &= C \varepsilon_n | (J \psi + t Jn) \nabla T_a (\Psi(y,t)) w_n (y,t) + T_a(\Psi(y,t)) \nabla_y w_n (y,t) |^2\nonumber\\
    &\leq C \left( 1 + \frac1d \right) \varepsilon_n |\nabla T_a(\Psi(y,t))|^2. \label{eq:change_tangential}
\end{align}
Here, we do not consider the term $\nabla_y w_n (y,t)$, since it turns out that, given our construction of the recovery sequence, we get
\begin{equation*}
    w_n (y,t) = T_a^{-1}(\Psi(y,t)) v_n(y,t) = T_a (\psi(y_i)) \gamma_i (g_n^i(t)) \implies \nabla_y w_n = 0.
\end{equation*}
Using \eqref{eq:change_normal} and \eqref{eq:change_tangential} in \eqref{eq:limsup_estimate_dontknow}, we get
\begin{align}
    &\sum_{i \in I_n} \int_{Q'(y_i,r_n)} \int_{-\ell_n}^{\ell_n} \bigg[ \frac{1}{\varepsilon_n} W(\psi(y_i), T_a(\psi(y_i)) T_a^{-1}(\Psi(y,t)) v_n) + \varepsilon_n | \partial_t v_n|^2 \nonumber \\
    &\hspace{8cm}+ \varepsilon_n \nabla_y v_n G^{-1} \nabla_y v_n \bigg] | \mathrm{det}(J\Psi)| \dd y \dd t \nonumber \\
    &\leq (1+d) \sum_{i \in I_n} \int_{Q'(y_i,r_n)} \int_{-\ell_n}^{\ell_n} \bigg[ \frac{1}{\varepsilon_n} W(\psi(y_i), T_a(\psi(y_i))w_n) \nonumber \\
    &\hspace{8cm}+ \varepsilon_n |a(\Psi(y,t))|^2 | \partial_t w_n|^2 \bigg] | \mathrm{det}(J\Psi)| \dd y \dd t \label{eq:limsup_remaining_term} \\
    &\qquad+ \left( 1 + \frac1d \right) \sum_{i \in I_n} \int_{Q'(y_i,r_n)} \int_{-\ell_n}^{\ell_n} \varepsilon_n | \nabla T_a (\Psi(y,t)) |^2 | \mathrm{det}(J\Psi)| \dd y \dd t. \label{eq:term_last_line}
\end{align}
The term in \eqref{eq:term_last_line} vanishes since
$$ \left( 1 + \frac1d \right) \sum_{i \in I_n} \int_{Q'(y_i,r_n)} \int_{-\ell_n}^{\ell_n} \varepsilon_n | \nabla T_a(\Psi(y,t)) |^2 | \mathrm{det}(J\Psi)| \dd y \dd t \leq \left( 1 + \frac1d \right) \varepsilon_n \| \nabla T_a \|^2_{L^2(\Omega)} \to 0. $$
Noticing that $T_a(\psi(y_i)) w_n(y,t) = \gamma_i (g_n^i(t)) \eqqcolon z_n^i (t)$, we can rewrite \eqref{eq:limsup_remaining_term} as follows:
\begin{align}
    &(1+d) \sum_{i \in I_n} \int_{Q'(y_i,r_n)} \int_{-\ell_n}^{\ell_n} \left[ \frac{1}{\varepsilon_n} W(\psi(y_i), T_a(\psi(y_i))w_n) + \varepsilon_n |a(\Psi(y,t)|^2 | \partial_t w_n|^2 \right] | \mathrm{det} (J\Psi)| \dd y \dd t \nonumber \\
    &= (1+d) \sum_{i \in I_n} \int_{Q'(y_i,r_n)} \int_{-\ell_n}^{\ell_n} \left[ \frac{1}{\varepsilon_n} W(\psi(y_i), z_n^i(t)) + \varepsilon_n \frac{|a(\Psi(y,t))|^2}{|a(\psi(y_i))|^2} |(z_n^i)'|^2 \right] | \mathrm{det} (J\Psi)| \dd y \dd t \nonumber \\
    &= (1+d) \sum_{i \in I_n} \int_{Q'(y_i,r_n)} \int_{-\ell_n}^{\ell_n} \left[ \frac{1}{\varepsilon_n} W(\psi(y_i), z_n^i(t)) + \varepsilon_n | (z_n^i)' |^2 \right] | \mathrm{det} (J\Psi)| \dd y \dd t \label{eq:limsup_remaining_term_last}\\
    &\qquad+ (1+d) \sum_{i \in I_n} \int_{Q'(y_i,r_n)} \int_{-\ell_n}^{\ell_n} \varepsilon_n \left( \frac{|a(\Psi(y,t))|^2}{|a(\psi(y_i))|^2} - 1 \right) | (z_n^i)'|^2 | \mathrm{det} (J\Psi)| \dd y \dd t \label{eq:limsup_afraction}.
\end{align}
Since by \ref{H6} we know that $a$ is continuous outside of the jump set, we get that $a(\Psi)$ is also continuous. By \ref{H3}, we have $|a(x)| \geq \delta > 0$, therefore
\begin{align*}
    \cfrac{|a(\Psi(y,t))|^2}{|a(\psi(y_i))|^2} - 1 &= \cfrac{| a(\Psi(y,t))|^2 - |a(\Psi(y_i,0))|^2}{|a(\psi(y_i))|^2} \\
    &\leq \cfrac{1}{\delta^2} \omega_{|a|^2} (|\Psi(y,t) - \Psi(y_i,0)|) \\
    &\leq \cfrac{1}{\delta^2} \omega_{|a|^2 \circ \Psi} (|(y,t) - (y_i,0)|) \\
    &\leq \cfrac{1}{\delta^2} \omega_{|a|^2 \circ \Psi} (f_n) \to 0.
\end{align*}
Substituting in \eqref{eq:limsup_afraction} we get
\begin{align*}
    &(1+d) \sum_{i \in I_n} \int_{Q'(y_i,r_n)} \int_{-\ell_n}^{\ell_n} \varepsilon_n \left( \frac{|a(\Psi(y,t))|^2}{|a(\psi(y_i))|^2} - 1 \right) | (z_n^i)' |^2 | \mathrm{det} (J\Psi)| \dd y \dd t \\
    &\leq (1+d) \frac{1}{r_n^{N-1}} r_n^{N-1} \frac{\ell_n}{\varepsilon_n} \frac{1}{\delta^2} \omega_{|a|^2 \circ \Psi} (f_n) \to 0.
\end{align*}
Since $| \mathrm{det} (J\Psi)|$ is a continuous function, it admits locally a modulus of continuity, and we can rewrite \eqref{eq:limsup_remaining_term_last} as
\begin{align}
    &(1+d) \sum_{i \in I_n} \int_{Q'(y_i,r_n)} \int_{-\ell_n}^{\ell_n} \left[ \frac{1}{\varepsilon_n} W(\psi(y_i), z_n^i(t)) + \varepsilon_n | (z_n^i)'|^2 \right] | \mathrm{det} (J\Psi)|(y,t) \dd y \dd t \nonumber \\
    &= (1+d) \sum_{i \in I_n} \int_{Q'(y_i,r_n)} \int_{-\ell_n}^{\ell_n} \left[ \frac{1}{\varepsilon_n} W(\psi(y_i), z_n^i(t)) + \varepsilon_n | (z_n^i)'|^2 \right] | \mathrm{det} (J\Psi)|(y_i,0) \dd y \dd t \label{eq:limsup_det_est_other} \\
    &\qquad+ (1+d) \sum_{i \in I_n} \int_{Q'(y_i,r_n)} \int_{-\ell_n}^{\ell_n} \left[ \frac{1}{\varepsilon_n} W(\psi(y_i), z_n^i(t)) + \varepsilon_n | (z_n^i)' |^2 \right] \cdot \nonumber \\
    &\hspace{7cm} \cdot\left( | \mathrm{det} (J\Psi)|(y,t) - | \mathrm{det} (J\Psi)|(y_i,0) \right) \dd y \dd t \label{eq:limsup_det_est}.
\end{align}
Using the modulus of continuity of $|\mathrm{det}(J \Psi)|$, we estimate \eqref{eq:limsup_det_est} as
\begin{align*}
    &(1+d) \sum_{i \in I_n} \int_{Q'(y_i,r_n)} \int_{-\ell_n}^{\ell_n} \left[ \frac{1}{\varepsilon_n} W(\psi(y_i), z_n^i(t)) + \varepsilon_n | (z_n^i)' |^2 \right]\cdot\\
    &\hspace{7cm}\cdot\left( | \mathrm{det} (J\Psi)|(y,t) - | \mathrm{det} (J\Psi)|(y_i,0) \right) \dd y \dd t \\
    &\leq (1+d) \frac{1}{r_n^{N-1}} r_n^{N-1} \ell_n \left[ \frac{1}{\varepsilon_n} C + \varepsilon_n \frac{1}{\varepsilon_n^2} C \right] \omega_{|\mathrm{det}(J\Psi)|} (f_n) \to 0.
\end{align*}
We are therefore left with estimating \eqref{eq:limsup_det_est_other}, namely
$$ (1+d) \sum_{i \in I_n} \int_{Q'(y_i,r_n)} \int_{-\ell_n}^{\ell_n} \left[ \frac{1}{\varepsilon_n} W(\psi(y_i), z_n^i(t)) + \varepsilon_n | (z_n^i)'(t)|^2 \right] | \mathrm{det} (J\Psi)|(y_i,0) \dd y \dd t. $$
Let us now zoom in on the term $| \mathrm{det} (J\Psi) | (y_i, 0)$: this would be the determinant of the matrix $(J\psi + t Jn \;| \; n ) $ computed in $(y_i,0)$, or equivalently, the determinant of the matrix $(J \psi (y_i) \; | \; n(y_i) )$. In particular, this represents the $N$-dimensional volume of the hyper-rectangle spanned by the columns of the matrix. Here we can see that, since $\nu$ is a unitary vector by definition perpendicular to the other columns, this volume will also be equal to the $N-1$-dimensional volume spanned by the first $N-1$ columns, which is equal to $| \mathrm{det}(J\psi)|(y_i)$.\\
Using now \eqref{eq:reparameterization_limsup}, we can write
\begin{align}
    &(1+d) \sum_{i \in I_n} \int_{Q'(y_i,r_n)} \int_{-\ell_n}^{\ell_n} \left[ \frac{1}{\varepsilon_n} W(\psi(y_i), z_n^i(t)) + \varepsilon_n | (z_n^i)'(t)|^2 \right] | \mathrm{det}(J\psi)|(y_i) \dd y \dd t \nonumber \\
    &\leq (1+d) \sum_{i \in I_n} \mathcal{H}^{N-1} (Q'(y_i, r_n)) |\mathrm{det}(J\psi)|(y_i)\cdot \nonumber\\
    &\hspace{5cm}\cdot\int_{-\ell_n}^{\ell_n} \left[ \frac{1}{\varepsilon_n} W(\psi(y_i), \gamma_i(g_n^i(t))) + \varepsilon_n |\gamma_i'(g_n^i(t))|^2 |(g_n^i)'(t)|^2 \right] \dd t \nonumber \\
    &\leq (1+d) \sum_{i \in I_n} \mathcal{H}^{N-1} (Q'(y_i, r_n)) |\mathrm{det}(J\psi)|(y_i) \int_{-1}^1 2 \sqrt{W(\psi(y_i), \gamma_i(s))} | \gamma_i'(s)| \dd s \label{eq:remainder_reparameterization} \\
    &\hspace{8cm}+ (1+d) 2 \overline{L} \mathcal{H}^{N-1}(\Sigma) \varepsilon_n \nonumber
\end{align}
This last term vanishes, therefore we can conclude by estimating \eqref{eq:remainder_reparameterization} with
\begin{align*}
    &\lim_{n \to \infty} \int_{\Sigma_n} \left[ \frac{1}{\varepsilon_n} W(x,u_n) + \varepsilon_n |\nabla u_n|^2 \right] \dd x \\
    &\leq (1+d) \lim_{n \to \infty} \sum_{i \in I_n} \mathcal{H}^{N-1} (Q'(y_i, r_n)) | \mathrm{det}(J\psi)|(y_i) \int_{-1}^1 2 \sqrt{W(\psi(y_i), \gamma_i(s))} | \gamma_i'(s)| \dd s \\
    &= (1+d) \lim_{n \to \infty} \sum_{i \in I_n} \mathcal{H}^{N-1} (Q'(y_i, r_n)) | \mathrm{det}(J\psi)|(y_i) \dd_W(\psi(y_i), u^+(\psi(y_i)),u^-(\psi(y_i))) \\
    &= (1+d) \int_D \dd_W (\psi(y), u^+(\psi(y)),u^-(\psi(y))) | \mathrm{det}(J\psi)|(y) \dd \mathcal{H}^{N-1} (y),
\end{align*}
where the last step follows from the continuity of $ y \mapsto \dd_W(\psi(y), u^+(\psi(y)),u^-(\psi(y)))$ and the continuity of $y \mapsto | \mathrm{det}(J\psi)|(y)$. The proof is then concluded by letting $d \to 0$.\\[0.5em]
\textbf{Step 5: General case.} Take $u \in BV(\Omega; \{a,b\})$. By Proposition \ref{prop:approx_sets} we know that there exists a sequence of functions $\{ v_n \}_n \subset BV(\Omega; \{a,b\})$ with piecewise $C^2$ jump set, such that 
$$ v_n \to u \text{ in } L^1(\Omega; \R^M) \text{ as } n \to \infty, \qquad \lim_{n\to \infty} \mathcal{F}_\infty(v_n) = \mathcal{F}_\infty (u). $$
We can therefore apply the previous steps to the function $v_n$ to obtain a sequence of functions $\{w_m^n\}_m \subset W^{1,2}(\Omega; \R^M)$ such that  
$$ w_m^n \to v_n \text{ in } L^1(\Omega; \R^M) \text{ as } m \to \infty, \qquad \lim_{m\to \infty} \mathcal{F}_{\varepsilon_m}(w_m^n) = \mathcal{F}_\infty (v_n). $$
Using a diagonalization argument, we find a sequence $\{m_n\}_n$ such that the sequence of functions defined by $\{u_n\}_n = \{w_{m_n}^n\}_n$ has the following properties
$$ u_n \to u \text{ in } L^1(\Omega; \R^M) \text{ as } n \to \infty, \qquad \lim_{n\to \infty} \mathcal{F}_{\varepsilon_n} (u_n) = \mathcal{F}_\infty (u), $$
and we conclude.
\end{proof}

We now show that it is possible to slightly adjust the proof of Theorem \ref{thm:limsup} and obtain the recovery sequence for the mass constrained case.

\begin{theorem}\label{thm:limsup_mass}
    Assume that \ref{H1}-\ref{H8} hold, with the function $f$ from $\ref{H5}$ satisfying $f(t) \leq |t|^\alpha$ with $\alpha > 1$.
    Let $m\in\R^M$.
    Let $u \in L^1(\Omega;\R^M)$. Then, there exists $\{u_n\}_n \subset H^1(\Omega;\R^M)$ such that $u_n \to u$ in $L^1(\Omega;\R^M)$, and moreover
    \[
        \mathcal{F}^m_\infty(u) = \lim_{n \to \infty} \mathcal{F}^m_{\eps_n}(u_n).
    \]
    for every sequence $\eps_n \to 0$.
\end{theorem}

\begin{proof}
The compactness and liminf proof work without any modifications. We only need to slightly modify the limsup proof, and this revolves around modifying the jump set approximation and the recovery sequence so that they satisfy the mass constraint.\\[0.5em]
\textbf{Step 1: Fixing the approximation.} We now have $u \in \mathrm{BV}(\Omega; \{ a,b\})$ with 
$$ \int_\Omega u \dd x = m a + (1-m) b. $$
Then, defining as before $A \coloneqq \{ u = a \}$, this implies that $|A| = m$. Since this is a set of finite perimeter, thanks to Proposition \ref{prop:approx_sets} we approximated it with a sequence $\{A_n\}_n$ of appropriate $C^2$ sets. The problem here is that we did not require these sets to satisfy the constraint $|A_n|=m$ for every $n \in \N$. Therefore, this is the first change that needs to be addressed: we thus follow an idea originally due to Ryan Murray, appropriately modified for our case (see \cite{murray2016slow}).

Since $A$ is a set of finite perimeter, we can take its reduced boundary $\partial^* A$. Let us now take $x_0, x_1 \in \Omega$ be points of density $0$ and $1$ respectively for $A$, and for $k \in \N$ define 
$$ D_k \coloneqq \left( A \cup B \! \left( \! x_0, \frac{1}{k} \! \right) \right) \setminus B \! \left( \! x_1, \frac{1}{k} \! \right). $$

It is possible to see that
$$ \mathbbm{1}_{D_k} \to \mathbbm{1}_A \quad \text{in } L^1(\Omega; \R) \quad \text{as } k \to \infty. $$
Moreover, we can check that
\begin{align*}
	\left| D \mathbbm{1}_{D_k} \right| (\Omega) &= \mathcal{H}^{N-1} \left(\Omega \cap \partial^* D_k \right) \\
	&\leq \mathcal{H}^{N-1} \left(\Omega \cap \partial^* A \right) + \mathcal{H}^{N-1} \left(\Omega \cap \partial B \! \left( \! x_1, \frac{1}{k} \! \right) \right) + \mathcal{H}^{N-1} \left(\Omega \cap \partial B \! \left( \! x_2, \frac{1}{k} \! \right) \right) \\
	&\to \mathcal{H}^{N-1} \left(\Omega \cap \partial^* A \right) \\
	&= \left| D \mathbbm{1}_{A} \right| (\Omega).
\end{align*}
Therefore we also have
$$ \lim_{k \to \infty} \mathrm{Per}(D_k;\Omega) = \mathrm{Per}(A;\Omega). $$
By definition of density for $x_0, x_1$ with respect to $A$, we have
$$ \lim_{r \to 0} \frac{|A \cap B(x_0,r)|}{|B(x_0,r)|} = 0, \qquad \lim_{r \to 0} \frac{|A \cap B(x_1,r)|}{|B(x_1,r)|} = 1. $$
Therefore, this implies that there exists $k$ large enough such that
$$ \left| A \cap B \left( x_0, \frac{1}{k} \right) \right| \leq \frac{1}{4} \left| B \left( x_0, \frac{1}{k} \right) \right|, \qquad \left| A \cap B \left( x_1, \frac{1}{k} \right) \right| \geq \frac{3}{4} \left| B \left( x_1, \frac{1}{k} \right) \right|. $$
Let now $(D_k^n)_n$ be the sequence of $C^2$ sets obtained by applying Proposition \ref{prop:approx_sets} to the set $D_k$. Since we have convergence in perimeter and $L^1$, this implies that for a fixed $k$, there exists $\widetilde{n}_1(k) \in \N$ such that for every $n \geq \widetilde{n}_1(k)$ we have
$$ \left| \mathrm{Per} \left( D_k; \Omega \right) - \mathrm{Per} \left( D_k^n; \Omega \right) \right| \leq \frac{1}{k}, \qquad \int_\Omega \left| \mathbbm{1}_{D_k} (x) - \mathbbm{1}_{D_k^n}(x) \right| \dd x \leq \frac{1}{k}. $$
Since these sets $D_k^n$ are obtained by using a standard mollifying procedure and taking a super-level set, we know that there exists $\widetilde{n}_2(k) \in \N$ such that for every $n \geq \max \{ \widetilde{n}_1(k), \widetilde{n}_2(k) \}$ we have both the previous inequalities and also
$$ B \left( x_0, \left( \frac{4}{5} \right)^{\frac{1}{N}} \frac{1}{k} \right) \subset D_k^n, \qquad B \left( x_1, \left( \frac{4}{5} \right)^{\frac{1}{N}} \frac{1}{k} \right) \subset \Omega \setminus D_k^n. $$
This holds because 
$$ \left( \frac{4}{5} \right)^{\frac{1}{N}} < 1 \quad \forall N \in \N. $$

We have now three cases to distinguish between. If $| D_k^n | = m$ then we do not have to do anything. Assume now that $ | D_k^n | > m $. Define $r_k^n >0$ to be the radius such that
$$ \left| B \left( x_0, r_k^n \right) \right| = \left| D_k^n \right| - m > 0, $$
and define
$$ A_k^n \coloneqq D_k^n \setminus B \left( x_0, r_k^n \right). $$
We now want to prove that $$ r_k^n < \left( \frac{4}{5} \right)^\frac{1}{N} \frac{1}{k}. $$
Since we know that $|A|=m$ and
$$ \left| A \cap B \left( x_0, \frac{1}{k} \right) \right| \geq 0, \qquad \left| A \cap B \left( x_1, \frac{1}{k} \right) \right| \geq \frac{3}{4} \left| B \left( x_1, \frac{1}{k} \right) \right|, $$
we have
\begin{align*}
	\left| D_k \right| &= |A| + \left| B \left( x_0, \frac{1}{k} \right) \setminus A \right| - \left| B \left( x_1, \frac{1}{k} \right) \cap A \right| \\
	&\leq |A| + \left| B \left( x_0, \frac{1}{k} \right) \right| - \left| B \left( x_0, \frac{1}{k} \right) \cap A \right| - \left| B \left( x_1, \frac{1}{k} \right) \cap A \right| \\
	&\leq |A| + \left| B \left( x_0, \frac{1}{k} \right) \right| - \frac{3}{4} \left| B \left( x_1, \frac{1}{k} \right) \right| \\
	&= |A| + \frac{1}{4} \left| B \left( x_0, \frac{1}{k} \right) \right| \\
	&= m + \left| B \left( x_0, \left( \frac{1}{4} \right)^\frac{1}{N} \frac{1}{k} \right) \right|.
\end{align*}
Therefore, for $n$ large enough we also get 
$$ \left| B \left( x_0, r_k^n \right) \right| = \left| D_k^n \right| - m \leq \left| B \left( x_0, \left( \frac{1}{4} \right)^\frac{1}{N} \frac{1}{k} \right) \right| < \left| B \left( x_0, \left( \frac{4}{5} \right)^\frac{1}{N} \frac{1}{k} \right) \right|, $$
therefore we got the estimate on $r_k^n$.
This in particular implies that the set $A_k^n$ is also $C^2$, and it follows that
$$ |A_k^n| = |D_k^n| - \left| B \left( x_0, r_k^n \right) \right| = |D_k^n| - |D_k^n| + m = m. $$
The case of $|D_k^n| < m$ follows analogously. 

Therefore we just proved that for every $k$ we can modify the sequence of sets obtained from Proposition \ref{prop:approx_sets} applied to $D_k$ such that every element of the sequence satisfies the mass constraint. Using a diagonal argument, we can therefore obtain the desired conclusion.\\[0.5em]
\textbf{Step 2: Fixing the recovery sequence.}
Let $u_n$ be defined as in \eqref{eq:recovery_sequence}. In general it is not true that this function satisfies the mass constraint, therefore we need to modify it accordingly.\\
Let $N \geq 2$ and define
$$ m_n \coloneqq \int_\Omega u_n (x) \dd x. $$
If, for a given $n \in \N$, we have $m_n = m$, then we are done. Let us suppose that $m_n \neq m$. 
Let us recall the definition of $\Sigma_n$:
$$ \Sigma_n \coloneqq \{ x \in \Omega : | \mathrm{dist}(x, \Sigma) | \leq \tau_n \}, $$
which is the set where $u_n(x) \notin \{a, b\}$. 

Let $x_0 \in \Omega \setminus \Sigma_n$ such that $u_n (x) = u(x) = a(x)$ in a neighborhood of $x_0$. Let now $(r_n)_{n \in \N}$ be an infinitesimal sequence and define $B_n \coloneqq B (x_0, r_n)$ to be such neighborhood. Since $r_n \leq \mathrm{dist}(x_0,\Sigma_n) $, we can modify $u_n$ as follows
\begin{gather*}
    v_n (x) \coloneqq \begin{cases}
        u_n(x) \qquad &x \in \Omega \setminus B_n, \\
        a(x) + c_n (m_n - m) \left( 1 - \frac{|x - x_0|}{r_n} \right) \qquad &x \in B_n,
    \end{cases}
\end{gather*}
where $c_n \in \R$ is to be determined by enforcing the mass constraint.

Therefore we must have
\begin{align*}
    m &= \int_\Omega v_n (x) \dd x \\
    &= \int_{\Omega \setminus B_n} u_n(x) \dd x + \int_{B_n} \left[ a(x) + c_n (m_n - m) \left( 1 - \frac{|x - x_0|}{r_n} \right) \right] \dd x \\
    &= \int_\Omega u_n (x) \dd x - \int_{B_n} u_n (x) \dd x + \int_{B_n} a(x) \dd x + c_n (m_n - m) \int_{B_n} \left[ 1 - \frac{|x - x_0|}{r_n}\right] \dd x \\
    &= m_n + c_n (m_n - m) N \omega_N r_n^N \int_0^1 ( 1 - s) s^{N-1} \dd s \\
    &= m_n + c_n (m_n - m) r_n^N \frac{\omega_N}{N+1},
\end{align*}
which implies
$$ c_n = - \frac{N+1}{\omega_N} \ \frac{1}{r_n^N}.$$
Therefore with this choice of $c_n$, the function $v_n$ satisfies the mass constraint. 
Moreover, we also have
$$ | m_n - m| \leq \int_\Omega |u_n - u| \dd x = \int_{\Sigma_n} |u_n - u| \dd x \leq |b-a| \ln(\Sigma_n) \leq C \varepsilon_n \mathcal{H}^{N-1}(\Sigma). $$
We need to check that $v_n$ still converges in $L^1$, and that is does not change the energy in the limit. Checking the $L^1$ convergence is easy, since the way we chose the constant $c_n$ implies
$$ \left| \int_{B_n} c_n (m_n - m) \left( 1 - \frac{|x - x_0|}{r_n} \right) \dd x  \right| = | m - m_n | \leq C \varepsilon_n \to 0. $$
For the energy convergence, we need to check that
$$ \lim_{n \to \infty} F_n^{(1)} (v_n, B_n) = 0. $$
Let us check first the potential energy term: we can use Assumption \ref{H5} with $f(t) \leq t^\alpha$ and $\alpha > 1$ to see that
\begin{align}
	\int_{B_n} \frac{1}{\varepsilon_n} W \left( x, v_n(x) \right) \dd x &\leq \int_{B_n} \frac{1}{\varepsilon_n} | v_n(x) - a(x)|^\alpha \dd x \\
    &\leq C \frac{r_n^N}{\varepsilon_n} | c_n |^\alpha |m_n - m|^\alpha \\
    &\leq C \left( \frac{\varepsilon_n}{r_n^N} \right)^\alpha.
\end{align}
Therefore, for this term to vanish in the limit, it is enough to have 
$$ r_n = \varepsilon_n^\beta, \qquad 0 < \beta < \frac{1}{N}, $$
and we conclude.
\end{proof}

\begin{remark}
    The necessary assumption on the wells, namely \ref{H5} with $f(t) \leq |t|^\alpha$ for a $\alpha > 1$, implies uniform differentiability of $p \mapsto W(x,p)$ for $p \in \{a(x), b(x)\}$.
\end{remark}

\begin{remark}
    Note that a similar procedure has been carried out in \cite{cristoferi2025phaseseparationmultiplyperiodic} for fixing the approximation, with the difference that the points $x_0, x_1$ lay on the reduced boundary. Both methods work, but for our setting it was easier to have the points lay outside of the reduced boundary.
\end{remark}



%%%%%%%%%%%%%%%%%%%%%%%%%%%%%%%%%%%%%%%%%%%%
%%%%%%%%%%%%%%%%%%%%%%%%%%%%%%%%%%%%%%%%%%%%
%%%%%%%%%%%%%%%%%%%%%%%%%%%%%%%%%%%%%%%%%%%%
%%%%%%%%%%%%%%%%%%%%%%%%%%%%%%%%%%%%%%%%%%%%



\appendix
\section{Auxiliary results}


Finally, we prove two technical results that we will need in the proof of the liminf inequality (see Theorem \ref{thm:liminf}).

\begin{lemma}\label{lemma:technical}
        Suppose that $g: \R^N \times \R^M \to [0, \infty)$ is continuous and there exist $p \in \R^N, (p_n) \subset \R^N$ and $L_n \in \R^{m \times m}$ with $p_n \to p$ in $\R^N$ and $L_n \to Id_m$ in $\R^{m \times m}$. Then,
        \[
            \inf_{r \in \R^M} g(p_n, L_n r) \xrightarrow{n \to \infty} \inf_{r \in \R^M} f(p, r).
        \]
    \end{lemma}
    
    \begin{proof}
        Consider the functionals
        \[
            G_n(r) := g(p_n, L_n r),
        \]
        \[
            G(r) := g(p, r).
        \]
        Let $r_n \to r$. By continuity 
        \[
            F_n(r_n) = f(p_n, L_n r_n) \xrightarrow{n \to \infty} f(p, r) = F(r),
        \]
        which, in particular, implies that the $\Gamma$-limit of $G_n$ is $G$. By the fundamental property of $\Gamma$-convergence, the infimum of $G_n$ converges to the infimum of $G$. This concludes the proof.
    \end{proof}
    
    \begin{lemma}\label{lemma:technical_approx_limits}
        Let $u \in BV(\Omega; \R^M)$, $0 \in J_u$ with $\nu_{0} = e_n$, and $F \subset (0,1)$ be a set with $|F| = 1$. Then, there exists a sequence $\{\rho_m\} \subset F$ with $\rho_m \to 0^+$ such that 
        \[
            \H^{N-1}(J_u \cap \partial (-\rho_m/2, \rho_m/2)) = 0,
        \]
        and for the rescaled trace $u|_{\partial Q_{\rho_m}(x_0)}$ restricted to $(-\rho_m/2, \rho_m/2)\times \{\pm \rho_m\}$ we have 
        \[
            u(\rho_m \cdot, \pm \rho_m/2) \to u^\pm(x_0)
        \]
        in $L^1((-1/2, 1/2)^{N-1})$.
    \end{lemma}
    
    \begin{proof}
        Since $u^\pm(x_0)$ are one side approximate limits we have 
        \[
            \lim_{\rho \to 0^+} \int_{(0, 1)^N} |u(\rho x) - u^+(x_0)| \, dx = 0
        \]
        and
        \[
            \lim_{\rho \to 0^+} \int_{(-1, 0)^N} |u(\rho x) - u^-(x_0)| \, dx = 0.
        \]
        In particular, we have by Fubini's Theorem
        \[
            \lim_{\rho \to 0} \int_0^1 \int_{(0, 1)^{N}} |u(\rho x', \rho t) - u^+(x_0)| \, dx' \, dt= 0,
        \]
        where $u(\rho x', \rho t)$ agrees with the two-sided trace at $\partial (0,1)^N$ of $u$ that exists since 
        \begin{align}\label{eq:eligible_rho}
            \H^{N-1}(J_u \cap \partial (-\rho/2, \rho/2)) = 0
        \end{align}
        for all but countably many $\rho > 0$. In particular, for a set $E_\rho \subset (0,1)$ of full measure we obtain $t \in E_\rho$
        \[
            \lim_{\rho \to 0} \int_{(0, 1)^{N-1}} |u(\rho x', \rho t) - u^+(x_0)| \, dy' = 0,
        \]
        and, therefore, it follows that 
        \[
            \lim_{\rho \to 0} \int_{(0, 1)^{N-1}} |u(\rho t x', \rho t) - u^+(x_0)| \, dy' = \lim_{\rho \to 0} \frac1 {t^{N-1}} \int_{(0, t)^{N-1}} |u(\rho x', \rho t) - u^+(x_0)| \, dy' = 0.
        \]
        Now, choose a sequence $\{ \tilde \rho_m\} \subset (0,1)$ with $\tilde \rho_m \to 0$ for which \eqref{eq:eligible_rho} holds. Then, with an $t_m \in E_{\rho_m}$ such that $t_m\rho_m \in F$ holds, the sequence $\rho_m := t_m \tilde \rho_m$ satisfies the sought-after requirements.
    \end{proof}
    


%%%%%%%%%%%%%%%%%%%%%%%%%%%%%%%%%%%%%%%%%%%%
%%%%%%%%%%%%%%%%%%%%%%%%%%%%%%%%%%%%%%%%%%%%
%%%%%%%%%%%%%%%%%%%%%%%%%%%%%%%%%%%%%%%%%%%%
%%%%%%%%%%%%%%%%%%%%%%%%%%%%%%%%%%%%%%%%%%%%



\subsection*{Acknowledgment}
The authors were partially supported under NWO-OCENW.M.21.336, MATHEMAMI - Mathematical Analysis of phase Transitions in HEterogeneous MAterials with Materials Inclusions.



%%%%%%%%%%%%%%%%%%%%%%%%%%%%%%%%%%%%%%%%%%%%
%%%%%%%%%%%%%%%%%%%%%%%%%%%%%%%%%%%%%%%%%%%%
%%%%%%%%%%%%%%%%%%%%%%%%%%%%%%%%%%%%%%%%%%%%
%%%%%%%%%%%%%%%%%%%%%%%%%%%%%%%%%%%%%%%%%%%%



\bibliographystyle{siam}
\bibliography{Bibliography}

\end{document}
